\section{Bessel Poincar\'e series}\label{sec:bessel}

This section proves the three automorphic inputs (d) and (e) of Theorem~\ref{r3:A3} from the statements listed in it: the near orthogonality of Bessel Poincar\'e series, the share of irreducible rows, and the Bessel identity for the period maps.



\subsection{The three statements}\label{w2:sec:statements}

We keep the notation of Section~\ref{r3:sec}. The field $K$ is a totally complex Galois number field, $v$ is a complex place of $K$, $Y$ is a $K$-vector space with a hermitian lattice $R$ of rank $D_{\mathrm{rows}}$, and for each word $u$ of the finite set $W$ the period map is a $K$-linear map $B_u\colon V_u\to Y^*$. Write $Y_v=Y\otimes_{K,\sigma_v}\C$ and $V_{u,v}=V_u\otimes_{K,\sigma_v}\C$, with the hermitian norms of $R$ and of $N_u$ at $v$.

On the automorphic side $F$ is a totally real field of degree $f$, $E/F$ is a CM extension, $V$ is a nondegenerate hermitian space over $E$, and $G=\mathrm U(V)$ is its unitary group over $F$. Write $\A$ and $\Af$ for the adeles and the finite adeles of $F$, $G_\infty=G(F\otimes_{\mathbb Q}\R)$ and $[G]=G(F)\backslash G(\A)$. The space $\mathcal H$ is the finite dimensional space of automorphic forms on $G$ whose archimedean component is a fixed discrete series and whose finite part is fixed by the level idempotent of the construction, with the Petersson inner product, which is conjugate linear in the first variable. Its dimension is $D_{\mathrm{all}}$. The number $D^*$ is the product of the formal degree of the discrete series, the volume of $[G]$ and the value at $1$ of the level idempotent.

For each word $u$ the column space $\mathcal C_u$ has dimension $d^S_u$, the Poincar\'e map $\mathcal P_u\colon\mathcal C_u\to\mathcal H$ sends a column form to its Bessel Poincar\'e series (Section~\ref{w2:sec:unfold}), $c_u>0$ is its normalizing volume, and
\[
K_{uu'}=(c_uc_{u'})^{-1/2}\mathcal P_u^*\mathcal P_{u'}\colon\mathcal C_{u'}\to\mathcal C_u .
\]
The family parameters are $x$, $\varepsilon$, $b=e^{x^\varepsilon}$, $L=\log b=x^\varepsilon$, the ranks $m,k,r$ of the datum and $\Gamma_0=\exp(C_\Gamma m^5fL)$.

\begin{theorem}[near orthogonality]\label{w2:S}
There are absolute constants $C_0\ge1$ and $C_\Gamma>0$ with the following property. For every $\varepsilon$, every $x$ large in terms of $\varepsilon$, and every family of data of Part~II at the scale $x$ whose words satisfy
\begin{equation}\label{w2:window}
C_0^{-1}\le d^S_u\Gamma_0^2/D^*\le C_0\qquad\text{for every }u\in W,
\end{equation}
the Gram operators satisfy
\[
\lVert K_{uu'}\rVert_\HS^2\le(1+\eta_0)\frac{d^S_ud^S_{u'}}{D^*}\ (u\ne u'),\qquad\lVert K_{uu}-1\rVert_\HS^2\le e_0\,d^S_u,\qquad D_{\mathrm{all}}\le(1+\eta_1)D^*,
\]
with $e_0\ge0$ and $\eta_0,\eta_1,\sqrt{e_0}\le\exp(-x^{\varepsilon/8})$.
\end{theorem}

\begin{theorem}[share of irreducible rows]\label{w2:I}
Let $J\ge1$ be the fixed integer of the construction and $\varepsilon_z=\varepsilon/J$. Let $E_{\mathrm{row}}\subset\mathcal H$ be the span of the image of $R\otimes_{K,\sigma_v}\C$ under the row embedding, and define $\delta$ by $D_{\mathrm{rows}}=\dim E_{\mathrm{row}}=(1-\delta)D_{\mathrm{all}}$. For every family of data of Part~II and every $x$ large in terms of the family constants,
\[
0\le\delta\le\exp\bigl(-x^{\varepsilon_z/7}\bigr).
\]
\end{theorem}

\begin{theorem}[Bessel identity]\label{w2:B}
For every family of data of Part~II and every $x$ large, the following holds at the complex place $v$. There are a unitary operator $\mathcal V$ of $Y_v$, and for each $u\in W$ a unitary operator $\mathcal U_u$ of $V_{u,v}$ and a number $a_u\ne0$, such that
\begin{equation}\label{w2:besselid}
(B_{u,v}\varphi)(F)=a_u\,\bigl\langle\mathbf s(\mathcal VF),\,\mathcal P_u\iota_u\mathcal U_u\varphi\bigr\rangle_{\mathcal H}\qquad(\varphi\in V_{u,v},\ F\in Y_v),
\end{equation}
where $\mathbf s\colon Y_v\to\mathcal H$ is the conjugate linear isometric row embedding and $\iota_u\colon V_{u,v}\to\mathcal C_u$ is the linear isometric column embedding of Definition~\ref{r3:bessel}. In addition $B_u$ is $K$-rational in the sense that $(B_{u,v}(\varphi\otimes1))(y\otimes1)=\sigma_v\bigl(B_u(\varphi)(y)\bigr)$ for $\varphi\in V_u$ and $y\in Y$.
\end{theorem}

Theorem~\ref{w2:B} together with Theorem~\ref{w2:I} is hypothesis (d) of Theorem~\ref{r3:A3}, and Theorem~\ref{w2:S} is hypothesis (e) there. 

\subsection{The Bessel identity}\label{w2:sec:bessel}

The identity \eqref{w2:besselid} is the end of a chain with four links. The period map is defined by a period of the chain of theta lifts on the definite side. The seesaw relation transfers every linear relation among those periods to the local Bessel periods of the theta lifts on the indefinite side. Unfolding the Poincar\'e series transfers every linear relation among local Bessel periods to the pairings with Poincar\'e series. One nonzero pairing then turns the transfer of relations into proportionality. Section~\ref{w2:sec:unfold} proves the unfolding, Section~\ref{w2:sec:prop} proves the passage from transfer of relations to proportionality, and Section~\ref{w2:sec:assembleB} assembles the theorem.

\subsubsection{Bessel Poincar\'e series and their unfolding}\label{w2:sec:unfold}

Fix a word $u$ and drop it from the notation where no confusion arises. The construction gives the following objects.

\begin{enumerate}[label=(P\arabic*),leftmargin=*]
\item\label{w2:P1} A closed $F$-subgroup $S=N\rtimes H$ of $G$, the Bessel subgroup of the flag of $u$. Here $N$ is a unipotent $F$-group and $H=\mathrm U(W)$ for a hermitian space $W$ over $E$ that is positive definite at every real place of $F$. A unitary character $\psi\colon N(\A)\to\C^\times$, trivial on $N(F)$ and invariant under conjugation by $H(\A)$.
\item\label{w2:P2} A compact open subgroup $K_2\subset G(\Af)$ with a unitary character $\chi_2$, and an element $x_u\in G(\Af)$.
\item\label{w2:P3} A finite dimensional unitary representation $(\sigma,V_\sigma)$ of the compact group $H_\infty=H(F\otimes\R)$, extended to $S_\infty=N_\infty H_\infty$ by $\sigma(nh)=\psi_\infty(n)\sigma(h)$, and a measurable function $W_u\colon G_\infty\to V_\sigma$ with $W_u(sy)=\sigma(s)W_u(y)$ for $s\in S_\infty$ and
\begin{equation}\label{w2:Wint}
\int_{S_\infty\backslash G_\infty}\lVert W_u(y)\rVert\,d\nu(y)<\infty .
\end{equation}
\item\label{w2:P4} The space $\mathcal C$ of column forms. A column form is a linear map $T$ from $V_\sigma$ to functions on $H(\A)$ with $T(\xi)(\gamma a)=T(\xi)(a)$ for $\gamma\in H(F)$, $T(\sigma(h)\xi)(a)=T(\xi)(ah)$ for $h\in H_\infty$, and $\lvert T(\xi)(a)\rvert\le C_T\lVert\xi\rVert$.
\item\label{w2:P5} For each $T\in\mathcal C$ a measurable seed $\mathrm{Ph}_T$ on $G(\A)$, supported on $G_\infty\times S(\Af)x_uK_2$, with $\mathrm{Ph}_T(gk)=\chi_2(k)\mathrm{Ph}_T(g)$ for $k\in K_2$, and
\begin{equation}\label{w2:seed}
\mathrm{Ph}_T\bigl(na\cdot(y,x_u)\bigr)=\psi(n)\,T\bigl(W_u(y)\bigr)(a)\qquad(n\in N(\A),\ a\in H(\A),\ y\in G_\infty).
\end{equation}
\end{enumerate}

Measures are fixed as follows. $G(\A)=G_\infty\times G(\Af)$ has a product Haar measure, $S(\A)$, $N(\A)$ and $H(\A)$ have Haar measures with $ds=dn\,da$, discrete groups have counting measure, and $\nu$ is the quotient measure on $S_\infty\backslash G_\infty$. Lemma~\ref{w2:unimod} below shows that these quotient measures exist.

For a bounded measurable function $\theta$ on $G(\A)$ that is left $N(F)$-invariant put
\begin{equation}\label{w2:per}
\per(\Psi,\theta)=\int_{H(F)\backslash H(\A)}\int_{N(F)\backslash N(\A)}\theta(na)\,\psi(n)\,\Psi(a)\,dn\,da
\end{equation}
for every bounded measurable left $H(F)$-invariant function $\Psi$ on $H(\A)$. This is the Bessel period of $\theta$ against $\Psi$.

\begin{lemma}\label{w2:unimod}
Under \ref{w2:P1} the groups $N(\A)$, $H(\A)$ and $S(\A)$ are unimodular, the quotients $N(F)\backslash N(\A)$ and $H(F)\backslash H(\A)$ are compact, and $S(F)\backslash S(\A)$ is compact.
\end{lemma}

\begin{proof}
The modular function of $S(\A)=N(\A)\rtimes H(\A)$ is $nh\mapsto\lvert\det\mathrm{Ad}(h)|_{\mathrm{Lie}\,N}\rvert_\A^{-1}$ times the modular functions of $N(\A)$ and $H(\A)$. The map $h\mapsto\det\mathrm{Ad}(h)|_{\mathrm{Lie}\,N}$ is an $F$-rational character of $H=\mathrm U(W)$. Every $F$-rational character of $\mathrm U(W)$ factors through the determinant $\mathrm U(W)\to\mathrm U(1)$, and $\mathrm U(1)=\ker(N_{E/F}\colon R_{E/F}\mathbb G_m\to\mathbb G_m)$ is an anisotropic torus, so it has no nontrivial $F$-rational character. Hence $\det\mathrm{Ad}(h)|_{\mathrm{Lie}\,N}=1$ identically, so at every place and on the adelic points. The same argument shows that $S_\infty$ and $S(\Af)$ are unimodular. Unipotent groups and reductive groups have unimodular adelic points, so $S(\A)$ is unimodular. The quotient $N(F)\backslash N(\A)$ is compact for every unipotent $F$-group \cite{w2borel}. Since $W$ is positive definite at a real place, it has no isotropic vector over $F$, so $H$ is $F$-anisotropic and $H(F)\backslash H(\A)$ is compact by the theorem of Borel and Harish-Chandra and of Mostow and Tamagawa \cite{w2bhc}. The image of the compact set $C_N\times C_H$ under multiplication, with $C_N$ and $C_H$ compact sets mapping onto the two quotients, maps onto $S(F)\backslash S(\A)$, because $S(F)=N(F)H(F)$ and $H(\A)$ normalizes $N(\A)$.
\end{proof}

\begin{proposition}[unfolding]\label{w2:unfold}
Assume \ref{w2:P1} to \ref{w2:P5} and let $T\in\mathcal C$. The series
\[
P_T(g)=\sum_{\gamma\in S(F)\backslash G(F)}\mathrm{Ph}_T(\gamma g)
\]
converges absolutely for almost every $g\in G(\A)$, defines an element of $L^1([G])$ that does not depend on the choice of representatives, and is right $(K_2,\chi_2)$-equivariant. There is a number $c_u>0$, depending on the Haar measures and on $u$ and not on $T$, such that for every bounded measurable function $\theta$ on $G(\A)$ that is left $G(F)$-invariant and right $(K_2,\overline{\chi_2})$-equivariant,
\begin{equation}\label{w2:perlaw}
\int_{[G]}\theta(g)\,P_T(g)\,dg=c_u\int_{S_\infty\backslash G_\infty}\per\Bigl(T\bigl(W_u(y)\bigr),\ R_{(y,x_u)}\theta\Bigr)\,d\nu(y),
\end{equation}
where $R_{(y,x_u)}\theta(g)=\theta(g\cdot(y,x_u))$, and both sides converge absolutely.
\end{proposition}

\begin{proof}
\emph{Step 1, integration in stages.} Since $G(\A)$, $S(\A)$, $S(F)$ and $G(F)$ are unimodular (Lemma~\ref{w2:unimod}), the quotient integral formula of Weil \cite[\S2.6]{w2folland} applies to $S(F)\subset S(\A)\subset G(\A)$ and to $S(F)\subset G(F)\subset G(\A)$. For every measurable $\Phi\ge0$ on $G(\A)$ that is left $S(F)$-invariant,
\begin{align}
\int_{S(F)\backslash G(\A)}\Phi&=\int_{S(\A)\backslash G(\A)}\int_{S(F)\backslash S(\A)}\Phi(s\dot g)\,ds\,d\dot g,\label{w2:stage1}\\
\int_{S(F)\backslash G(\A)}\Phi&=\int_{[G]}\sum_{\gamma\in S(F)\backslash G(F)}\Phi(\gamma g)\,dg.\label{w2:stage2}
\end{align}
The second identity is the unfolding. To see it directly, let $\mathfrak F$ be a measurable fundamental domain for $G(F)$ in $G(\A)$, which exists because $G(F)$ is a countable discrete subgroup of a second countable locally compact group, and let $D$ be a set of representatives of $S(F)\backslash G(F)$. The sets $\delta\gamma\mathfrak F$ with $\delta\in S(F)$ and $\gamma\in D$ are disjoint and cover $G(\A)$, so $\bigcup_{\gamma\in D}\gamma\mathfrak F$ is a fundamental domain for $S(F)$. Left invariance of Haar measure and Tonelli's theorem give \eqref{w2:stage2}.

\emph{Step 2, the finite part.} Put $\Omega=S(\Af)x_uK_2$, an open subset of $G(\Af)$. Let $\Phi\ge0$ be measurable on $G(\A)$ with $\Phi(sg)=\Phi(g)$ for $s\in S(\A)$, right $K_2$-invariant and supported on $G_\infty\times\Omega$. Then $\Phi$ factors through $S(\A)\backslash G(\A)=(S_\infty\backslash G_\infty)\times(S(\Af)\backslash G(\Af))$ and is a function of the first coordinate on the subset $S(\Af)\backslash\Omega$, which is the image of $K_2$ under $k\mapsto S(\Af)x_uk$. The quotient measure on a product of quotients is the product of the quotient measures, so
\begin{equation}\label{w2:finite}
\int_{S(\A)\backslash G(\A)}\Phi(\dot g)\,d\dot g=c_u\int_{S_\infty\backslash G_\infty}\Phi\bigl((y,x_u)\bigr)\,d\nu(y),
\end{equation}
where $c_u$ is the quotient measure of $S(\Af)\backslash\Omega$. It is finite because $S(\Af)\backslash\Omega$ is the continuous image of the compact group $K_2$, and positive because it is open.

\emph{Step 3, absolute convergence.} Let $\theta$ be as in the statement. The product $\Phi_0=\lvert\theta\,\mathrm{Ph}_T\rvert$ is left $S(F)$-invariant, because $\theta$ is left $G(F)$-invariant and $\mathrm{Ph}_T$ is left $S(F)$-invariant. For the latter, the support $G_\infty\times\Omega=S(\A)\cdot(G_\infty\times\{x_u\})\cdot K_2$ is stable under $S(F)$, and for $\gamma=n_0a_0\in S(F)$, \eqref{w2:seed} gives $\mathrm{Ph}_T(\gamma na(y,x_u))=\psi(n_0)\psi(a_0na_0^{-1})T(W_u(y))(a_0a)=\psi(n)T(W_u(y))(a)$, by the triviality of $\psi$ on $N(F)$, its invariance under $H(\A)$, and the $H(F)$-invariance in \ref{w2:P4}. It is right $K_2$-invariant, because the characters $\chi_2$ and $\overline{\chi_2}$ cancel, and it is supported on $G_\infty\times\Omega$. By \eqref{w2:stage1}, then \eqref{w2:finite} applied to the left $S(\A)$-invariant function $g\mapsto\int_{S(F)\backslash S(\A)}\Phi_0(sg)\,ds$, then \eqref{w2:seed} and the integration formula $ds=dn\,da$ on $S(F)\backslash S(\A)$,
\begin{align*}
\int_{S(F)\backslash G(\A)}\Phi_0
&=c_u\int_{S_\infty\backslash G_\infty}\int_{[H]}\int_{[N]}\bigl\lvert\theta(na(y,x_u))\bigr\rvert\,\bigl\lvert T(W_u(y))(a)\bigr\rvert\,dn\,da\,d\nu(y)\\
&\le c_u\lVert\theta\rVert_\infty\vol[N]\vol[H]\,C_T\int_{S_\infty\backslash G_\infty}\lVert W_u\rVert\,d\nu,
\end{align*}
which is finite by Lemma~\ref{w2:unimod} and \eqref{w2:Wint}. The same computation applies to $\Phi_0=\lvert\mathrm{Ph}_T\rvert$, which is right $K_2$-invariant because $\lvert\chi_2\rvert=1$, and with \eqref{w2:stage2} it shows that $\sum_\gamma\lvert\mathrm{Ph}_T(\gamma g)\rvert$ is integrable on $[G]$, hence finite almost everywhere, so $P_T$ converges absolutely almost everywhere and lies in $L^1([G])$. Independence of representatives is the left $S(F)$-invariance of $\mathrm{Ph}_T$, and right equivariance under $(K_2,\chi_2)$ passes to the sum.

\emph{Step 4, the identity.} Since $\theta\,\mathrm{Ph}_T$ is absolutely integrable on $S(F)\backslash G(\A)$ by Step 3, the identities \eqref{w2:stage2}, \eqref{w2:stage1} and \eqref{w2:finite} hold for it by Fubini's theorem applied to the positive and negative parts of its real and imaginary parts. The left side of \eqref{w2:stage2} is $\int_{[G]}\theta P_T$ because $\theta$ is left $G(F)$-invariant. On the right, \eqref{w2:seed} gives
\[
\begin{aligned}
\int_{S(F)\backslash S(\A)}(\theta\,\mathrm{Ph}_T)(s(y,x_u))\,ds&=\int_{[H]}\int_{[N]}\theta(na(y,x_u))\,\psi(n)\,T(W_u(y))(a)\,dn\,da\\
&=\per\bigl(T(W_u(y)),R_{(y,x_u)}\theta\bigr).
\end{aligned}
\]
This is \eqref{w2:perlaw}.
\end{proof}

\begin{remark}\label{w2:rem:perlaw}
The argument shows $c_u>0$. The identity \eqref{w2:perlaw} does not need $P_T\in L^2([G])$, nor that $P_T$ is fixed by the archimedean projector. Both facts are needed to know that $P_T\in\mathcal H$, which Definition~\ref{w2:anchor} uses.
\end{remark}

\begin{corollary}[transfer of relations from Bessel periods to Poincar\'e pairings]\label{w2:bottom}
Assume \ref{w2:P1} to \ref{w2:P5}. Let $T_1,\dots,T_\ell\in\mathcal C$, let $\theta_1,\dots,\theta_\ell$ be bounded measurable functions on $G(\A)$, left $G(F)$-invariant and right $(K_2,\overline{\chi_2})$-equivariant, and let $\alpha_1,\dots,\alpha_\ell\in\C$. If
\begin{equation}\label{w2:bothyp}
\sum_{i=1}^\ell\alpha_i\,\per\bigl(T_i(\xi),\,R_{(y,x_u)}\theta_i\bigr)=0\qquad\text{for every }\xi\in V_\sigma\text{ and every }y\in G_\infty,
\end{equation}
then $\sum_i\alpha_i\int_{[G]}\theta_iP_{T_i}=0$. If \eqref{w2:bothyp} holds for $\xi$ in a spanning set of $V_\sigma$ and every $y$, it holds for every $\xi$.
\end{corollary}

\begin{proof}
The period \eqref{w2:per} is linear in $\Psi$ and $T_i$ is linear, so the left side of \eqref{w2:bothyp} is linear in $\xi$, which proves the last sentence. Apply Proposition~\ref{w2:unfold} to each $i$ with the same $c_u$, and put $\xi=W_u(y)$ in \eqref{w2:bothyp}. The integrand of the sum of the right sides of \eqref{w2:perlaw} vanishes for every $y$.
\end{proof}

\subsubsection{From transfer of relations to proportionality}\label{w2:sec:prop}

\begin{lemma}\label{w2:prop}
Let $X$ be a set and $\lambda,\mu\colon X\to\C$ functions. Assume that for every finite family $x_1,\dots,x_\ell\in X$ and $\alpha_1,\dots,\alpha_\ell\in\C$ with $\sum_i\alpha_i\lambda(x_i)=0$ one has $\sum_i\alpha_i\mu(x_i)=0$. Assume $\mu(x_0)\ne0$ for some $x_0\in X$. Then $\lambda(x_0)\ne0$, and $\mu=\kappa\lambda$ on $X$ with $\kappa=\mu(x_0)/\lambda(x_0)\ne0$.
\end{lemma}

\begin{proof}
If $\lambda(x_0)=0$, the family consisting of $x_0$ with coefficient $1$ satisfies the hypothesis, so $\mu(x_0)=0$, a contradiction. For $x\in X$ the family $(x_0,x)$ with coefficients $(\lambda(x),-\lambda(x_0))$ satisfies $\lambda(x)\lambda(x_0)-\lambda(x_0)\lambda(x)=0$, so $\lambda(x)\mu(x_0)-\lambda(x_0)\mu(x)=0$.
\end{proof}

\begin{lemma}[central translations]\label{w2:central}
Let $z$ be an element of the center of $G(\Af)$ and $R_z\theta(g)=\theta(gz)$. Then $R_z$ is a unitary operator of $L^2([G])$ that commutes with the right action of $G(\A)$. It maps the right $(K_2,\chi_2)$-equivariant functions to themselves and preserves every $G(\A)$-isotypic subspace, in particular $\mathcal H$.
\end{lemma}

\begin{proof}
$R_z$ is well defined on $[G]$ and preserves Haar measure because $G(\A)$ is unimodular, and it is invertible with inverse $R_{z^{-1}}$. Since $z$ is central, $R_zR_h=R_hR_z$ for $h\in G(\A)$, which gives the last two assertions, and $\theta(gkz)=\theta(gzk)=\chi_2(k)\theta(gz)$ for $k\in K_2$.
\end{proof}

\subsubsection{Assembly of the Bessel identity}\label{w2:sec:assembleB}

The theorem uses four inputs of the construction. The first is a definition, the second is the seesaw relation of the theta tower, the third is a pair of dictionaries between the abstract realization of Definition~\ref{r3:bessel} and the automorphic objects of Section~\ref{w2:sec:unfold}, and the fourth is rationality. They are stated here in the form used, for one word $u$ at the place $v$.

\begin{hypothesis}[the period map and the seesaw relation]\label{w2:H-see}
There are a central element $z\in G(\Af)$, a linear map $U_u$ of $V_{u,v}$ with $\lVert U_u\varphi\rVert=r_u\lVert\varphi\rVert$ for some $r_u>0$, a set $X_u$ containing $V_{u,v}\times Y_v$, a function $\lambda_u\colon X_u\to\C$ with $\lambda_u(\varphi,F)=(B_{u,v}\varphi)(F)$ on $V_{u,v}\times Y_v$, and for each $x\in X_u$ a column form $T_x\in\mathcal C_u$ and a function $\theta_x\in\mathcal H$, with $T_{(\varphi,F)}=\iota_u(U_u\varphi)$ and $\theta_{(\varphi,F)}=\mathbf s(F)$ on $V_{u,v}\times Y_v$, such that the following holds. For every finite family $x_i\in X_u$ and $\alpha_i\in\C$ with $\sum_i\alpha_i\lambda_u(x_i)=0$,
\[
\sum_i\alpha_i\,\per\Bigl(T_{x_i}(\xi),\,R_{(y,x_u)}R_z\overline{\theta_{x_i}}\Bigr)=0\qquad\text{for every }\xi\in V_\sigma\text{ and }y\in G_\infty.
\]
For some $x_0\in X_u$ the pairing $\int_{[G]}R_z\overline{\theta_{x_0}}\,P_{T_{x_0}}$ is nonzero.
\end{hypothesis}

\begin{hypothesis}[dictionaries]\label{w2:H-dict}
The space $\mathcal H$ consists of bounded continuous functions on $G(\A)$, left $G(F)$-invariant and right $(K_2,\chi_2)$-equivariant. The column space $\mathcal C_u$ of Definition~\ref{r3:bessel} is the space $\mathcal C$ of \ref{w2:P4} with the norm for which $\mathcal P_u$ is the map $T\mapsto P_T$ of Proposition~\ref{w2:unfold}, and $P_T\in\mathcal H$ for every $T$. The central translation $R_z$ maps $\mathbf s(Y_v)$ onto itself.
\end{hypothesis}

\begin{hypothesis}[rationality]\label{w2:H-rat}
For $\varphi\in V_u$ and $y\in Y$ the number $(B_{u,v}(\varphi\otimes1))(y\otimes1)$ lies in $\sigma_v(K)$ and equals $\sigma_v(B_u(\varphi)(y))$.
\end{hypothesis}

\begin{definition}\label{w2:anchor}
For $x\in X_u$ put $\mu_u(x)=\int_{[G]}R_z\overline{\theta_x}\,P_{T_x}$.
The last condition of Hypothesis~\ref{w2:H-see} asks for a pair $x_0$
with $\mu_u(x_0)\ne0$.
\end{definition}

\begin{lemma}\label{w2:punz}
Assume $\lVert K_{uu}-1\rVert_\HS^2\le e_0\,d^S_u$ and $\dim\iota_u(V_{u,v})>e_0\,d^S_u$. Then there is $\varphi_0\in V_{u,v}$ with $\mathcal P_u\iota_u(U_u\varphi_0)\ne0$.
\end{lemma}

\begin{proof}
For $w\in\mathcal C_u$ one has $\lVert\mathcal P_uw\rVert^2=c_u\langle w,K_{uu}w\rangle$, so $\ker\mathcal P_u=\ker K_{uu}$, because $K_{uu}$ is positive semidefinite. On this kernel $K_{uu}-1$ acts as $-1$, so $\dim\ker\mathcal P_u\le\lVert K_{uu}-1\rVert_\HS^2\le e_0d^S_u$. Since $U_u$ is invertible, $\iota_uU_u(V_{u,v})=\iota_u(V_{u,v})$ has dimension larger than $\dim\ker\mathcal P_u$ and is not contained in it.
\end{proof}

Under the conclusion of Theorem~\ref{w2:S}, $e_0\le\tfrac12$ for large $x$, and $\dim\iota_u(V_{u,v})=d_u\ge(1-2f/T_u)d^S_u>\tfrac12d^S_u$ by Lemma~\ref{r3:ramshare}, since $T_u\ge b>8f$. So Lemma~\ref{w2:punz} applies.

\begin{theorem}\label{w2:Bproof}
Assume \ref{w2:P1} to \ref{w2:P5} for the word $u$, Hypotheses~\ref{w2:H-see}, \ref{w2:H-dict} and \ref{w2:H-rat}. Then \eqref{w2:besselid} holds for $u$ with $\mathcal V$ the unitary operator of $Y_v$ defined by $\mathbf s(\mathcal VF)=R_z\mathbf s(F)$, with $\mathcal U_u=r_u^{-1}U_u$ and with a number $a_u\ne0$. If the hypotheses hold for every $u\in W$ with one $z$, then $\mathcal V$ does not depend on $u$ and Theorem~\ref{w2:B} holds.
\end{theorem}

\begin{proof}
By Lemma~\ref{w2:central} and Hypothesis~\ref{w2:H-dict}, $R_z$ restricts to a unitary operator of $\mathbf s(Y_v)$, and since $\mathbf s$ is a conjugate linear isometry onto its image, $\mathcal V=\mathbf s^{-1}R_z\mathbf s$ is a linear unitary operator of $Y_v$.

Let $\mu_u$ be as in Definition~\ref{w2:anchor}. Every $\theta_x\in\mathcal H$ is bounded, and so is $R_z\overline{\theta_x}$, which is left $G(F)$-invariant and right $(K_2,\overline{\chi_2})$-equivariant by Lemma~\ref{w2:central}. Hypothesis~\ref{w2:H-see} and Corollary~\ref{w2:bottom} show that every linear relation among the values of $\lambda_u$ holds among the values of $\mu_u$. By Hypothesis~\ref{w2:H-see}, $\mu_u(x_0)\ne0$. Lemma~\ref{w2:prop} gives $\kappa_u\ne0$ with $\mu_u=\kappa_u\lambda_u$ on $X_u$. On $V_{u,v}\times Y_v$ this reads
\[
\kappa_u\,(B_{u,v}\varphi)(F)=\int_{[G]}\overline{R_z\mathbf s(F)}\;\mathcal P_u\iota_u(U_u\varphi)=\bigl\langle\mathbf s(\mathcal VF),\mathcal P_u\iota_u(U_u\varphi)\bigr\rangle_{\mathcal H}=r_u\bigl\langle\mathbf s(\mathcal VF),\mathcal P_u\iota_u\mathcal U_u\varphi\bigr\rangle_{\mathcal H}.
\]
So \eqref{w2:besselid} holds with $a_u=r_u/\kappa_u\ne0$, and $\mathcal U_u=r_u^{-1}U_u$ is unitary because it is a linear isometry of a finite dimensional space. The rationality clause is Hypothesis~\ref{w2:H-rat}.
\end{proof}

\begin{remark}[what Theorem~\ref{w2:Bproof} leaves open]\label{w2:rem:Bopen}
Hypothesis~\ref{w2:H-see} combines a definition and a relation. The period map is defined from the period $\mathrm{perCol}$ of the chain of theta lifts on the definite side. The seesaw relation transfers relations among the numbers $c(t)\,\mathrm{perCol}$ to relations among the local Bessel periods of the indefinite theta lifts translated by $z$, with $c(t)=c_{\mathrm{top}}(t)s/c_{\mathrm{bot}}(t)$. The constants of the double step come from the volumes and local scalar factors in Theorem~\ref{w3s:thm:T2}. The Siegel--Weil and doubling identities enter the surjectivity in Hypothesis~\ref{w2:see:H-links}.
\end{remark}

\subsubsection{The chain of seesaw transfers}\label{w2:sec:seesaw}

The transfer follows from the theta tower identities in Hypothesis~\ref{w2:see:H-laws} and the period identifications in Hypothesis~\ref{w2:see:H-links}.

\paragraph{The tower.} Let $m=2n+2$ be the (even) number of lifts of the chain. The construction gives hermitian spaces $V_0,\dots,V_m$ and $T$ over $E$, definite hermitian spaces $X_0,\dots,X_m$ with an orthogonal decomposition $X_m=T\oplus L$, and $h\mapsto\epsilon(h)$ the induced embedding $\mathrm U(T)\to\mathrm U(X_m)$. For $0\le j$ and $\rho\ge0$ with $j+2\rho+2=m$ there is a Bessel datum of rank $\rho$ at level $j$, and $\mathcal B^\rho_j(\phi,f)$ denotes the Bessel period \eqref{w2:per} of a function $f$ on $\mathrm U(V_j)(\A)$ against a function $\phi$ on $\mathrm U(X_j)(\A)$, with the fixed additive character $\psi$. For $e=\sum_ge(g)\,g$ a finitely supported combination of finite adelic points put $\mathrm{tr}(e)f=\sum_ge(g)R_gf$.

For $j\le m-1$ let $\mathcal S_j$ be the space of smooth vectors of the archimedean Fock model at level $j$ and $\mathcal S_T$ the one of $T$. The space $\mathcal C_j$ consists of linear maps $G\colon\mathcal S_j\to\mathcal A(\mathrm U(V_j))$, and $\mathcal A$ denotes the space of automorphic functions. Define $\mathcal C_T$ for $T$ in the same way. The theta lifts of the chain are linear maps
\[
\iota_j\colon\mathcal C_{j+1}\to\mathcal C_j\quad(j\le m-2),\qquad\iota_T\colon\mathcal C_T\to\mathcal C_{m-1},
\]
and $\Theta_j(f;P)$ denotes the theta integral of a function $f$ on $\mathrm U(V_{j+1})$ against the theta kernel attached to a polynomial Fock vector $P$, with $\Theta_T$ the same for $T$. The maps $\mathrm{It}_j\colon\mathcal P_j\otimes\mathcal S_{j+1}\to\mathcal S_j$ and $\mathrm{It}_T\colon\mathcal P_T\otimes\mathcal S_T\to\mathcal S_{m-1}$ are the tensor maps of the Fock models.

The tower is indexed by a set $\mathcal T$ of tuples. For each $s\in\mathcal T$ and each level $j$ the construction gives a finite set $\mathrm{Ch}_j$, combinations $e_{j,\alpha}$ with finite support for $\alpha\in\mathrm{Ch}_j$, a set $\mathrm{Col}_j$ of columns with a subset $\mathrm{Ram}_j$, maps $S_j\colon\mathrm{Col}_j\to\mathrm{Col}_j$ and $K_j\colon\mathrm{Col}_j\to\mathrm{Col}_{j+1}$ preserving the subsets $\mathrm{Ram}$, a set $\mathrm{Fr}_j$ of frames, reading functions $\mathrm{rd}_j(\alpha,\psi',y)$ on $\mathrm U(X_j)(\A)$ for $\psi'\in\mathrm{Col}_j$ and $y\in\mathrm{Fr}_j$, a tail $\mathrm{tail}(\psi',F)\in\C$ for $\psi'\in\mathrm{Col}_m$ and $F\in\mathcal C_T$, and complex numbers $c_I(j),c_V(j),c_K(j),c_P$. The dependence on $s$ is not written. Put
\begin{equation}\label{w2:see:per}
\mathrm{per}^\rho_j(\psi',G;y,v)=\sum_{\alpha\in\mathrm{Ch}_j}\mathcal B^\rho_j\bigl(\mathrm{rd}_j(\alpha,\psi',y),\ \mathrm{tr}(e_{j,\alpha})\,G(v)\bigr)\qquad(G\in\mathcal C_j,\ v\in\mathcal S_j),
\end{equation}
$\mathrm{up}_0\psi'=\psi'$ and $\mathrm{up}_{k+1}\psi'=K_{2k+1}K_{2k}S_{2k}\,\mathrm{up}_k\psi'$ for $\psi'\in\mathrm{Col}_0$, and define the top and bottom constants
\begin{equation}\label{w2:see:const}
W_k=\prod_{l=0}^{k-1}c_V(2l)\,c_K(2l)\,c_K(2l+1),\qquad c_{\mathrm{top}}=W_{n+1}\,c_P,\qquad c_{\mathrm{bot}}=\prod_{j=0}^{m-1}c_I(j).
\end{equation}
The lift $\iota_T$ has the factor $c_I(m-1)$, so $c_{\mathrm{bot}}$ is the product of the constants of all $m$ lifts.

\begin{hypothesis}[laws of the theta tower]\label{w2:see:H-laws}
For every tuple the following hold.
\begin{enumerate}[label=(T\arabic*),leftmargin=*]
\item\label{w2:see:T1} (Lift law.) For $G\in\mathcal C_{j+1}$, $P\in\mathcal P_j$ and $v\in\mathcal S_{j+1}$, $\Theta_j(G(v);P)=c_I(j)\,(\iota_jG)(\mathrm{It}_j(P\otimes v))$, and the same holds for $\iota_T$ with $c_I(m-1)$. All $c_I(j)$ are nonzero.
\item\label{w2:see:T2} (Double-step seesaw identity.) Let $j+2(\sigma+1)+2=m$, $P_1\in\mathcal P_j$, $P_2\in\mathcal P_{j+1}$, and let $a$ and $a'$ be archimedean points of $\mathrm U(V_{j+1})$ and $\mathrm U(V_{j+2})$. There is a bivacuum kernel $C$, depending on $(j,\sigma,P_1,P_2,a,a')$ and not on the tuple, such that for $\psi'\in\mathrm{Ram}_j$, $y$, $G\in\mathcal C_{j+2}$ and $v$, the integrand of $\mathcal B^{\sigma+1}_j$ applied to the double lift $\Theta_j(\Theta_{j+1}(G(v);P_2);P_1)$, integrated first over the Jacobi quotient at level $j+1$ translated by $a$ and then over the Bessel quotient at level $j+2$ translated by $a'$, converges absolutely and equals
\[
c_V(j)\sum_{\beta}\mathcal B^\sigma_{j+2}\bigl(\mathrm{two}_j(\mathrm{rd}_j(\cdot,S_j\psi',y),C)_\beta,\ \mathrm{tr}(e_{j+2,\beta})\,R_{a'}G(v)\bigr),
\]
where $\mathrm{two}_j$ is the bilift of the reading functions through the kernels of levels $j$ and $j+1$ composed with $C$.
\item\label{w2:see:T3} (Kernel law.) For every bivacuum kernel $C$ and every $y$ there are coefficients $a(y')$, finitely many nonzero and independent of the tuple and of $\psi'\in\mathrm{Ram}_j$, with
\[
\mathrm{two}_j(\mathrm{rd}_j(\cdot,\psi',y),C)_\beta=c_K(j)\,c_K(j+1)\sum_{y'}a(y')\,\mathrm{rd}_{j+2}\bigl(\beta,K_{j+1}K_j\psi',y'\bigr).
\]
\item\label{w2:see:T4} (Last-step identity.) For $j=m-2$, $P_1\in\mathcal P_j$, $P_2\in\mathcal P_T$ there is a bivacuum kernel $C$ independent of the tuple with
\[
\mathrm{per}^0_j\bigl(\psi',\Theta_j(\Theta_T(F(w);P_2);P_1);y\bigr)=c_V(j)\,\mathrm{perT}\bigl(\mathrm{two}_j(\mathrm{rd}_j(\cdot,S_j\psi',y),C),F(w)\bigr),
\]
where $\mathrm{perT}(\phi,f)=\sum_\beta\int_{[\mathrm U(T)]}\mathrm{tr}(e_{T,\beta})f(h)\,\phi_\beta(\epsilon(h))\,dh$.
\item\label{w2:see:T5} (Top period.) For $y$ and $w$ there is $b(y,w)\in\C$, independent of the tuple, of $\psi'\in\mathrm{Ram}_m$ and of $F$, with $\mathrm{perT}(\mathrm{rd}_m(\cdot,\psi',y),F(w))=c_P\,b(y,w)\,\mathrm{tail}(\psi',F)$.
\item\label{w2:see:T6} (Unfolding of the double step.) If for every pair $(a,a')$ of archimedean points the iterated integrals in \ref{w2:see:T2} of a finite family of tuples, with coefficients $\alpha_i$, have values $v_i$ with $\sum_i\alpha_iv_i=0$, then $\sum_i\alpha_i\,\mathcal B^{\sigma+1}_j$ of the corresponding double lifts vanishes.
\item\label{w2:see:T7} (Density.) The vectors $\mathrm{It}_j(P_1\otimes\mathrm{It}_{j+1}(P_2\otimes v))$, and the vectors $\mathrm{It}_j(P_1\otimes\mathrm{It}_T(P_2\otimes v))$, span a subspace on which a finite linear relation among the functionals $v\mapsto\mathrm{per}^\rho_j(\psi_i,G_i;y,v)$ implies the same relation on all of $\mathcal S_j$.
\item\label{w2:see:T8} (Archimedean translation.) For an archimedean point $a'$ of $\mathrm U(V_{j})$ there is a linear map $\tau_{a'}$ of $\mathcal S_j$ with $R_{a'}G(v)=G(\tau_{a'}v)$ for every $G\in\mathcal C_j$ and every tuple, and $\mathcal B^\rho_j$ is linear in each argument.
\end{enumerate}
\end{hypothesis}

\begin{proposition}[one double step]\label{w2:see:step}
Assume Hypothesis~\ref{w2:see:H-laws}. Let $j+2(\sigma+1)+2=m$, let $s_1,\dots,s_\ell$ be tuples, $\alpha_i\in\C$, $\psi_i\in\mathrm{Ram}_j$ and $G_i\in\mathcal C_{j+2}$, all for the tuple $s_i$. If
\[
\sum_i\alpha_i\,c_V(j)c_K(j)c_K(j+1)\,\mathrm{per}^\sigma_{j+2}\bigl(K_{j+1}K_jS_j\psi_i,G_i;y,v\bigr)=0\qquad\text{for all }y,v,
\]
then
\[
\sum_i\alpha_i\,c_I(j)c_I(j+1)\,\mathrm{per}^{\sigma+1}_j\bigl(\psi_i,\iota_j\iota_{j+1}G_i;y,v\bigr)=0\qquad\text{for all }y,v.
\]
\end{proposition}

\begin{proof}
Fix $y$. By \ref{w2:see:T7} it suffices to prove the second identity at $v=\mathrm{It}_j(P_1\otimes\mathrm{It}_{j+1}(P_2\otimes w))$. Two applications of \ref{w2:see:T1} give
\[
c_I(j)c_I(j+1)\,(\iota_j\iota_{j+1}G_i)(v)=\Theta_j\bigl(\Theta_{j+1}(G_i(w);P_2);P_1\bigr),
\]
so the left side of the second identity is $\sum_i\alpha_i\,\mathrm{per}^{\sigma+1}_j$ of the double lifts. By \ref{w2:see:T6} it suffices to show, for fixed archimedean points $(a,a')$, that the iterated integrals $v_i$ satisfy $\sum_i\alpha_iv_i=0$. Let $C$ be the kernel of \ref{w2:see:T2} for $(j,\sigma,P_1,P_2,a,a')$. Then \ref{w2:see:T2}, \ref{w2:see:T3} applied to $S_j\psi_i\in\mathrm{Ram}_j$, and \ref{w2:see:T8} give
\[
v_i=c_V(j)c_K(j)c_K(j+1)\sum_{y'}a(y')\,\mathrm{per}^\sigma_{j+2}\bigl(K_{j+1}K_jS_j\psi_i,G_i;y',\tau_{a'}w\bigr),
\]
with $a(y')$ and $\tau_{a'}$ independent of $i$. Hence $\sum_i\alpha_iv_i=\sum_{y'}a(y')\cdot0=0$ by the hypothesis at $(y',\tau_{a'}w)$.
\end{proof}

\begin{proposition}[last step]\label{w2:see:last}
Assume Hypothesis~\ref{w2:see:H-laws} and let $j=m-2$. Let $s_i$, $\alpha_i$, $\psi_i\in\mathrm{Ram}_j$ and $F_i\in\mathcal C_T$ be finitely many data. If
\[
\sum_i\alpha_i\,c_V(j)c_K(j)c_K(j+1)\,c_P\ \mathrm{tail}\bigl(K_{j+1}K_jS_j\psi_i,F_i\bigr)=0,
\]
then $\sum_i\alpha_i\,c_I(j)c_I(j+1)\,\mathrm{per}^0_j(\psi_i,\iota_j\iota_TF_i;y,v)=0$ for all $y,v$.
\end{proposition}

\begin{proof}
By \ref{w2:see:T7} take $v=\mathrm{It}_j(P_1\otimes\mathrm{It}_T(P_2\otimes w))$. By \ref{w2:see:T1} the $i$-th term is
\[
 \alpha_i\,\mathrm{per}^0_j(\psi_i,\Theta_j(\Theta_T(F_i(w);P_2);P_1);y).
\]
By \ref{w2:see:T4}, \ref{w2:see:T3} and \ref{w2:see:T5} this equals
\[
\sum_{y'}a(y')\,b(y',w)\,\alpha_i\,c_V(j)c_K(j)c_K(j+1)\,c_P\ \mathrm{tail}\bigl(K_{j+1}K_jS_j\psi_i,F_i\bigr).
\]
The sum over $i$ vanishes for each $y'$ by hypothesis.
\end{proof}

\begin{theorem}[transfer along the tower]\label{w2:see:main}
Assume Hypothesis~\ref{w2:see:H-laws}. Let $s_i$, $\alpha_i$, $\psi_i\in\mathrm{Ram}_0$ and $F_i\in\mathcal C_T$ be finitely many data. If
\[
\sum_i\alpha_i\,c_{\mathrm{top}}(s_i)\,\mathrm{tail}\bigl(\mathrm{up}_{n+1}\psi_i,F_i\bigr)=0,
\]
then for all $y\in\mathrm{Fr}_0$ and $v\in\mathcal S_0$
\[
\sum_i\alpha_i\,c_{\mathrm{bot}}(s_i)\,\mathrm{per}^n_0\bigl(\psi_i,\ \iota_0\iota_1\cdots\iota_{m-2}\iota_TF_i;y,v\bigr)=0 .
\]
\end{theorem}

\begin{proof}
Since $S$ and $K$ preserve the subsets $\mathrm{Ram}$, $\mathrm{up}_k\psi_i\in\mathrm{Ram}_{2k}$ for $2k\le m$. Since $\mathrm{up}_{n+1}\psi_i=K_{2n+1}K_{2n}S_{2n}\,\mathrm{up}_n\psi_i$ and $c_{\mathrm{top}}=W_n\,c_V(2n)c_K(2n)c_K(2n+1)\,c_P$, Proposition~\ref{w2:see:last} with $j=2n$, columns $\mathrm{up}_n\psi_i$ and coefficients $\alpha_iW_n$ gives
\begin{equation}\label{w2:see:start}
\sum_i\alpha_iW_n\,c_I(2n)c_I(2n+1)\,\mathrm{per}^0_{2n}\bigl(\mathrm{up}_n\psi_i,\iota_{2n}\iota_TF_i;y,v\bigr)=0 .
\end{equation}
We show by induction on $k$ the following statement $(\mathrm A_k)$. If $2k+2\rho+2=m$, $G_i\in\mathcal C_{2k}$ and $\sum_i\beta_iW_k\,\mathrm{per}^\rho_{2k}(\mathrm{up}_k\psi_i,G_i;y,v)=0$ for all $y,v$, then $\sum_i\beta_i\prod_{j<2k}c_I(j)\,\mathrm{per}^{k+\rho}_0(\psi_i,\iota_0\cdots\iota_{2k-1}G_i;y,v)=0$ for all $y,v$. For $k=0$ there is nothing to prove. If $(\mathrm A_k)$ holds and the hypothesis of $(\mathrm A_{k+1})$ holds with $\rho$, then $W_{k+1}=W_k\,c_V(2k)c_K(2k)c_K(2k+1)$, and Proposition~\ref{w2:see:step} at $j=2k$, $\sigma=\rho$ with coefficients $\beta_iW_k$ gives $\sum_i\beta_iW_k\,c_I(2k)c_I(2k+1)\,\mathrm{per}^{\rho+1}_{2k}(\mathrm{up}_k\psi_i,\iota_{2k}\iota_{2k+1}G_i;y,v)=0$. Then $(\mathrm A_k)$ with $\rho+1$ and coefficients $\beta_ic_I(2k)c_I(2k+1)$ gives $(\mathrm A_{k+1})$. Finally $(\mathrm A_n)$ with $\rho=0$, $\beta_i=\alpha_ic_I(2n)c_I(2n+1)$ and $G_i=\iota_{2n}\iota_TF_i$ applies to \eqref{w2:see:start}, and $\prod_{j<2n}c_I(j)\cdot c_I(2n)c_I(2n+1)=c_{\mathrm{bot}}$ because $m=2n+2$.
\end{proof}

\paragraph{The word $u$.} Fix a word $u$. The tower attached to $u$ has its tuples indexed by the set of tuples $t$ of the chain datum, with the base tuple $\mu_0$, and $\mathrm{perCol}_t(\varphi,F)$ denotes the period of the chain on the definite side at the tuple $t$, for $\varphi\in V_{u,v}$ and a $K$-rational vector $F$ of the good part at $t$. By the definition of the period map, $(B_{u,v}\varphi)(\iota_0(q))=\mathrm{perCol}_{\mu_0}(\varphi,q)$ for $q$ in the good part at $\mu_0$, and these vectors $\iota_0(q)$ span $Y_v$. Let $\mathcal R_0$ be the space of row forms at level $0$, the span of the theta lifts of level $0$ that the construction uses, and for $F$ as above let $\mathcal R_F\in\mathcal R_0$ be the row form of $F$. Let $b\in\mathrm U(V_0)(\A)$ be the translation of the Bessel datum of the word, with $z$ and $x_u$ as in Hypothesis~\ref{w2:H-see} related to it by the pinning relation of the construction, and write
\[
\mathfrak b_u(\varphi,\mathcal R;y,v)=\mathcal B^n_0\bigl(\mathrm{rd}^u_0(\varphi,y),\ R_b\,\mathcal R(v)\bigr)
\]
for the local Bessel period at level $0$ of the translated row form $\mathcal R$ against the column $\varphi$.

\begin{hypothesis}[dictionaries of the tower]\label{w2:see:H-links}
There are maps $\mathrm{col}_0(t,\cdot)\colon V_{u,v}\to\mathrm{Col}_0$ and $\mathrm{col}_m\colon\mathrm{Col}_0\to\mathrm{Col}_m$ such that the following hold for every tuple $t$.
\begin{enumerate}[label=(D\arabic*),leftmargin=*]
\item\label{w2:see:D1} $\mathrm{col}_0(t,\varphi)\in\mathrm{Ram}_0$ and $\mathrm{per}^n_0(\mathrm{col}_0(t,\varphi),G;y,v)=\mathfrak b_u(\varphi,G;y,v)$ for every $G\in\mathcal C_0$.
\item\label{w2:see:D2} If $G\in\mathcal C_T$ is the leading map of $F$, then $\mathrm{tail}(\mathrm{col}_m\mathrm{col}_0(t,\varphi),G)=\mathrm{perCol}_t(\varphi,F)$.
\item\label{w2:see:D3} For every $F$ there is $G\in\mathcal C_T$, the leading map of $F$, with $\iota_0\cdots\iota_{m-2}\iota_TG=\mathcal R_F$.
\item\label{w2:see:D4} Every $\mathcal R\in\mathcal R_0$ equals $\iota_0\cdots\iota_{m-2}\iota_TG$ for some $G\in\mathcal C_T$ at the base tuple $\mu_0$.
\item\label{w2:see:D5} There are a linear bijection $S$ of $V_{u,v}$, $r>0$ and $\varsigma\ne0$ with $\lVert S\varphi\rVert=r\lVert\varphi\rVert$ and $\mathrm{tail}(\mathrm{up}_{n+1}\mathrm{col}_0(t,\varphi),G)=\varsigma\,\mathrm{tail}(\mathrm{col}_m\mathrm{col}_0(t,S\varphi),G)$ for all $t$, $\varphi$ and $G$.
\item\label{w2:see:D6} $c_{\mathrm{top}}(t)\ne0$.
\end{enumerate}
\end{hypothesis}

\begin{corollary}[the seesaw relation]\label{w2:see:cor}
Assume Hypotheses~\ref{w2:see:H-laws} and~\ref{w2:see:H-links}. Put $U_u=S^{-1}$, $r_u=r^{-1}$ and
\begin{equation}\label{w2:see:ct}
c(t)=\frac{c_{\mathrm{top}}(t)\,\varsigma}{c_{\mathrm{bot}}(t)}\ne0,\qquad p_{\mathcal R}(\varphi_0)=\frac{c_{\mathrm{top}}(\mu_0)}{c_{\mathrm{bot}}(\mu_0)}\,\mathrm{tail}\bigl(\mathrm{up}_{n+1}\mathrm{col}_0(\mu_0,S^{-1}\varphi_0),G_{\mathcal R}\bigr)
\end{equation}
for $\mathcal R\in\mathcal R_0$, with $G_{\mathcal R}$ chosen by \ref{w2:see:D4}. Then $\lVert U_u\varphi\rVert=r_u\lVert\varphi\rVert$, and for finitely many data $(\varphi_{0,i},\mathcal R_i)$, $(t_k,\varphi_k,F_k)$ and coefficients $\alpha_i,\alpha'_k$ with
\[
\sum_i\alpha_i\,p_{\mathcal R_i}(\varphi_{0,i})+\sum_k\alpha'_k\,c(t_k)\,\mathrm{perCol}_{t_k}(\varphi_k,F_k)=0
\]
one has, for all $y$ and $v$,
\[
\sum_i\alpha_i\,\mathfrak b_u(U_u\varphi_{0,i},\mathcal R_i;y,v)+\sum_k\alpha'_k\,\mathfrak b_u(U_u\varphi_k,\mathcal R_{F_k};y,v)=0 .
\]
\end{corollary}

\begin{proof}
The norm statement follows from $\lVert S\psi\rVert=r\lVert\psi\rVert$ with $\psi=S^{-1}\varphi$, and $c(t)\ne0$ follows from \ref{w2:see:D6}, $\varsigma\ne0$ and \ref{w2:see:T1}. For each $k$ choose $G_k$ by \ref{w2:see:D3}. Apply Theorem~\ref{w2:see:main} to the family consisting of the tuples $\mu_0$ (indices $i$) and $t_k$ (indices $k$), the columns $\mathrm{col}_0(\mu_0,S^{-1}\varphi_{0,i})$ and $\mathrm{col}_0(t_k,S^{-1}\varphi_k)$, the elements $G_{\mathcal R_i}$ and $G_k$, and the coefficients $\alpha_i/c_{\mathrm{bot}}(\mu_0)$ and $\alpha'_k/c_{\mathrm{bot}}(t_k)$. Its hypothesis is the displayed relation, because by \ref{w2:see:D5} and \ref{w2:see:D2}
\[
\frac{c_{\mathrm{top}}(t_k)}{c_{\mathrm{bot}}(t_k)}\,\mathrm{tail}\bigl(\mathrm{up}_{n+1}\mathrm{col}_0(t_k,S^{-1}\varphi_k),G_k\bigr)=\frac{c_{\mathrm{top}}(t_k)\,\varsigma}{c_{\mathrm{bot}}(t_k)}\,\mathrm{perCol}_{t_k}(\varphi_k,F_k)=c(t_k)\,\mathrm{perCol}_{t_k}(\varphi_k,F_k),
\]
and the $i$-terms are $\alpha_ip_{\mathcal R_i}(\varphi_{0,i})$ by definition. Its conclusion, after the factors $c_{\mathrm{bot}}$ cancel, is $\sum_i\alpha_i\,\mathrm{per}^n_0(\mathrm{col}_0(\mu_0,S^{-1}\varphi_{0,i}),\mathcal R_i;y,v)+\sum_k\alpha'_k\,\mathrm{per}^n_0(\mathrm{col}_0(t_k,S^{-1}\varphi_k),\mathcal R_{F_k};y,v)=0$, using \ref{w2:see:D3} and \ref{w2:see:D4}, and \ref{w2:see:D1} turns this into the assertion.
\end{proof}

\begin{remark}[passage to Hypothesis~\ref{w2:H-see}]\label{w2:see:rem}
Take $X_u$ to consist of the pairs $(\varphi,F)\in V_{u,v}\times Y_v$ and the pairs $(\varphi_0,\mathcal R)$ with $\mathcal R\in\mathcal R_0$, with $\lambda_u(\varphi,F)=(B_{u,v}\varphi)(F)$ and $\lambda_u(\varphi_0,\mathcal R)=p_{\mathcal R}(\varphi_0)/c(\mu_0)$. Only the base tuple $t=\mu_0$ is needed. A relation among the values of $\lambda_u$ with complex $F$ is a relation with $K$-rational $F$ after expanding $F$ in a $K$-basis of the good part, because both sides are linear in $F$, and multiplying the relation by $c(\mu_0)$ turns it into the hypothesis of Corollary~\ref{w2:see:cor} with the same coefficients. The conclusion of the corollary is the conclusion of Hypothesis~\ref{w2:H-see}, once the local Bessel period $\mathfrak b_u(U_u\varphi,\mathcal R;y,v)$ is identified with $\mathrm{per}(T_x(\xi),R_{(y,x_u)}R_z\overline{\theta_x})$ through the pinning relation between $b$, $z$ and $x_u$ and the row dictionary $\mathcal R_F\leftrightarrow\mathbf s(F)$ of Hypothesis~\ref{w2:H-dict}. Lemma~\ref{w2:prop} uses only relations with one or two terms, so Hypothesis~\ref{w2:H-see} may be weakened to such relations without changing Theorem~\ref{w2:Bproof}. The anchor of Lemma~\ref{w2:anchor} may be any $x_0=(\varphi_0,\mathcal R)$ with $\mu_u(x_0)\ne0$. The exact values of $c_{\mathrm{top}}$, $c_{\mathrm{bot}}$ and $\varsigma$ do not enter Theorem~\ref{w2:B}, because the number $a_u=r_u/\kappa_u$ absorbs them. Only $c(\mu_0)\ne0$ is used there, while the unit property of $c(t)/c(\mu_0)$ at the primes of the chain is used for the integral structure of the period maps.
\end{remark}

\begin{lemma}[rationality]\label{w2:see:rat}
Assume that $B_u(\varphi)(y)=\mathrm{Pt}_{\mu_0}(\varphi,q_y)$ for $\varphi\in V_u$ and $y\in Y$, where $\mathrm{Pt}_{\mu_0}$ is the $K$-bilinear period of the chain on the definite side at the base tuple and $q_y$ is the $K$-rational vector of the good part with $\iota_0(q_y)=y$. Assume also $\mathrm{perCol}_{\mu_0}(\varphi\otimes1,q\otimes1)=\sigma_v(\mathrm{Pt}_{\mu_0}(\varphi,q))$ and $(B_{u,v}\varphi)(\iota_0q)=\mathrm{perCol}_{\mu_0}(\varphi,q)$. Then Hypothesis~\ref{w2:H-rat} holds.
\end{lemma}

\begin{proof}
For $\varphi\in V_u$ and $y\in Y$, $(B_{u,v}(\varphi\otimes1))(y\otimes1)=\mathrm{perCol}_{\mu_0}(\varphi\otimes1,q_y\otimes1)=\sigma_v(\mathrm{Pt}_{\mu_0}(\varphi,q_y))=\sigma_v(B_u(\varphi)(y))$, which lies in $\sigma_v(K)$.
\end{proof}

\paragraph{Rationality on the selected record.}\label{rat:selected}
We give the descent producing this identity in the original argument.
Fix one selected record, its coefficient field $K$, its embedding
$\sigma:K\hookrightarrow\mathbb C$, its base tuple $\mu_0$ and a word $u$.
Write $N_{u,K}$ and $R_K$ for its rational column and row spaces and
$j_N,j_R$ for their $\sigma$-semilinear maps to the complex spaces.
The images of the chosen bases are complex bases. Let $B^K_u$
be its rational bilinear pairing and $\mathcal B_u$ its complex pairing.
For these data the assertion is
\begin{equation}\label{rat:base}
 \mathcal B_u(j_N\varphi,j_R y)=\sigma\big(B^K_u(\varphi,y)\big)
 \qquad(\varphi\in N_{u,K},\ y\in R_K).
\end{equation}

Choose the rational column coordinates of $\varphi$. On the row side
let $e_R:K^d\longrightarrow R_K$ be the coordinate isomorphism and let
$p:K^n\to K^d$ be the composite of the good-part projection and the
row-coordinate map. It is onto by surjectivity of the first projection
and bijectivity on the good part. Let $A(F)$ be the corresponding
complex level row. Its values span the degree-zero space $V^0$,
since the level operator has precisely this image and fixes its image.
Put $\beta(y')=B^K_u(\varphi,e_Ry')$. The row comparison on this
realization gives
\[
 \iota A(F)=j_R e_R(pF),\qquad
 \mathrm{perCol}_{\mu_0}(j_N\varphi,A(F))=\sigma(\beta(pF)).
\]
For the second equality, the comparison of a rational definite period
with the complex column gives $\sigma(\mathrm{Pt}_{\mu_0})$ on the left.
The definite period factors as $\zeta(u,\mu_0)$ times the pairing on the
projected row, and $\zeta(u,\mu_0)=1$. Transport by the rational column
and row isomorphisms gives $\beta(pF)$. At other tuples the source
specifies $\zeta(u,t)\ne0$ and valuation zero at the designated places.
Identifying this scalar with $c(t)/c(\mu_0)$ requires a separate
comparison of period normalizations.

Extend $\beta e_R^{-1}$ to the complex linear functional $\Lambda$
by prescribing its $\sigma$-embedded values on the chosen row basis.
The last two identities show that $\mathrm{perCol}_{\mu_0}(j_N\varphi,-)$
and $\Lambda\iota$ agree on every $A(F)$, hence on $V^0$. Put
$U=V^0\cap\ker\iota$ and $Q=V^0\cap U^\perp$. On $Q$, $\iota$ is
injective. The section at $\mu_0$ first projects onto $\iota(Q)$ and
then inverts $\iota|_Q$. Given $y=e_R y'$, take $F$ with $pF=y'$.
Now $q=A(F)-\operatorname{proj}_U A(F)$ belongs to $Q$ and
$\iota q=j_R y$. Thus the section takes $j_Ry$ to $q$.
The defining tie for $\mathcal B_u$ evaluates the period on this
section. Its value is $\Lambda(j_R y)=\sigma(B^K_u(\varphi,y))$.
This proves \eqref{rat:base}. It also proves vanishing on $U$ for the
columns just treated. For a general column, extending the good-part
identity to all of $V^0$ requires vanishing on $U$, by the same
orthogonal decomposition.

In a selected family this argument holds at every index $i$ once $x$
exceeds a common threshold. Take the maximum of the threshold for
the row facts, the selection threshold, and the initial scale plus one.
The realization supplies the assertion that all tuples are used, which
allows the definite-period comparison to be applied to these rows.
For condition (d) of Theorem~\ref{r3:A3} one place suffices for the
whole family. Take the place of $\sigma$ and the same chosen record for
all words. Use its rational pairing for $B_u$. For bases $(b_i)$ of
$N_{u,K}$ and $(c_l)$ of $R_K$, equation~\eqref{rat:base} gives the
matrix of the analytic pairing as
$(\sigma(B_u(b_i)(c_l)))_{il}$. Bilinearity now identifies the two
maps on the full complex spaces and gives Hypothesis~\ref{w2:H-rat}
for this realization at the place required by (d). If the seesaw
calculation uses different coordinates, transport its columns and rows
by their respective basis isomorphisms. In common coordinates the
remaining comparison is equality of the two resulting matrices, or
equivalently equality on every pair $(b_i,c_l)$. The height comparison
in Section~\ref{r4:sec:R4} still uses the metrics at all infinite places.



\subsubsection{The double-step seesaw identity}\label{w3s:sec}

This subsection proves the double-step identity (T2) of Hypothesis~\ref{w2:see:H-laws}. The identity is an iterated integral over two quotients of finite adelic groups. We show that the integral collapses to an integral over the finitely many places of the tuple, with a volume factor $\kappa_j$, and we identify the result with the bilift of the reading functions through one kernel $C$ that does not depend on the tuple. The analytic unfolding over the archimedean and automorphic variables is not part of (T2). It is item (T6), where rapid decay of cusp forms enters.

\paragraph{Quotient integrals.} Let $\mathbf G$ be a second countable locally compact group with a left Haar measure $\mu$, and let $R\subset\mathbf G$ be a closed subgroup with a left Haar measure $\rho$. We say that a function $f\colon\mathbf G\to\C$ has the \emph{absolutely convergent quotient integral} $z$ over $R\backslash\mathbf G$, and write $\int^{\mathrm{abs}}_{R\backslash\mathbf G}f=z$, if there is $F\in L^1(\mathbf G,\mu)$ such that for every $g\in\mathbf G$ the function $r\mapsto F(rg)$ is $\rho$-integrable, $f(g)=\int_RF(rg)\,d\rho(r)$, and $z=\int_{\mathbf G}F\,d\mu$. For two such pairs $(\mathbf G_1,R_1)$, $(\mathbf G_2,R_2)$ and a function $I$ on $\mathbf G_1\times\mathbf G_2$, we say that $I$ has the \emph{absolutely convergent iterated quotient integral} $z$ if there is $f_1\colon\mathbf G_1\to\C$ with $\int^{\mathrm{abs}}_{R_1\backslash\mathbf G_1}f_1=z$ and $\int^{\mathrm{abs}}_{R_2\backslash\mathbf G_2}I(q_1,\cdot)=f_1(q_1)$ for every $q_1\in\mathbf G_1$. These are the formal notions \texttt{IsQuotIntL1} and \texttt{Iter2}.

\begin{lemma}[uniqueness]\label{w3s:lem:uniq}
Assume that $\mathbf G$ and $R$ are unimodular. If $\int^{\mathrm{abs}}_{R\backslash\mathbf G}f=z$ and $\int^{\mathrm{abs}}_{R\backslash\mathbf G}f=z'$, then $z=z'$. The value is linear in $f$, it vanishes for $f=0$, and $\int^{\mathrm{abs}}_{R\backslash\mathbf G}R_bf=\int^{\mathrm{abs}}_{R\backslash\mathbf G}f$ for every $b\in\mathbf G$, where $R_bf(g)=f(gb)$.
\end{lemma}

\begin{proof}
Since the modular functions of $\mathbf G$ and $R$ are both trivial, there is a right $\mathbf G$-invariant Radon measure $\dot\mu$ on $R\backslash\mathbf G$ with $\int_{\mathbf G}F\,d\mu=\int_{R\backslash\mathbf G}\int_RF(rg)\,d\rho(r)\,d\dot\mu(g)$ for every $F\in L^1(\mathbf G)$ \cite[\S2.6]{w2folland}. Hence $z=\int_{R\backslash\mathbf G}f\,d\dot\mu$ depends only on $f$, and linearity follows. If $F$ is a witness for $f$, then $R_bF$ is a witness for $R_bf$, and $\int R_bF\,d\mu=\int F\,d\mu$ because $\mu$ is right invariant.
\end{proof}

\begin{lemma}[collapse to a factor]\label{w3s:lem:collapse}
Let $\mathbf G$ be a second countable locally compact unimodular group, and let $G_1,G_2\subset\mathbf G$ be closed unimodular subgroups whose elements commute with each other and for which $(a,b)\mapsto ab$ is a homeomorphism $G_1\times G_2\to\mathbf G$. Let $R\subset\mathbf G$ be a closed unimodular subgroup such that $ab\in R$ with $a\in G_1$, $b\in G_2$ implies $a,b\in R$. Put $R_i=R\cap G_i$. Let $K_i\subset G_i$ be compact open subgroups and $K=K_1K_2$. Let $\mu,\rho,\mu_1,\rho_1$ be Haar measures on $\mathbf G,R,G_1,R_1$ with $\mu_1(K_1)=1$ and $\rho_1(R_1\cap K_1)=1$. Let $f\colon\mathbf G\to\C$ satisfy
\begin{enumerate}[label=(\roman*),leftmargin=*]
\item $f(rg)=f(g)$ for $r\in R$ and $g\in\mathbf G$;
\item $f(gk)=f(g)$ for $k\in K_2$ and $g\in\mathbf G$;
\item $f(g)\ne0$ implies $g\in RG_1K_2$.
\end{enumerate}
If $\int^{\mathrm{abs}}_{R_1\backslash G_1}f|_{G_1}=z$, then
\[
\int^{\mathrm{abs}}_{R\backslash\mathbf G}f=\frac{\mu(K)}{\rho(R\cap K)}\,z .
\]
\end{lemma}

\begin{proof}
The product map is a continuous homomorphism because $G_1$ and $G_2$ commute, so it is an isomorphism of topological groups $G_1\times G_2\cong\mathbf G$. By the hypothesis on $R$, it restricts to an isomorphism $R_1\times R_2\cong R$, and $R_1,R_2$ are closed. The modular function of $R_1\times R_2$ is the product of those of $R_1$ and $R_2$, so $R_1$ and $R_2$ are unimodular. By uniqueness of Haar measure there are Haar measures $\mu_2$ on $G_2$ and $\rho_2$ on $R_2$ with $\mu=\mu_1\otimes\mu_2$ and $\rho=\rho_1\otimes\rho_2$ under these isomorphisms. The group $R_2\cap K_2$ is compact and open in $R_2$, so $0<\rho_2(R_2\cap K_2)<\infty$, and likewise $0<\mu_2(K_2)<\infty$.

\emph{Step 1.} We show $f(ab)=f(a)\,\mathbf 1_{R_2K_2}(b)$ for $a\in G_1$ and $b\in G_2$. If $b=r_2k$ with $r_2\in R_2$ and $k\in K_2$, then $ab=r_2ak$, so (ii) and (i) give $f(ab)=f(r_2a)=f(a)$. If $f(ab)\ne0$, then (iii) gives $ab=r_1r_2xk$ with $r_i\in R_i$, $x\in G_1$, $k\in K_2$, that is $ab=(r_1x)(r_2k)$, and uniqueness of the decomposition gives $b=r_2k\in R_2K_2$.

\emph{Step 2.} Let $F_1\in L^1(G_1)$ be a witness for $f|_{G_1}$, and put $F(ab)=F_1(a)\,\mathbf 1_{K_2}(b)/\rho_2(R_2\cap K_2)$. Then $F\in L^1(\mathbf G)$ and $\int_{\mathbf G}F\,d\mu=z\,\mu_2(K_2)/\rho_2(R_2\cap K_2)$. For $g=ab$ and $r=r_1r_2$ we have $F(rg)=F_1(r_1a)\mathbf 1_{K_2}(r_2b)/\rho_2(R_2\cap K_2)$. The second factor is the indicator of the set $R_2\cap K_2b^{-1}$, of measure at most $\rho_2(R_2\cap K_2)$, so $r\mapsto F(rg)$ is $\rho$-integrable and Fubini's theorem gives
\[
\int_RF(rg)\,d\rho(r)=f(a)\,\frac{\rho_2(R_2\cap K_2b^{-1})}{\rho_2(R_2\cap K_2)} .
\]
If $b=r_0k\in R_2K_2$, then $R_2\cap K_2b^{-1}=(R_2\cap K_2)r_0^{-1}$, which has measure $\rho_2(R_2\cap K_2)$ because $\rho_2$ is right invariant. If $R_2\cap K_2b^{-1}$ is nonempty, then $b\in R_2K_2$. So the right side equals $f(a)\mathbf 1_{R_2K_2}(b)$, which is $f(ab)$ by Step 1. Hence $\int^{\mathrm{abs}}_{R\backslash\mathbf G}f=z\,\mu_2(K_2)/\rho_2(R_2\cap K_2)$.

\emph{Step 3.} By the hypothesis on $R$, $R\cap K=(R_1\cap K_1)(R_2\cap K_2)$. Hence $\mu(K)=\mu_1(K_1)\mu_2(K_2)=\mu_2(K_2)$ and $\rho(R\cap K)=\rho_1(R_1\cap K_1)\rho_2(R_2\cap K_2)=\rho_2(R_2\cap K_2)$.
\end{proof}

\paragraph{The double-step integrand.} We use the notation of Subsection~\ref{w2:sec:seesaw}. For $0\le i<m$ the pair $(\mathrm U(V_i),\mathrm U(V_{i+1}))$ acts on functions on $W_i(\A)$, $W_i=\Hom(V_{i+1},V_i)$, by the Weil representation $\omega_i$ of the tuple attached to $\psi$ and to the splitting characters of the tuple. The pairs $(\mathrm U(X_i),\mathrm U(X_{i+1}))$ of definite spaces act on functions on $W'_i(\A)$ in the same way, by $\omega'_i$. For a finitely supported function $e$ on $W_i(\Af)$ put $\mathrm{comb}(e)\Phi=\sum_ue(u)\,\mathrm{heis}(u)\Phi$, where $\mathrm{heis}(u)$ is the Heisenberg translation by $u$ with the character $\psi$. We say that $e$ is \emph{supported at} a set $P$ of places if every $u$ in its support has trivial components outside $P$. For self-dual lattices $N_i\subset V_i$ and a set $P$ of finite places, a theta function $\Phi$ on $W_i(\A)$ is \emph{vacuum at} $P$ if it is invariant under the Heisenberg action of the lattice $\Hom(N_{i+1},N_i)\otimes\prod_{p\in P}\cO_p$, so that its component at $p\in P$ is a multiple of the characteristic function of that lattice. For finitely supported functions $e$ on $W'_i(\Af)$ and $e'$ on $W'_{i+1}(\Af)$ and a function $C$ on $W'_i(\A)\times W'_{i+1}(\A)$ write $\mathrm{biComb}(e,e')C(u,v)=\mathrm{comb}(e)\bigl(u'\mapsto\mathrm{comb}(e')C(u',\cdot)(v)\bigr)(u)$.

Fix $j$ and $\sigma$ with $j+2(\sigma+1)+2=m$. The Bessel datum of rank $\sigma+1$ at level $j$, the Fourier--Jacobi datum of rank $\sigma+1$ at level $j+1$ and the Bessel datum of rank $\sigma$ at level $j+2$ give closed subgroups of $\mathrm U(V_{j+1})(\A)$ and $\mathrm U(V_{j+2})(\A)$. Their intersections with the finite adelic groups $\mathbf G_{j+1}=\mathrm U(V_{j+1})(\Af)$ and $\mathbf G_{j+2}=\mathrm U(V_{j+2})(\Af)$ are denoted $R_J$ and $R_B$, with Haar measures $\rho_J$, $\rho_B$, and $\mu_{j+1},\mu_{j+2}$ are the Haar measures of $\mathbf G_{j+1},\mathbf G_{j+2}$ fixed by the construction. For theta functions $\Phi_1$ on $W_j(\A)$ and $\Phi_2$ on $W_{j+1}(\A)$, a function $\phi$ on $\mathrm U(X_j)(\A)$ and points $g,h_1,h_2$ of $\mathrm U(V_j)(\A)$, $\mathrm U(V_{j+1})(\A)$, $\mathrm U(V_{j+2})(\A)$, let
\[
\mathbf H(\Phi_1,\Phi_2;\phi;g,h_1,h_2)=\mathrm{hand}_{\mathrm{FB}}\Bigl(\mathrm{hand}_{\mathrm{BF}}\bigl(\phi,\mathrm{flat}_{\mathrm{BF}}(\omega_j(g,h_1)\Phi_1)\bigr),\ \mathrm{flat}_{\mathrm{FB}}(\omega_{j+1}(h_1,h_2)\Phi_2)\Bigr)
\]
be the function on $\mathrm U(X_{j+2})(\A)$ obtained by the Bessel-to-Fourier--Jacobi step at level $j$ followed by the Fourier--Jacobi-to-Bessel step at level $j+1$ (formal \texttt{Chain.hand}). Here $\mathrm{flat}_{\mathrm{BF}}$ and $\mathrm{flat}_{\mathrm{FB}}$ are the partial integrations of a theta function along the unipotent directions of the two flags, taken with Haar measures of covolume one, $\mathrm{hand}_{\mathrm{BF}}$ integrates $\phi$ over $[\mathrm U(X_j)]$ against the restricted kernel, and $\mathrm{hand}_{\mathrm{FB}}$ contracts the result with the second kernel. The map $\mathbf H$ is bilinear in $(\Phi_1,\Phi_2)$.

The contracted kernel of two theta functions $\Psi_1,\Psi_2$ is the function $\mathrm{ctr}(\Psi_1,\Psi_2)$ on $W'_j(\A)\times W'_{j+1}(\A)$ given by the integral over the link variable of the product $\mathrm{flat}_{\mathrm{FB}}\Psi_2\cdot\mathrm{flat}_{\mathrm{BF}}\Psi_1$ (formal \texttt{Comm.ctr}). For a function $C$ on $W'_j(\A)\times W'_{j+1}(\A)$ and a function $\phi$ on $\mathrm U(X_j)(\A)$ the bilift is
\[
\mathrm{bl}_j(\phi,C)(h_2)=\int_{[\mathrm U(X_{j+1})]}\int_{[\mathrm U(X_j)]}\phi(h_0)\,\Bigl(\omega'_{j+1}(h_1,h_2)\bigl(v\mapsto\omega'_j(h_0,h_1)C(\cdot,v)(0)\bigr)\Bigr)(0)\,dh_0\,dh_1,
\]
and for a family $\phi=(\phi_\alpha)_{\alpha\in\mathrm{Ch}_j}$ the double kernel lift of Hypothesis~\ref{w2:see:H-laws} is $\mathrm{two}_j(\phi,C)_\beta=\sum_{\alpha,\gamma}\mathrm{bl}_j\bigl(\phi_\alpha,\mathrm{biComb}(\mathrm{dd}_j(\alpha,\gamma),\mathrm{dd}_{j+1}(\gamma,\beta))C\bigr)$ for $\beta\in\mathrm{Ch}_{j+2}$, where $\mathrm{dd}_i$ are the kernel combinations of the tuple. A \emph{bivacuum kernel} is a continuous $C$ that is a theta function vacuum at every finite place in each variable, and it is \emph{real} if it is moreover smooth in the archimedean directions of both variables.

For a tuple $s$, archimedean points $a\in\mathrm U(V_{j+1})_\infty$, $a'\in\mathrm U(V_{j+2})_\infty$, polynomial Fock vectors $P_1,P_2$, a column $\psi'\in\mathrm{Ram}_j$, a frame $y$, $G\in\mathcal C_{j+2}$ and $v\in\mathcal S_{j+2}$, the \emph{double-step integrand} is the function on $\mathbf G_{j+1}\times\mathbf G_{j+2}$
\begin{equation}\label{w3s:integrand}
I_s(q_1,q_2)=\sum_{\alpha\in\mathrm{Ch}_j}\sum_ge_{j,\alpha}(g)\,\mathcal B^\sigma_{j+2}\Bigl(\mathbf H\bigl(\mathrm{pt}_j(P_1),\mathrm{pt}_{j+1}(P_2);\mathrm{rd}_j(\alpha,\psi',y);g,aq_1,a'q_2\bigr),\ R_{a'q_2}G(v)\Bigr),
\end{equation}
where $\mathrm{pt}_i(P)$ is the theta function of the tuple at level $i$ attached to $P$ (formal \texttt{Tup.integrand}). It is the integrand of $\mathcal B^{\sigma+1}_j$ on the double lift after the two unfoldings, as a function of the two finite adelic variables.

\begin{hypothesis}[general laws of the Weil representation]\label{w3s:hyp:gen}
\begin{enumerate}[label=(W\arabic*),leftmargin=*]
\item\label{w3s:W1} (Disjoint places.) Heisenberg combinations supported at disjoint sets of places commute. A combination supported at $P$ maps theta functions vacuum at a set $P'$ disjoint from $P$ to theta functions vacuum at $P'$, and it commutes with $\omega_i(g,h)$ for $g,h$ with trivial components at $P$. A theta function is vacuum at $P\cup P'$ if and only if it is vacuum at $P$ and at $P'$.
\item\label{w3s:W2} (Normalized steps.) The maps $\mathrm{flat}_{\mathrm{BF}}$ and $\mathrm{flat}_{\mathrm{FB}}$ on theta functions depend only on the Bessel and Fourier--Jacobi data and not on the remaining data of the step of the tuple. Consequently $\mathrm{ctr}(\Psi_1,\Psi_2)$ does not depend on the tuple.
\item\label{w3s:W3} (Representation.) $\omega_i$ is a representation of $\mathrm U(V_i)(\A)\times\mathrm U(V_{i+1})(\A)$, and archimedean and finite adelic points commute.
\item\label{w3s:W4} (Bi-action.) For theta data on $(X_j,X_{j+1})$ and $(X_{j+1},X_{j+2})$ and a bikernel $C$ that is a theta function in each variable, the two actions $\omega'_j$, $\omega'_{j+1}$ on the two variables of $C$ commute and preserve this property (formal \texttt{BiAct}, Lean theorem \texttt{seeSawWeilBiAction\_commutes}).
\item\label{w3s:W5} (Flat maps.) The flat maps send theta functions to theta functions and their defining integrals converge absolutely (formal \texttt{FlatLaw}, Lean theorem \texttt{flatLaw}).
\item\label{w3s:W6} (Seesaw exchange.) For $\phi$ continuous and left $\mathrm U(X_j)(F)$-invariant and theta functions $\Xi_1,\Xi_2$ on the reduced spaces, $\mathrm{hand}_{\mathrm{FB}}\bigl(\mathrm{hand}_{\mathrm{BF}}(\phi,\Xi_1),\Xi_2\bigr)=\mathrm{bl}_j\bigl(\phi,\mathrm{contr}(\Xi_1,\Xi_2)\bigr)$, where $\mathrm{contr}$ is the contraction over the link variable, so that $\mathrm{ctr}(\Psi_1,\Psi_2)=\mathrm{contr}(\mathrm{flat}_{\mathrm{BF}}\Psi_1,\mathrm{flat}_{\mathrm{FB}}\Psi_2)$
(formal \texttt{ContrLaw}, Lean theorem \texttt{contrLaw}). Its proof is Fubini's theorem for a bounded continuous integrand on the product of a compact fundamental domain of $[\mathrm U(X_j)]$ and a fundamental domain of covolume one of the link lattice, followed by the factorization of the Weil representation of an orthogonal sum $V^0=A\oplus L$ as the tensor product of the Weil representations of $A$ and $L$, with splitting characters multiplying (\texttt{weil\_datum\_product\_compatibility}).
\item\label{w3s:W7} (Linearity.) $\mathcal B^\sigma_{j+2}(c\phi',c'f)=cc'\,\mathcal B^\sigma_{j+2}(\phi',f)$ and $\mathbf H(c\Phi_1,c'\Phi_2;\dots)=cc'\,\mathbf H(\Phi_1,\Phi_2;\dots)$.
\end{enumerate}
\end{hypothesis}

\begin{hypothesis}[local laws of the tuple]\label{w3s:hyp:loc}
There are data common to all tuples, namely finite disjoint sets $g\mathrm P$ and $t\mathrm P$ of finite places, self-dual lattices $N_i\subset V_i$, base points $b_i\in\mathbf G_i$, compact open subgroups $J_i$ of $\mathrm U(V_i)(F_{g\mathrm P})$, theta functions $\Psi_i(P;a,a')$ on $W_i(\A)$ vacuum at every finite place, combinations $e^G_i$, $d^G_i$ supported at $g\mathrm P$, combinations $d^T_i$ and an operator $\mathrm{op}_j=\sum_k\mathrm{biComb}(u_k,u'_k)$ supported at $t\mathrm P$, a number $c_G(j)$, a set $\mathrm{RamAd}_j$ of continuous automorphic functions on $\mathrm U(X_j)(\A)$ with a map $\mathrm{Sad}_j$ of it into itself. Put $K_i=\{\kappa k:\kappa\text{ stabilizes }N_i\text{ away from }g\mathrm P,\ k\in J_i\}$ and
\[
\kappa_j=\frac{\mu_{j+1}(K_{j+1})}{\rho_J(R_J\cap K_{j+1})}\cdot\frac{\mu_{j+2}(K_{j+2})}{\rho_B(R_B\cap K_{j+2})} .
\]
For each tuple $s$ there are a finite set $Q=Q_s$ of finite places disjoint from $g\mathrm P\cup t\mathrm P$, numbers $\mathrm{sc}_i$, $c_Q(j)$, $c_T(j)$, combinations $e^Q_i$, $d^Q_i(\alpha,\gamma)$ supported at $Q$, and finitely supported functions $\varepsilon_{i,\alpha}$ on $\mathrm U(V_i)(F_Q)$, such that the following hold.
\begin{enumerate}[label=(L\arabic*),leftmargin=*]
\item\label{w3s:L1} (Vacuum form.) $\omega_i(ab_i,a'b_{i+1})\,\mathrm{pt}_i(P)=\mathrm{sc}_i\,\mathrm{comb}(e^G_i)\,\mathrm{comb}(e^Q_i)\,\Psi_i(P;a,a')$ for archimedean $a,a'$. Moreover $e_{i,\alpha}(g)=\varepsilon_{i,\alpha}(gb_i^{-1})$, the combination $\mathrm{dd}_i(\alpha,\gamma)$ acts on theta functions as $\mathrm{comb}(d^Q_i(\alpha,\gamma))\,\mathrm{comb}(d^G_i)\,\mathrm{comb}(d^T_i)$, for $\psi'\in\mathrm{Ram}_j$ one has $\mathrm{rd}_j(\alpha,\psi',y)\in\mathrm{RamAd}_j\cap\mathrm{FX}_{j,\alpha}$ and $\mathrm{rd}_j(\alpha,S_j\psi',y)=\mathrm{Sad}_j(\mathrm{rd}_j(\alpha,\psi',y))$, and for $G\in\mathcal C_i$ and archimedean $a$ the function $R_{ab_i}G(v)$ is automorphic, lies in $\mathrm{FY}_i$ and is invariant under the stabilizer of $N_i$ at the places outside $g\mathrm P\cup Q$. Here $\mathrm{FX}_{j,\alpha}$ and $\mathrm{FY}_i$ are sets of functions given with the tuple (formal \texttt{VTup.FX}, \texttt{VTup.FY}).
\item\label{w3s:L2} (Cells.) Let $\Psi_1,\Psi_2$ be theta functions vacuum at every finite place, put $\Phi^0_i=\mathrm{comb}(e^G_{j+i-1})\mathrm{comb}(e^Q_{j+i-1})\Psi_i$, let $\phi_\alpha\in\mathrm{RamAd}_j$, let $f=R_{a'b_{j+2}}G(v)$ and
\[
I^0(q_1,q_2)=\sum_\alpha\sum_g\varepsilon_{j,\alpha}(g)\,\mathcal B^\sigma_{j+2}\bigl(\mathbf H(\Phi^0_1,\Phi^0_2;\phi_\alpha;g,q_1,q_2),R_{q_2}f\bigr).
\]
Write $\mathbf G_{i,Q}$ for the points of $\mathrm U(V_i)$ at $Q$ and $K_i=K_{i,Q}K'_i$ with $K_{i,Q}$ the stabilizer of $N_i$ at $Q$ and $K'_i$ the part of $K_i$ away from $Q$. Then $I^0$ is left invariant under $R_J$ in $q_1$ and under $R_B$ in $q_2$, right invariant under $K'_{j+1}$ in $q_1$ and under $K'_{j+2}$ in $q_2$, and $I^0(q_1,q_2)\ne0$ implies $q_1\in R_J\mathbf G_{j+1,Q}K'_{j+1}$, and also $q_2\in R_B\mathbf G_{j+2,Q}K'_{j+2}$ when $q_1\in\mathbf G_{j+1,Q}$.
\item\label{w3s:L3} (Twist places.) Let $\rho_{J,Q},\mu_{j+1,Q},\rho_{B,Q},\mu_{j+2,Q}$ be Haar measures on $R_J\cap\mathbf G_{j+1,Q}$, $\mathbf G_{j+1,Q}$, $R_B\cap\mathbf G_{j+2,Q}$, $\mathbf G_{j+2,Q}$ giving $K_{j+1,Q}$, $K_{j+2,Q}$ and their intersections with $R_J$, $R_B$ the volume one. For theta functions $\Psi'_1,\Psi'_2$ vacuum at $Q$, continuous automorphic $\phi_\alpha\in\mathrm{FX}_{j,\alpha}$ and automorphic $f\in\mathrm{FY}_{j+2}$, the function
\[
(q_1,q_2)\mapsto\sum_\alpha\sum_g\varepsilon_{j,\alpha}(g)\,\mathcal B^\sigma_{j+2}\bigl(\mathbf H(\mathrm{comb}(e^Q_j)\Psi'_1,\mathrm{comb}(e^Q_{j+1})\Psi'_2;\phi_\alpha;g,q_1,q_2),R_{q_2}f\bigr)
\]
on $\mathbf G_{j+1,Q}\times\mathbf G_{j+2,Q}$ has the absolutely convergent iterated quotient integral
\[
c_Q(j)\sum_\beta\mathcal B^\sigma_{j+2}\Bigl(\sum_{\alpha,\gamma}\mathrm{bl}_j\bigl(\phi_\alpha,\mathrm{biComb}(d^Q_j(\alpha,\gamma),d^Q_{j+1}(\gamma,\beta))\,\mathrm{ctr}(\Psi'_1,\Psi'_2)\bigr),\ \mathrm{tr}(\varepsilon_{j+2,\beta})f\Bigr).
\]
\item\label{w3s:L4} (Level places.) For theta functions $\Psi_1,\Psi_2$ vacuum at $g\mathrm P$,
$\mathrm{ctr}(\mathrm{comb}(e^G_j)\Psi_1,\mathrm{comb}(e^G_{j+1})\Psi_2)=c_G(j)\,\mathrm{biComb}(d^G_j,d^G_{j+1})\,\mathrm{ctr}(\Psi_1,\Psi_2)$.
\item\label{w3s:L5} (Places of $t\mathrm P$.) For theta functions $\Psi_1,\Psi_2$ vacuum at $t\mathrm P$, combinations $d,d'$ supported away from $t\mathrm P$ and $\phi\in\mathrm{RamAd}_j$,
\begin{multline*}
\mathrm{bl}_j\bigl(\phi,\mathrm{biComb}(d,d')\,\mathrm{ctr}(\Psi_1,\Psi_2)\bigr)\\
=c_T(j)\,\mathrm{bl}_j\bigl(\mathrm{Sad}_j\phi,\ \mathrm{biComb}(d,d')\,\mathrm{biComb}(d^T_j,d^T_{j+1})\,\mathrm{op}_j\,\mathrm{ctr}(\Psi_1,\Psi_2)\bigr),
\end{multline*}
and $\mathrm{op}_j\,\mathrm{ctr}(\Psi_1,\Psi_2)$ is a real bivacuum kernel when $\Psi_1,\Psi_2$ are vacuum at every finite place.
\item\label{w3s:L6} (Normalization.) $c_V(j)=\kappa_j\,\mathrm{sc}_j\,\mathrm{sc}_{j+1}\,c_Q(j)\,c_G(j)\,c_T(j)$.
\end{enumerate}
\end{hypothesis}

\begin{theorem}[the double-step identity (T2)]\label{w3s:thm:T2}
Assume Hypotheses~\ref{w3s:hyp:gen} and~\ref{w3s:hyp:loc}. Let $j+2(\sigma+1)+2=m$, let $P_1,P_2$ be polynomial Fock vectors and $a\in\mathrm U(V_{j+1})_\infty$, $a'\in\mathrm U(V_{j+2})_\infty$. Put $\Psi_1=\Psi_j(P_1;1,a)$, $\Psi_2=\Psi_{j+1}(P_2;a,a')$ and $C=\mathrm{op}_j\,\mathrm{ctr}(\Psi_1,\Psi_2)$. Then $C$ is a real bivacuum kernel that depends only on $(j,\sigma,P_1,P_2,a,a')$ and not on the tuple. For every tuple $s$, every $\psi'\in\mathrm{Ram}_j$, $y$, $G\in\mathcal C_{j+2}$ and $v$, the integrand \eqref{w3s:integrand} has an absolutely convergent iterated quotient integral over $R_J\backslash\mathbf G_{j+1}$ and then $R_B\backslash\mathbf G_{j+2}$, equal to
\begin{equation}\label{w3s:value}
c_V(j)\sum_{\beta\in\mathrm{Ch}_{j+2}}\mathcal B^\sigma_{j+2}\Bigl(\mathrm{two}_j\bigl(\mathrm{rd}_j(\cdot,S_j\psi',y),C\bigr)_\beta,\ \mathrm{tr}(e_{j+2,\beta})\,R_{a'}G(v)\Bigr),
\end{equation}
with $c_V(j)=c_V(j)(s)=\kappa_j\,\mathrm{sc}_j(s)\,\mathrm{sc}_{j+1}(s)\,c_Q(j)(s)\,c_G(j)\,c_T(j)(s)$ and $\kappa_j>0$ independent of $s$.
\end{theorem}

\begin{proof}
\emph{Step 1, the kernel.} The functions $\Psi_1,\Psi_2$ belong to the common data, and by \ref{w3s:W2} so does $\mathrm{ctr}(\Psi_1,\Psi_2)$. Hence $C$ depends only on $(j,\sigma,P_1,P_2,a,a')$. It is a real bivacuum kernel by \ref{w3s:L5}. The groups $K_{j+1}$ and $K_{j+2}$ are compact and open, and their intersections with $R_J$ and $R_B$ are compact open subgroups of $R_J$ and $R_B$, so the four volumes in $\kappa_j$ are positive and finite and $\kappa_j>0$. It does not involve the tuple.

\emph{Step 2, the twist places.} Fix $s$, $\psi'$, $y$, $G$, $v$, and put $\phi_\alpha=\mathrm{rd}_j(\alpha,\psi',y)$, which lies in $\mathrm{RamAd}_j\cap\mathrm{FX}_{j,\alpha}$ by \ref{w3s:L1}, and $f=R_{a'b_{j+2}}G(v)$, which is automorphic and lies in $\mathrm{FY}_{j+2}$ by \ref{w3s:L1}. The functions $\Psi_i$ are vacuum at every finite place, so by \ref{w3s:W1} the functions $\Psi'_i=\mathrm{comb}(e^G_{j+i-1})\Psi_i$ are vacuum at $Q$, because $Q$ and $g\mathrm P$ are disjoint, and $\mathrm{comb}(e^Q)\Psi'_i=\Phi^0_i$ in the notation of \ref{w3s:L2}, because combinations at disjoint places commute. By \ref{w3s:L3} the restriction of the function $I^0$ of \ref{w3s:L2} to $\mathbf G_{j+1,Q}\times\mathbf G_{j+2,Q}$ has an absolutely convergent iterated quotient integral $z_Q$ for the measures normalized at $Q$, with $z_Q$ the value displayed in \ref{w3s:L3} for $\mathrm{ctr}(\Psi'_1,\Psi'_2)$. Let $f_Q$ be the corresponding inner function on $\mathbf G_{j+1,Q}$.

\emph{Step 3, collapse of the inner integral.} Apply Lemma~\ref{w3s:lem:collapse} to $\mathbf G=\mathbf G_{j+2}$, $G_1=\mathbf G_{j+2,Q}$, $G_2$ the points away from $Q$, $R=R_B$, $K_1=K_{j+2,Q}$, $K_2=K'_{j+2}$. The groups involved are the finite adelic points of unitary groups, which are unimodular, and the Bessel and Fourier--Jacobi subgroups, which are semidirect products of a unipotent group by a unitary group and are unimodular by the argument of Lemma~\ref{w2:unimod}. The finite adelic unitary group is the restricted product of its local groups, which gives the product decomposition and the splitting of $R_B$, since $R_B$ is the group of finite adelic points of an $F$-subgroup. For $q_1\in\mathbf G_{j+1,Q}$ the function $I^0(q_1,\cdot)$ satisfies (i) to (iii) by \ref{w3s:L2}, so its quotient integral over $R_B\backslash\mathbf G_{j+2}$ converges absolutely with value $\nu_{j+2}f_Q(q_1)$, where $\nu_{j+2}=\mu_{j+2}(K_{j+2})/\rho_B(R_B\cap K_{j+2})$. For an arbitrary $q_1\in\mathbf G_{j+1}$, either $I^0(q_1,\cdot)=0$, and the inner integral is $0$, or by \ref{w3s:L2} $q_1=r\,x\,k$ with $r\in R_J$, $x\in\mathbf G_{j+1,Q}$, $k\in K'_{j+1}$, and $I^0(q_1,\cdot)=I^0(x,\cdot)$ by the invariances of \ref{w3s:L2}. Let $f_1(q_1)$ be the value of the inner integral. By Lemma~\ref{w3s:lem:uniq} it is well defined, and the invariances and support of $I^0$ in $q_1$ pass to $f_1$. On $\mathbf G_{j+1,Q}$ one has $f_1=\nu_{j+2}f_Q$.

\emph{Step 4, collapse of the outer integral.} Lemma~\ref{w3s:lem:collapse}, applied to $\mathbf G_{j+1}$, $R_J$, $K_{j+1}=K_{j+1,Q}K'_{j+1}$ and the function $f_1$, shows that $f_1$ has the absolutely convergent quotient integral $\nu_{j+1}\nu_{j+2}z_Q=\kappa_jz_Q$ over $R_J\backslash\mathbf G_{j+1}$, with $\nu_{j+1}=\mu_{j+1}(K_{j+1})/\rho_J(R_J\cap K_{j+1})$. So $I^0$ has the absolutely convergent iterated quotient integral $\kappa_jz_Q$.

\emph{Step 5, base points.} For $p_1\in\mathbf G_{j+1}$, $p_2\in\mathbf G_{j+2}$ and $g'$ in the support of $\varepsilon_{j,\alpha}$, \ref{w3s:W3} gives $(g'b_j,a\,p_1b_{j+1})=(g',p_1)(b_j,ab_{j+1})$ and $(a\,p_1b_{j+1},a'\,p_2b_{j+2})=(p_1,p_2)(ab_{j+1},a'b_{j+2})$. By \ref{w3s:L1},
\begin{align*}
\omega_j(g'b_j,a\,p_1b_{j+1})\,\mathrm{pt}_j(P_1)&=\mathrm{sc}_j\,\omega_j(g',p_1)\Phi^0_1,\\
\omega_{j+1}(a\,p_1b_{j+1},a'\,p_2b_{j+2})\,\mathrm{pt}_{j+1}(P_2)&=\mathrm{sc}_{j+1}\,\omega_{j+1}(p_1,p_2)\Phi^0_2,
\end{align*}
and $R_{a'p_2b_{j+2}}G(v)=R_{p_2}f$. Since $e_{j,\alpha}(g)=\varepsilon_{j,\alpha}(gb_j^{-1})$, \ref{w3s:W7} gives $I_s(p_1b_{j+1},p_2b_{j+2})=\mathrm{sc}_j\mathrm{sc}_{j+1}\,I^0(p_1,p_2)$. The right translations by $b_{j+1}^{-1}$ and $b_{j+2}^{-1}$ preserve absolutely convergent quotient integrals with their values (Lemma~\ref{w3s:lem:uniq}). Hence $I_s$ has the absolutely convergent iterated quotient integral $\mathrm{sc}_j\mathrm{sc}_{j+1}\kappa_jz_Q$.

\emph{Step 6, the value.} By \ref{w3s:L4} and \ref{w3s:W1}, \ref{w3s:W4},
\begin{multline*}
\mathrm{biComb}(d^Q_j(\alpha,\gamma),d^Q_{j+1}(\gamma,\beta))\,\mathrm{ctr}(\Psi'_1,\Psi'_2)\\
=c_G(j)\,\mathrm{biComb}\bigl(d^Q_j(\alpha,\gamma)d^G_j,\ d^Q_{j+1}(\gamma,\beta)d^G_{j+1}\bigr)\,\mathrm{ctr}(\Psi_1,\Psi_2),
\end{multline*}
where products of combinations denote composition of the operators. These combinations are supported away from $t\mathrm P$. By \ref{w3s:L5}, then \ref{w3s:L1} for $\mathrm{dd}$ and for $\mathrm{Sad}_j$,
\[
\sum_{\alpha,\gamma}\mathrm{bl}_j\bigl(\phi_\alpha,\mathrm{biComb}(d^Q_j(\alpha,\gamma),d^Q_{j+1}(\gamma,\beta))\,\mathrm{ctr}(\Psi'_1,\Psi'_2)\bigr)=c_G(j)\,c_T(j)\,\mathrm{two}_j\bigl(\mathrm{rd}_j(\cdot,S_j\psi',y),C\bigr)_\beta .
\]
Finally $\mathrm{tr}(e_{j+2,\beta})R_{a'}G(v)=\mathrm{tr}(\varepsilon_{j+2,\beta})R_{a'b_{j+2}}G(v)=\mathrm{tr}(\varepsilon_{j+2,\beta})f$, because $e_{j+2,\beta}(g)=\varepsilon_{j+2,\beta}(gb_{j+2}^{-1})$ and $a'$ commutes with $b_{j+2}$. By \ref{w3s:W7} the iterated integral of Step 5 equals $\kappa_j\mathrm{sc}_j\mathrm{sc}_{j+1}c_Q(j)c_G(j)c_T(j)$ times the sum in \eqref{w3s:value}, and \ref{w3s:L6} identifies the constant with $c_V(j)$.
\end{proof}

\begin{proposition}[nonvanishing of $c_V(j)$]\label{w3s:prop:cV}
Under \ref{w3s:L6}, $c_V(j)(s)\ne0$ if and only if $\mathrm{sc}_j(s)$, $\mathrm{sc}_{j+1}(s)$, $c_Q(j)(s)$, $c_G(j)$ and $c_T(j)(s)$ are all nonzero.
\end{proposition}

\begin{proof}
This is $\kappa_j>0$ (Step 1 of the proof of Theorem~\ref{w3s:thm:T2}) together with \ref{w3s:L6}.
\end{proof}

\begin{remark}[normalization]\label{w3s:rem:corr}
The constant $c_V(j)$ depends on the tuple through $\mathrm{sc}_j$, $\mathrm{sc}_{j+1}$, $c_Q(j)$ and $c_T(j)$. The object that does not depend on the tuple is the kernel $C$, together with the volume factor $\kappa_j$. In Proposition~\ref{w2:see:step} the constants $c_V(j)c_K(j)c_K(j+1)$ stand inside the sum over the tuples $s_i$, and the proof needs only that one kernel $C$ serves all tuples, so that the coefficients $a(y')$ of the kernel law (T3) can be pulled out of the sum. Nonvanishing of $c_V(j)(s)$ is not used in Proposition~\ref{w2:see:step} or Theorem~\ref{w2:see:main}. It enters only through $c_{\mathrm{top}}(t)\ne0$ in (D6), and Proposition~\ref{w3s:prop:cV} reduces it to the nonvanishing of the five scalars of the construction.
\end{remark}

\begin{remark}[what remains]\label{w3s:rem:open}
Theorem~\ref{w3s:thm:T2} reduces (T2) to Hypotheses~\ref{w3s:hyp:gen} and~\ref{w3s:hyp:loc}. Items \ref{w3s:W1} to \ref{w3s:W7} are general properties of the Weil representation, the flat maps and the Bessel period, and \ref{w3s:W6} is Fubini's theorem combined with the tensor factorization of the Weil representation of an orthogonal sum. The proof above uses \ref{w3s:W1} to \ref{w3s:W4} and \ref{w3s:W7} directly. Items \ref{w3s:W5} and \ref{w3s:W6} are used in the formal derivation of the cell properties \ref{w3s:L2} and of the twist-place law \ref{w3s:L3}. Items \ref{w3s:L1} to \ref{w3s:L6} are local computations of the construction at the twist places $Q$, at the level places $g\mathrm P$, at the places $t\mathrm P$ and at the places where every datum is unramified. Of these, \ref{w3s:L3} is the local double-step identity at the twist places. The Siegel--Weil formula and the doubling method do not enter (T2). They enter Hypothesis~\ref{w2:see:H-links} only through the surjectivity (D4).
\end{remark}

\subsection{The share of irreducible rows}\label{w2:sec:irr}

\subsubsection{Labels and weights}

Let $g$ be a prime and $\Gamma_N=\mathrm{GL}_N(\F_g)$. A \emph{label} is a monic polynomial $c\in\F_g[X]$ of degree $N$ with $c(0)\ne0$. Every irreducible representation $\pi$ of $\Gamma_N$ has a label, the characteristic polynomial of the semisimple class $s_\pi$ of its Deligne--Lusztig series. For a label $c$ let $E_c\in\C[\Gamma_N]$ be the sum of the central idempotents of the irreducible representations with label $c$. The $E_c$ are orthogonal central idempotents with $\sum_cE_c=1$. For $e\in\C[\Gamma_N]$ with $e(1)\ne0$ put
\[
\mathrm{wt}(e,c)=\frac{\operatorname{Re}(eE_c)(1)}{\operatorname{Re}e(1)},
\]
where $a(1)$ is the coefficient of the identity. For a squarefree label put $\sq(c)=\prod_{P\mid c}(g^{\deg P}-1)^{-1}$, the product over the monic irreducible factors.

\begin{lemma}\label{w2:wtprob}
If $e=e^*=e^2$ and $e(1)>0$, then $\mathrm{wt}(e,c)\ge0$ for every label and $\sum_c\mathrm{wt}(e,c)=1$.
\end{lemma}

\begin{proof}
In the left regular representation $\lambda$ on $\ell^2(\Gamma_N)$ one has $a(1)=\langle\lambda(a)\delta_1,\delta_1\rangle$. Central idempotents of $\C[\Gamma_N]$ are self-adjoint, so $eE_c=E_ceE_c$ and $(eE_c)(1)=\lVert\lambda(e)\lambda(E_c)\delta_1\rVert^2\ge0$. Summing over $c$ gives $e(1)$.
\end{proof}

Let $U\subset\Gamma_N$ be the upper unitriangular group, $\psi_0$ a nontrivial additive character of $\F_g$, and $e_U=|U|^{-1}\sum_{u\in U}\overline{\psi_0(\sum_iu_{i,i+1})}\,u$ the Gelfand--Graev idempotent. It is a self-adjoint idempotent with $e_U(1)=|U|^{-1}$, because $u\mapsto\psi_0(\sum_iu_{i,i+1})$ is a character of $U$.

\begin{proposition}[Gelfand--Graev weights]\label{w2:GG}
For every squarefree label $c$, $\mathrm{wt}(e_U,c)=\sq(c)$.
\end{proposition}

\begin{proof}
By Frobenius reciprocity $(e_UE_\pi)(1)=|\Gamma_N|^{-1}\dim\pi\cdot m_\pi$, where $m_\pi$ is the multiplicity of $\pi$ in the Gelfand--Graev representation $\mathrm{Ind}_U^{\Gamma_N}\psi$, so $\mathrm{wt}(e_U,c)=[\Gamma_N:U]^{-1}\sum_{\pi\text{ of label }c}m_\pi\dim\pi$. By Gelfand--Graev multiplicity one and the theory of Deligne and Lusztig \cite[\S10]{w2dl}, \cite[Ch.~8]{w2carter}, each series contains exactly one constituent of the Gelfand--Graev representation, with multiplicity one, the regular character $\chi_{s,\mathrm{reg}}$, and $\chi_{s,\mathrm{reg}}(1)=[\Gamma_N:C(s)]_{g'}$, the prime-to-$g$ part of the index of the centralizer. For $\Gamma_N=\mathrm{GL}_N(\F_g)$ this is also Green's theorem \cite{w2green}. If $c$ is squarefree with irreducible factors of degrees $d_1,\dots,d_j$, the centralizer of $s$ is the torus $\prod_i\F_{g^{d_i}}^\times$, of order $\prod_i(g^{d_i}-1)$, prime to $g$. Since $[\Gamma_N:U]=|\Gamma_N|_{g'}$, $\mathrm{wt}(e_U,c)=|\Gamma_N|_{g'}/(\prod_i(g^{d_i}-1)\,|\Gamma_N|_{g'})=\sq(c)$.
\end{proof}

\subsubsection{Good tuples}

Fix a prime $g$, an integer $N\ge100$, a finite set $\Pi$ of $f$ places, a finite set $J$ of $m$ indices and elements $\chi_{w,j}\in\F_g^\times$ for $w\in\Pi$, $j\in J$. Let $\bar\F_g$ be an algebraic closure. For a prime $\ell\ne g$ and $x\in\bar\F_g^\times$ let $x_{\ell'}$ be the prime-to-$\ell$ component of $x$ in the torsion group $\bar\F_g^\times$, and let
\[
BB(g,N)=\{g\}\cup\bigl\{\ell\le g^N\text{ prime}:\ g^d<\ell^{4v_\ell(g^d-1)}\text{ for some }1\le d\le N\bigr\}.
\]
A tuple $(c_w)_{w\in\Pi}$ of labels is \emph{good} if the following hold.
\begin{enumerate}[label=(G\arabic*),leftmargin=*]
\item\label{w2:G1} Every $c_w$ is squarefree.
\item\label{w2:G2} For every prime $\ell\notin BB(g,N)$ and every $w$, the map $x\mapsto x_{\ell'}$ is injective on the roots of $c_w$.
\item\label{w2:G3} For every prime $\ell\ne g$, every $w$, $j$, every sign and every root $x$ of $c_w$, $x_{\ell'}\ne(\chi_{w,j}^{\pm1})_{\ell'}$.
\item\label{w2:G4} For every $1\le a\le N/2$ there is $w$ such that $c_w$ has no divisor of degree $a$.
\item\label{w2:G5} There is $w$ such that $c_w$ has an irreducible factor of prime degree $p$ with $N<2p$ and $p+3\le N$.
\end{enumerate}


\subsubsection{Counting lemmas}

\begin{lemma}\label{w2:counts}
Let $N\ge2$ and $g\ge2$. (i) The number $I_d$ of monic irreducible polynomials of degree $d$ in $\F_g[X]$ other than $X$ satisfies $dI_d\le g^d-1$. (ii) At least $g^N-2g^{N-1}$ labels of degree $N$ are squarefree. (iii) For every squarefree label $c$ of degree $N$, $\sq(c)\ge g^{-N}$.
\end{lemma}

\begin{proof}
(i) The roots of the counted polynomials are distinct elements of $\F_{g^d}^\times$, so $dI_d\le g^d-1$. (ii) The number of squarefree monic polynomials of degree $N\ge2$ is $g^N-g^{N-1}$ \cite[Prop.~2.3]{w2rosen}, and those with $c(0)=0$ are $X$ times a squarefree monic polynomial of degree $N-1$, at most $g^{N-1}$ of them. (iii) $\prod_P(g^{\deg P}-1)^{-1}\ge\prod_Pg^{-\deg P}=g^{-N}$.
\end{proof}

\begin{lemma}[types versus cycle types]\label{w2:types}
Let $\mathcal A$ be a property of partitions of $N$. Then
\[
\sum_{\substack{c\text{ squarefree label}\\ \mathcal A(\text{degrees of the factors of }c)}}\sq(c)\le\frac{\#\{\tau\in S_N:\ \mathcal A(\text{cycle type of }\tau)\}}{N!}.
\]
In particular $\sum_{c\text{ squarefree}}\sq(c)\le1$.
\end{lemma}

\begin{proof}
A squarefree label with $n_d$ irreducible factors of degree $d$ is a choice of $n_d$ distinct irreducibles of each degree $d$, none equal to $X$. The total weight of the labels of this type is $\prod_d\binom{I_d}{n_d}(g^d-1)^{-n_d}\le\prod_d\frac{1}{n_d!}\bigl(\frac{I_d}{g^d-1}\bigr)^{n_d}\le\prod_d\frac{1}{d^{n_d}n_d!}$ by Lemma~\ref{w2:counts}(i). By Cauchy's formula the last product is the proportion of permutations of cycle type $(n_d)$. Sum over the types with property $\mathcal A$.
\end{proof}

\begin{lemma}[multiples]\label{w2:mult}
Let $h$ be a squarefree monic polynomial of degree at most $N$ with $h(0)\ne0$. Then $\sum\sq(c)\le\sq(h)$, the sum over squarefree labels $c$ of degree $N$ divisible by $h$.
\end{lemma}

\begin{proof}
Such $c$ are $c=hk$ with $k$ squarefree, prime to $h$, of degree $N-\deg h$ and $k(0)\ne0$, and $\sq(hk)=\sq(h)\sq(k)$. By Lemma~\ref{w2:types} in degree $N-\deg h$ (or trivially if the degree is $0$), $\sum_k\sq(k)\le1$.
\end{proof}

\begin{lemma}[invariant sets of permutations]\label{w2:78}
For $1\le a<n$, the proportion of $\tau\in S_n$ that leave some $a$-element subset invariant is at most $7/8$.
\end{lemma}

\begin{proof}
Let $X(\tau)$ be the number of $\tau$-invariant $a$-subsets. Counting pairs $(\tau,A)$ with $A$ invariant gives $\sum_\tau X(\tau)=\binom na\,a!\,(n-a)!=n!$. Among these pairs consider those for which $\tau|_A$ and $\tau|_{A^c}$ both have a fixed point. With $D_j$ the number of derangements of $j$ points, $2D_j\le j!$ for $j\ge1$, so the proportion of such pairs is $(1-D_a/a!)(1-D_{n-a}/(n-a)!)\ge1/4$. If $(\tau,A)$ is such a pair, with fixed points $p\in A$ and $q\notin A$, then $A\setminus\{p\}\cup\{q\}$ is another invariant $a$-subset, so $X(\tau)\ge2$. Let $B=\{\tau:X(\tau)\ge2\}$ and $n_1=\#\{\tau:X(\tau)=1\}$. Then $\sum_{\tau\in B}X(\tau)\ge n!/4$, so $n_1=n!-\sum_{\tau\in B}X(\tau)\le\tfrac34n!$, and $\#B\le\tfrac12\sum_{\tau\in B}X(\tau)=\tfrac12(n!-n_1)$. Hence $\#\{X\ge1\}=n_1+\#B\le\tfrac12n!+\tfrac12n_1\le\tfrac78n!$.
\end{proof}

\begin{lemma}[long prime cycles]\label{w2:primecycles}
For $N\ge100$ the proportion of $\tau\in S_N$ that have a cycle of prime length $p$ with $N/2<p\le N-3$ is $\sum_{N/2<p\le N-3}1/p\ge1/(5\log N)$.
\end{lemma}

\begin{proof}
A permutation has at most one cycle of length greater than $N/2$, and the number of permutations with a cycle of a given length $L>N/2$ is $N!/L$, so the events are disjoint and the proportion is the sum. For $100\le N\le2\cdot10^6$ the inequality was checked by direct computation (the smallest value of $5\log N\sum1/p$ is $2.94$, at $N=226$). For $N>2\cdot10^6$, the bounds $x/\log x<\pi(x)<1.25506\,x/\log x$ of Rosser and Schoenfeld \cite[Cor.~1]{w2rs}, valid for $x\ge17$, give $\sum_{N/2<p\le N-3}1/p\ge(\pi(N-3)-\pi(N/2))/(N-3)$ with $\pi(N-3)-\pi(N/2)\ge\frac{N-3}{\log N}-0.6276\frac{N}{\log(N/2)}\ge\frac{0.339N-3}{\log N}$, because $\log(N/2)\ge0.95\log N$ for $N>2\cdot10^6$, and $(0.339N-3)/(N-3)\ge1/5$.
\end{proof}

\begin{lemma}[prime power parts of $g-1$]\label{w2:share}
Let $g$ be a prime with $16\mid g-1$ and let $q_1\ne d$ be odd primes dividing $g-1$ with $\max(q_1,d)\ge16$. Then $s(g):=\sum_{\ell\mid g-1}\ell^{v_\ell(g-1)}/(g-1)\le\frac{49}{256\min(q_1,d)}$.
\end{lemma}

\begin{proof}
Write $g-1=\prod_ia_i$ with $a_i=\ell_i^{v_{\ell_i}(g-1)}$, so $s(g)=\sum_i1/\prod_{j\ne i}a_j$. Here $a_2\ge16$, $a_{q_1}\ge q_1$, $a_d\ge d$, and the remaining factors $b_1,\dots,b_t\ge3$ have product $B$ with $\sum_ib_i\le B$ when $t\ge1$. The terms are at most $1/(q_1d)$, $1/(16d)$, $1/(16q_1)$, and $\sum_ib_i/(16q_1dB)\le1/(16q_1d)$. With $m_0=\min(q_1,d)$ and $\max(q_1,d)\ge16$ the total is at most $(1/16+1/8+1/256)/m_0=49/(256m_0)$.
\end{proof}

\subsubsection{The bad weight at one prime}

\begin{theorem}\label{w2:badweight}
Let $g\ge2^{12}$ be a prime, $N\ge100$, and let $q_1,d$ be as in Lemma~\ref{w2:share} with $m_0=\min(q_1,d)$. Let $P$ be a nonnegative function on labels of degree $N$ with total mass $1$ that equals $\sq$ on squarefree labels. Then the $P^{\otimes\Pi}$-mass of the tuples that are not good is at most
\[
\begin{aligned}
W(g)={}&\frac{2f}{g}+\frac{f\log_2g}{\sqrt g}\Bigl(1.31+\frac{2.02\,(1+\log N)}{g^{1/4}}\Bigr)\\
&+2fm\Bigl(\frac1{4m_0}+\frac{1+\log N}{g-1}\Bigr)+\frac N2\Bigl(\frac78\Bigr)^f+e^{-f/(5\log N)}.
\end{aligned}
\]
\end{theorem}

\begin{proof}
A tuple that is not good satisfies one of the following: (I) some $c_w$ is not squarefree; (II) some $c_w$ is squarefree and fails \ref{w2:G2}; (III) for some $(w,j,\pm)$, $c_w$ is squarefree and fails \ref{w2:G3} for that index; (IV) all $c_w$ are squarefree and \ref{w2:G4} fails; (V) all $c_w$ are squarefree and \ref{w2:G5} fails. The mass is at most the sum of the masses of the five events. For an event that concerns one place $w$, the mass is the $P$-mass of the corresponding set of labels at $w$, because $P$ has mass $1$ at the other places.

(I) By Lemma~\ref{w2:counts}(ii),(iii), the squarefree labels have $\sq$-mass at least $(g^N-2g^{N-1})g^{-N}=1-2/g$. So the non-squarefree labels have $P$-mass at most $2/g$, and the union over $f$ places at most $2f/g$.

(II) Fix $w$ and a prime $\ell\notin BB(g,N)$. For $1\le\delta\le N$ the Sylow $\ell$-subgroup of $\F_{g^\delta}^\times$ has order $\ell^{v_\ell(g^\delta-1)}\le g^{\delta/4}$: for $\ell\le g^N$ this is the definition of $BB$, and for $\ell>g^N$ the group is trivial. A squarefree label fails \ref{w2:G2} if it has two distinct roots with equal $\ell'$-components. Either both are roots of one irreducible factor, or they are roots of two distinct factors. By Lemma~\ref{w2:mult} the mass is at most the sum of $\sq(Q)$ over the irreducible $Q$ with a self-collision plus the sum of $\sq(QQ')$ over the ordered pairs $Q\ne Q'$ with a cross collision.

Self-collisions. Let $Q$ have degree $\delta\ge2$ and roots $x,x^{g^i}$ with $1\le i<\delta$ and $(x^{g^i})_{\ell'}=x_{\ell'}$, so $x^{g^i-1}$ lies in the Sylow $\ell$-subgroup of $\F_{g^\delta}^\times$, which is then nontrivial, so $\ell\mid g^\delta-1$. The kernel of $y\mapsto y^{g^i-1}$ on $\F_{g^\delta}^\times$ is $\F_{g^{\gcd(i,\delta)}}^\times$, of order less than $g^{\delta/2}$. So for fixed $i$ and $\ell$ there are at most $g^{3\delta/4}$ such $x$. There are at most $\delta\log_2g$ primes dividing $g^\delta-1$ and $\delta-1$ values of $i$, so at most $(\delta-1)\log_2g\,g^{3\delta/4}$ such $Q$ of degree $\delta$, summed over all $\ell\notin BB$. With $\rho=g^{-1/4}\le1/8$ their mass is at most
\[
\log_2g\sum_{\delta\ge2}(\delta-1)\frac{g^{3\delta/4}}{g^\delta-1}\le\bigl(1+2^{-23}\bigr)\log_2g\,\frac{\rho^2}{(1-\rho)^2}\le1.31\,\frac{\log_2g}{\sqrt g}.
\]
Cross collisions. Let $Q\ne Q'$ have degrees $\delta,\delta'$ and roots $x,y$ with $x_{\ell'}=y_{\ell'}=:z$. Write $x=zx_\ell$, $y=zy_\ell$ with $\ell$-power components $x_\ell\in\F_{g^\delta}^\times$, $y_\ell\in\F_{g^{\delta'}}^\times$. Since $x\ne y$, one of $x_\ell,y_\ell$ is not $1$. Let $\mathcal A$ be the set of $(Q,Q')$ that admit such $(x,y,\ell)$ with $y_\ell\ne1$, and $\mathcal B$ the set with $x_\ell\ne1$. The Frobenius $t\mapsto t^g$ maps such a triple $(x,y,\ell)$ to another one with the same property, and moves $x$ through all $\delta$ roots of $Q$. So every $(Q,Q')\in\mathcal A$ yields at least $\delta$ triples with distinct $x$. Given $x$, the prime $\ell$ divides $g^{\delta'}-1$ (at most $\delta'\log_2g$ choices) and $y=x_{\ell'}y_\ell$ is determined by $y_\ell$ (at most $g^{\delta'/4}$ choices). Hence $\#\mathcal A\le g^\delta\delta'\log_2g\,g^{\delta'/4}/\delta$, and symmetrically $\#\mathcal B\le g^{\delta'}\delta\log_2g\,g^{\delta/4}/\delta'$. The mass is at most
\[
\sum_{\delta,\delta'=1}^N\frac{\#\mathcal A+\#\mathcal B}{(g^\delta-1)(g^{\delta'}-1)}\le2\Bigl(1+\frac1{g-1}\Bigr)^2\log_2g\Bigl(\sum_{\delta\le N}\frac1\delta\Bigr)\sum_{k\ge1}kg^{-3k/4}\le2.02\,\frac{(1+\log N)\log_2g}{g^{3/4}}.
\]
Adding and summing over the $f$ places gives the second term of $W(g)$.

(III) Fix $(w,j,\pm)$ and put $\chi=\chi_{w,j}^{\pm1}\in\F_g^\times$. For $\ell\ne g$, $x_{\ell'}=\chi_{\ell'}$ holds exactly when $x/\chi$ has $\ell$-power order. Let $Q$ be irreducible of degree $\delta$ with such a root $x$. Then $x\in\chi\cdot\Syl_\ell(\F_{g^\delta}^\times)$, and every root of $Q$ has the same property because $t\mapsto t^g$ fixes $\chi$. Hence the number of such $Q$ is at most $\delta^{-1}\sum_{\ell\mid g^\delta-1}\ell^{v_\ell(g^\delta-1)}$; the primes $\ell\nmid g^\delta-1$ contribute the single element $x=\chi$, which lies in every coset counted. For $\ell\mid g-1$, the lifting of the exponent gives $v_\ell(g^\delta-1)=v_\ell(g-1)+v_\ell(\delta)$, valid for odd $\ell$ and for $\ell=2$ because $4\mid g-1$, so $\ell^{v_\ell(g^\delta-1)}\le\delta\,\ell^{v_\ell(g-1)}$. For $\ell\nmid g-1$ the prime powers $\ell^{v_\ell(g^\delta-1)}$ are pairwise coprime divisors of $(g^\delta-1)/(g-1)$, so their sum is at most $(g^\delta-1)/(g-1)$. With $s(g)$ of Lemma~\ref{w2:share}, the mass of the labels with a bad factor is at most, by Lemma~\ref{w2:mult},
\[
\sum_{\delta=1}^N\frac{\delta(g-1)s(g)+(g^\delta-1)/(g-1)}{\delta\,(g^\delta-1)}\le s(g)\Bigl(1+\frac1{g-1}\Bigr)+\frac{1+\log N}{g-1}\le\frac1{4m_0}+\frac{1+\log N}{g-1}.
\]
There are $2fm$ indices $(w,j,\pm)$.

(IV) A squarefree label has a divisor of degree $a$ exactly when the degrees of its irreducible factors have a sub-multiset with sum $a$, a property of the type. By Lemmas~\ref{w2:types} and~\ref{w2:78} the $\sq$-mass of the squarefree labels with a divisor of degree $a$ is at most $7/8$. For fixed $a$ the event that this holds at every place has mass at most $(7/8)^f$, and there are at most $N/2$ values of $a$.

(V) By Lemmas~\ref{w2:types} and~\ref{w2:primecycles} the $\sq$-mass of the squarefree labels without an irreducible factor of prime degree in $(N/2,N-3]$ is at most $1-1/(5\log N)$. At all $f$ places the mass is at most $(1-1/(5\log N))^f\le e^{-f/(5\log N)}$.
\end{proof}

\subsubsection{From weights to dimensions}

The construction attaches to each element $\tau$ of an orthonormal eigenbasis of $\mathcal H$ a tuple of labels $\mathrm{lab}_k(\tau)=(\mathrm{lab}_{k,w}(\tau))_{w\in\Pi_k}$ at each of two auxiliary primes $g_0,g_1$, where $\Pi_k$ is the set of $f$ places of $F$ above $g_k$ and the level at $w$ contains $\mathrm{GL}_N(\F_{g_k})$ as a quotient, with $N=r$. For each $(k,w)$ it fixes an element $e_{k,w}\in\C[\mathrm{GL}_N(\F_{g_k})]$, the local factor of the level idempotent, and characters $\chi_{w,j}\in\F_{g_k}^\times$, $j\in J$, $|J|=m$. For a set $S$ of label tuples at $k$ put
\[
w_k(S)=\sum_{t\in S}\prod_{w\in\Pi_k}\mathrm{wt}(e_{k,w},t_w).
\]

\begin{hypothesis}[local weights]\label{w2:H-wt}
Each $e_{k,w}$ is a self-adjoint idempotent with $e_{k,w}(1)>0$, and $(1-1/g_k)\,\mathrm{wt}(e_U,c)\le\mathrm{wt}(e_{k,w},c)$ for every label $c$.
\end{hypothesis}

\begin{hypothesis}[multiplicity formula]\label{w2:H-cnt}
There is $\varepsilon_C\in[0,\tfrac12]$ with $\lvert D_{\mathrm{all}}-D^*\rvert\le\varepsilon_CD^*$ and, for $k\in\{0,1\}$ and every set $S$ of label tuples at $k$,
\[
\Bigl\lvert\#\{\tau:\mathrm{lab}_k(\tau)\in S\}-D^*w_k(S)\Bigr\rvert\le\varepsilon_C\,D^*w_k(S).
\]
\end{hypothesis}

\begin{hypothesis}[rows are the good forms]\label{w2:H-rows}
$D_{\mathrm{rows}}$ equals the number of $\tau$ whose label tuples at $g_0$ and at $g_1$ are both good for the characters $\chi_{w,j}$.
\end{hypothesis}

\begin{lemma}[comparison of weights]\label{w2:transfer}
Let $w_U,w_J$ be probability weights on a finite set at each of $f$ places with $(1-1/g)w_U\le w_J$. For every set $A$ of tuples, $w_J^{\otimes f}(A)\le w_U^{\otimes f}(A)+f/g$.
\end{lemma}

\begin{proof}
Write $w_J=(1-1/g)w_U+\tfrac1g\nu$ with $\nu$ a probability weight. Expand the product. The term $(1-1/g)^fw_U^{\otimes f}(A)$ is at most $w_U^{\otimes f}(A)$, and the other terms have total mass $1-(1-1/g)^f\le f/g$.
\end{proof}

\begin{theorem}\label{w2:Iproof}
Assume Hypotheses~\ref{w2:H-wt}, \ref{w2:H-cnt} and~\ref{w2:H-rows}, that $16\mid g_k-1$ and $q_1,d$ divide $g_k-1$ for both $k$, with $q_1\ne d$ odd primes and $\max(q_1,d)\ge16$, and the sizes
\begin{equation}\label{w2:sizes}
g_k\ge e^{8x^{\varepsilon'}},\quad x^2\le f\le x^{C_s},\quad m\le x^{C_s},\quad100\le N\le x^{C_s},\quad\min(q_1,d)\ge e^{x^{\varepsilon'/6}},
\end{equation}
with $0<\varepsilon'\le1$ and $C_s\ge1$. Then $0\le\delta\le\exp(-x^{\varepsilon'/7})$ for $x\ge x_0(\varepsilon',C_s)$. With $\varepsilon'=\varepsilon_z$ this is Theorem~\ref{w2:I}.
\end{theorem}

\begin{proof}
Let $B_k$ be the set of label tuples at $g_k$ that are not good. By Hypothesis~\ref{w2:H-rows} every $\tau$ outside the rows has $\mathrm{lab}_0(\tau)\in B_0$ or $\mathrm{lab}_1(\tau)\in B_1$, so $0\le D_{\mathrm{all}}-D_{\mathrm{rows}}\le\sum_k\#\{\tau:\mathrm{lab}_k(\tau)\in B_k\}$, which gives $\delta\ge0$. By Hypothesis~\ref{w2:H-cnt},
\[
\delta=\frac{D_{\mathrm{all}}-D_{\mathrm{rows}}}{D_{\mathrm{all}}}\le\frac{(1+\varepsilon_C)D^*(w_0(B_0)+w_1(B_1))}{(1-\varepsilon_C)D^*}\le3\bigl(w_0(B_0)+w_1(B_1)\bigr).
\]
By Lemma~\ref{w2:wtprob} and Hypothesis~\ref{w2:H-wt} the weights $\mathrm{wt}(e_{k,w},\cdot)$ and $\mathrm{wt}(e_U,\cdot)$ are probability weights with $(1-1/g_k)\mathrm{wt}(e_U,\cdot)\le\mathrm{wt}(e_{k,w},\cdot)$, so Lemma~\ref{w2:transfer} gives $w_k(B_k)\le\mathrm{wt}(e_U,\cdot)^{\otimes f}(B_k)+f/g_k$. By Proposition~\ref{w2:GG} the weight $P=\mathrm{wt}(e_U,\cdot)$ satisfies the hypothesis of Theorem~\ref{w2:badweight}, and the sizes give $g_k\ge2^{12}$ for large $x$. Hence $w_k(B_k)\le W(g_k)+f/g_k$.

Under \eqref{w2:sizes} the terms of $W(g_k)+f/g_k$ are bounded as follows for large $x$. The terms $3f/g_k$ and $2fm(1+\log N)/(g_k-1)$ are at most $4x^{2C_s}(1+C_s\log x)e^{-8x^{\varepsilon'}}$. Since $t\mapsto\log t/\sqrt t$ decreases for $t\ge e^2$, the second term is at most $3x^{C_s}\cdot12x^{\varepsilon'}e^{-4x^{\varepsilon'}}$. The term $2fm/(4m_0)$ is at most $\tfrac12x^{2C_s}e^{-x^{\varepsilon'/6}}$. Finally $\tfrac N2(7/8)^f\le x^{C_s}e^{-x^2/8}$ and $e^{-f/(5\log N)}\le e^{-x^2/(5C_s\log x)}$. Since $\varepsilon'\le1$, every term is at most $x^{2C_s}e^{-x^{\varepsilon'/6}}$ for large $x$, so $\delta\le36\,x^{2C_s}e^{-x^{\varepsilon'/6}}\le e^{-x^{\varepsilon'/7}}$ once $x^{\varepsilon'/6}-x^{\varepsilon'/7}\ge2C_s\log x+\log36$.
\end{proof}

\begin{remark}\label{w2:rem:irr}
The irreducibility certificates \ref{w2:G2} to \ref{w2:G5} are used later, where Jordan\'s theorem enters. Theorem~\ref{w2:I} itself is a count and does not use them.
\end{remark}

\subsubsection{The level idempotents at the auxiliary primes}\label{w2x:sec:wt}
Let $g$ be a prime, $\F=\F_g$, and let $\psi_0\colon\F\to\C^\times$ be a nontrivial additive character. Put $G_n=\GL_n(\F)$, let $U_n\subset G_n$ be the group of upper unitriangular matrices, and index rows and columns by $0,\dots,n-1$. For a subgroup $H$ of a finite group $G$ and a character $\lambda\colon H\to\C^\times$ put
\[
e(H,\lambda)=\lvert H\rvert^{-1}\sum_{h\in H}\overline{\lambda(h)}\,h\in\C[G].
\]
This is a self-adjoint idempotent with $e(H,\lambda)(1)=\lvert H\rvert^{-1}$. For a representation $(\pi,V)$ of $G$ the operator $\pi(e(H,\lambda))$ is the orthogonal projection onto $V^{(H,\lambda)}=\{v\in V:\pi(h)v=\lambda(h)v\ (h\in H)\}$. For $a\in\C[G]$ and a class function $\chi$ on $G$ put $\tr(a,\chi)=\sum_{x\in G}a(x)\chi(x)$. If $\chi=\chi_\pi$ is the character of $\pi$, then $\tr(a,\chi_\pi)=\Tr\pi(a)$, so
\begin{equation}\label{w39:trmult}
\tr\bigl(e(H,\lambda),\chi_\pi\bigr)=\dim V^{(H,\lambda)}=\dim\Hom_H(\lambda,\pi|_H).
\end{equation}
Since the class functions are the linear combinations of the irreducible characters and $\tr(a,\cdot)$ is linear, an identity between the numbers \eqref{w39:trmult} for all irreducible $\pi$ is an identity of $\tr(a,\cdot)$ on all class functions.

\paragraph{Pattern groups.}

Let $I=\{0,\dots,N-1\}$ and let $\prec$ be a strict partial order on $I$, that is, an irreflexive transitive relation. Write $a\preceq b$ for $a=b$ or $a\prec b$. The \emph{pattern group} of $\prec$ is
\[
U_\prec=\{u\in\GL_N(\F):u_{aa}=1\text{ for all }a,\ u_{ab}=0\text{ unless }a\prec b\text{ or }a=b\}.
\]

\begin{lemma}\label{w39:pattern}
(i) $U_\prec$ is a subgroup of $\GL_N(\F)$ of order $g^{\#\{(a,b):a\prec b\}}$.
(ii) Let $E$ be a set of pairs $(a,b)$ with $a\prec b$, and let $c\colon E\to\F$. Suppose that for every $(a,b)\in E$ there is no $m$ with $a\prec m\prec b$. Then $\lambda_{E,c}(u)=\psi_0\bigl(\sum_{(a,b)\in E}c_{ab}u_{ab}\bigr)$ is a character of $U_\prec$.
\end{lemma}

\begin{proof}
(i) For $u,v\in U_\prec$, $(uv)_{ab}=\sum_mu_{am}v_{mb}$, and a nonzero term needs $a\preceq m\preceq b$. By transitivity $a\preceq b$, and if $a=b$ the only term is $m=a$, because $a\prec m\prec a$ would give $a\prec a$. So $uv\in U_\prec$ with unit diagonal. Write $u=1+x$. Every product $x^j$ is supported on pairs $a\prec b$, and $x^N=0$ because a strictly increasing chain in $I$ has at most $N$ elements. Hence $u^{-1}=\sum_{j<N}(-x)^j\in U_\prec$. The entries at the pairs $a\prec b$ are free, which gives the order.
(ii) For $u,v\in U_\prec$ and $(a,b)\in E$ we have $(uv)_{ab}=u_{ab}+v_{ab}+\sum_{a\prec m\prec b}u_{am}v_{mb}=u_{ab}+v_{ab}$.
\end{proof}

\begin{lemma}[one double coset]\label{w39:mackey}
Let $H$ be a finite group with subgroups $H_1,H_2$ and characters $\lambda_i\colon H_i\to\C^\times$ that agree on $H_1\cap H_2$. Suppose that for $i\in\{1,2\}$ and every $h\in H\setminus H_i$ there is $a\in H_i$ with $h^{-1}ah\in H_i$ and $\lambda_i(h^{-1}ah)\ne\lambda_i(a)$. Then $\operatorname{Ind}_{H_1}^H\lambda_1\cong\operatorname{Ind}_{H_2}^H\lambda_2$. Consequently, for every finite group $G\supseteq H$ and every representation $\pi$ of $G$, $\dim\pi^{(H_1,\lambda_1)}=\dim\pi^{(H_2,\lambda_2)}$, and $\tr(e(H_1,\lambda_1),\chi)=\tr(e(H_2,\lambda_2),\chi)$ for every class function $\chi$ on $G$.
\end{lemma}

\begin{proof}
For subgroups $A,B\subseteq H$ and characters $\alpha,\beta$, Mackey's formula (J.-P. Serre, \emph{Linear representations of finite groups}, \S7.4) gives
\[
\bigl\langle\operatorname{Ind}_A^H\alpha,\operatorname{Ind}_B^H\beta\bigr\rangle_H=\sum_{h\in A\backslash H/B}\bigl\langle\alpha,\beta^h\bigr\rangle_{A\cap hBh^{-1}},\qquad\beta^h(y)=\beta(h^{-1}yh).
\]
Each term is $1$ if $\alpha=\beta^h$ on $A\cap hBh^{-1}$ and $0$ otherwise. Take $A=B=H_i$ and $\alpha=\beta=\lambda_i$. For $h\notin H_i$ the element $a$ of the hypothesis lies in $H_i\cap hH_ih^{-1}$ and satisfies $\lambda_i^h(a)=\lambda_i(h^{-1}ah)\ne\lambda_i(a)$, so the term of $h$ vanishes. The double coset of $1$ contributes $1$. Hence $\operatorname{Ind}_{H_i}^H\lambda_i$ is irreducible. Next take $A=H_1$, $B=H_2$. The term of $h=1$ is $1$ because $\lambda_1=\lambda_2$ on $H_1\cap H_2$, and the other terms are nonnegative, so the two irreducible induced representations have a common constituent and are isomorphic. By Frobenius reciprocity and transitivity of induction, $\dim\pi^{(H_i,\lambda_i)}=\dim\Hom_G(\operatorname{Ind}_H^G\operatorname{Ind}_{H_i}^H\lambda_i,\pi)$, which does not depend on $i$. The last claim follows from \eqref{w39:trmult} and linearity.
\end{proof}

\begin{lemma}[exchange]\label{w39:exchange}
Let $\prec_1,\prec_2$ be strict partial orders on $I$ whose union $\prec$ is again a strict partial order. Let $p,q\in I$, $W\subseteq I\setminus\{p,q\}$, and let $E$ be a set of pairs with a function $c\colon E\to\F$. Assume the following.
\begin{enumerate}[label=(X\arabic*),leftmargin=*]
\item\label{w39:X1} The pairs with $a\prec_1b$ and not $a\prec_2b$ are the pairs $(p,w)$, $w\in W$. The pairs with $a\prec_2b$ and not $a\prec_1b$ are the pairs $(w,q)$, $w\in W$.
\item\label{w39:X2} $(p,q)\in E$ and $c_{pq}\ne0$.
\item\label{w39:X3} If $w\in W$, $x\preceq p$ and $w\preceq y$, then $x\prec_1y$, and $(x,y)\in E$ only for $(x,y)=(p,q)$.
\item\label{w39:X4} If $w\in W$, $x\preceq w$ and $q\preceq y$, then $x\prec_2y$, and $(x,y)\in E$ only for $(x,y)=(p,q)$.
\item\label{w39:X5} For $(a,b)\in E$ and $j\in\{1,2\}$ there is no $m$ with $a\prec_jm\prec_jb$.
\end{enumerate}
Then $\lambda=\lambda_{E,c}$ is a character of $U_{\prec_1}$ and of $U_{\prec_2}$, and $\dim\pi^{(U_{\prec_1},\lambda)}=\dim\pi^{(U_{\prec_2},\lambda)}$ for every representation $\pi$ of every finite group containing $U_\prec$.
\end{lemma}

\begin{proof}
By Lemma~\ref{w39:pattern} the groups $H=U_\prec$, $H_1=U_{\prec_1}$, $H_2=U_{\prec_2}$ satisfy $H_1,H_2\subseteq H$, and by \ref{w39:X5} $\lambda$ is a character of $H_1$ and of $H_2$. The two characters are given by one formula, so they agree on $H_1\cap H_2$. We verify the hypothesis of Lemma~\ref{w39:mackey}.

Let $h\in H\setminus H_1$. By \ref{w39:X1} the pairs of $\prec$ outside $\prec_1$ are the pairs $(w,q)$, so $h_{wq}\ne0$ for some $w\in W$. Since $p\prec_1w$, the element $a=1+tE_{pw}$ lies in $H_1$ for every $t\in\F$. Then $(h^{-1}ah)_{xy}=\delta_{xy}+t\,(h^{-1})_{xp}h_{wy}$. A nonzero off-diagonal entry needs $(h^{-1})_{xp}\ne0$ and $h_{wy}\ne0$, so $x\preceq p$ and $w\preceq y$ because $h,h^{-1}\in U_\prec$, and \ref{w39:X3} gives $x\prec_1y$. Hence $h^{-1}ah\in H_1$. By \ref{w39:X3} the only pair of $E$ among these is $(p,q)$, and $(h^{-1})_{pp}=1$, so $\lambda(h^{-1}ah)=\psi_0(c_{pq}t\,h_{wq})$. By \ref{w39:X3} with $(x,y)=(p,w)$ the pair $(p,w)$ lies in $E$ only if $w=q$, which is excluded, so $\lambda(a)=1$. Since $c_{pq}h_{wq}\ne0$ and $\psi_0\ne1$ there is $t$ with $\psi_0(c_{pq}t\,h_{wq})\ne1$.

Let $h\in H\setminus H_2$. Then $h^{-1}\notin H_2$, and by \ref{w39:X1} $(h^{-1})_{pw}\ne0$ for some $w\in W$. Put $a=1+tE_{wq}\in H_2$. The off-diagonal entries of $h^{-1}ah$ are $t\,(h^{-1})_{xw}h_{qy}$ with $x\preceq w$ and $q\preceq y$, so \ref{w39:X4} gives $h^{-1}ah\in H_2$ and $\lambda(h^{-1}ah)=\psi_0(c_{pq}t\,(h^{-1})_{pw})$, while $\lambda(a)=1$ because $(w,q)\in E$ would force $w=p$. Choose $t$ as before.
\end{proof}

\paragraph{Exchange of roots.}

Let $k\ge1$ and $s\ge0$ be integers and $N=k+2s+1$. Put $W=\{s,\dots,s+k-1\}$, $c_0=s+k$, and $O=I\setminus W=\{0,\dots,s-1\}\cup\{c_0,\dots,N-1\}$ with its natural order. Define
\[
U_J=\{u\in U_N:u_{a,c_0}=0\ (a\in W)\},\qquad U'=\{u\in U_N:u_{a,k}=0\ (0\le a<k)\},
\]
\[
\psi_J(u)=\psi_0\Bigl(\sum_{(a,b)\in E}u_{ab}\Bigr),\qquad\psi'(u)=\psi_0\Bigl(\sum_{l=0}^{N-2}u_{l,l+1}\Bigr),
\]
where $E$ is the set of pairs $(a,b)$ that are consecutive in $O$ (the pairs $(l,l+1)$ with $l+1<s$, the pair $(s-1,c_0)$ when $s\ge1$, and the pairs $(l,l+1)$ with $l\ge c_0$) together with the pairs $(l,l+1)$ with $l,l+1\in W$. Put $e_J=e(U_J,\psi_J)$, $e'=e(U',\psi')$ and $e_U=e(U_N,\psi')$. In the construction $(N,k,s)=(r,k,r_S)$ with $r_S=m/2-1$, and $e_J$ is the local factor of the level idempotent at every place above $g_0$ and $g_1$.

\begin{theorem}[exchange of roots]\label{w39:thmA}\label{w2x:exch}
For every class function $\chi$ on $G_N$, $\tr(e_J,\chi)=\tr(e',\chi)$. Equivalently $\dim\pi^{(U_J,\psi_J)}=\dim\pi^{(U',\psi')}$ for every representation $\pi$ of $G_N$.
\end{theorem}

\begin{proof}
For $0\le i\le s$ let $T_i=\{b\in O:b>i\}$ if $i<s$ and $T_s=\{c_0+1,\dots,N-1\}$. Define $a\prec_ib$ if one of the following holds: $a,b\in W$ and $a<b$; $a,b\in O$ and $a<b$; $a\in O$, $a<i$ and $b\in W$; $a\in W$ and $b\in T_i$.

\emph{$\prec_i$ is a strict partial order.} It is contained in no cycle because a relation from $W$ to $O$ ends in $T_i\subseteq\{b>i\}$ and a relation from $O$ to $W$ starts in $\{a<i\}$, so no chain passes from $W$ to $O$ and back to $W$. Transitivity is checked on the four types. For $a\prec_iw\prec_ib$ with $a,b\in O$ and $w\in W$ we have $a<i<b$ when $b<s$, and $a<s\le c_0\le b$ otherwise, so $a<b$. For $w\prec_ib\prec_ib'$ with $b,b'\in O$ we have $b'>b$, and $b\in T_i$ implies $b'\in T_i$ since $T_i$ is an upper set of $O$ (for $i<s$ it is $\{b\in O:b>i\}$, and $T_s$ is the top of $O$). The remaining cases are immediate.

\emph{The character.} Each pair of $E$ is a covering pair of $\prec_i$. For a pair $(a,b)$ consecutive in $O$ no element of $O$ lies between, and $a\prec_iw\prec_ib$ with $w\in W$ needs $a<i$ and $b\in T_i$. If $b<s$ then $b=a+1\le i$ contradicts $b>i$. If $b=c_0$ then $a=s-1<i$ forces $i=s$, and $c_0\notin T_s$. If $b>c_0$ then $a\ge c_0>i$. For a pair $(l,l+1)$ in $W$ an intermediate element of $O$ would need $l\prec_im\prec_il+1$, that is $m\in T_i$ and $m<i$, which is impossible. By Lemma~\ref{w39:pattern} $\psi_J=\lambda_{E,1}$ is a character of every $U_{\prec_i}$.

\emph{The end points.} For $i=s$ the relation $\prec_s$ consists of the pairs $a<b$ other than $(w,c_0)$, $w\in W$, so $U_{\prec_s}=U_J$. For $i=0$ let $\sigma$ be the permutation of $I$ with $\sigma(a)=s+a$ for $a<k$, $\sigma(a)=a-k$ for $k\le a<k+s$, and $\sigma(a)=a$ for $a\ge k+s$. It maps $\{0,\dots,k-1\}$ increasingly onto $W$ and $\{k,\dots,N-1\}$ increasingly onto $O$. A direct check of the four types gives $\sigma(a)\prec_0\sigma(b)$ if and only if $a<b$ and not ($a<k$ and $b=k$), and $(\sigma(a),\sigma(b))\in E$ if and only if $b=a+1\ne k$. With the permutation matrix $P$, $Pe_a=e_{\sigma(a)}$, we get $(P^{-1}uP)_{ab}=u_{\sigma(a)\sigma(b)}$, hence $P^{-1}U_{\prec_0}P=U'$ and $\psi_J(PuP^{-1})=\psi'(u)$ for $u\in U'$ (the entry $u_{k-1,k}$ vanishes on $U'$). Since $\chi$ is a class function, $\tr(e(U_{\prec_0},\psi_J),\chi)=\tr(e',\chi)$.

\emph{One step.} Let $1\le i\le s$. We apply Lemma~\ref{w39:exchange} to $\prec_1=\prec_i$, $\prec_2=\prec_{i-1}$, $p=i-1$, $q=i$ if $i<s$ and $q=c_0$ if $i=s$, the set $W$ and $E$ with $c\equiv1$. Note that $q$ is the successor of $p$ in $O$. The union $\prec$ is the relation of the same four types with $\{a<i\}$ in the third type and $T'=\{b\in O:b\ge i\}$ in the fourth, which is a strict partial order by the argument above. Condition \ref{w39:X1} holds because $\{a\in O:a<i\}\setminus\{a<i-1\}=\{i-1\}$ and $T_{i-1}\setminus T_i=\{q\}$. Condition \ref{w39:X2} holds because $(p,q)$ is consecutive in $O$. For \ref{w39:X3} let $x\preceq p$ and $w\preceq y$ with $w\in W$. Since $p\in O$ and no element of $W$ precedes $p$ (that would need $p\ge i$), $x\in O$ and $x\le i-1$. Either $y\in W$, and then $x\prec_iy$ by the third type, or $y\in O$ with $y\ge i>x$, and then $x\prec_iy$ by the second type. A pair of $E$ of this form joins two consecutive elements $x<y$ of $O$ with $x\le i-1<i\le y$, so it is $(p,q)$. For \ref{w39:X4} let $x\preceq w$ and $q\preceq y$. Then $y\in O$ with $y\ge q\ge i$, because nothing in $W$ follows $q\ge i$ in $\prec$. If $x\in W$ then $y\in T_{i-1}$, and if $x\in O$ then $x<i\le y$. In both cases $x\prec_{i-1}y$, and a pair of $E$ of this form has $x\in O$, $x<i\le q\le y$, $x,y$ consecutive in $O$, so it is $(p,q)$. Condition \ref{w39:X5} was verified above. Lemma~\ref{w39:exchange} gives $\tr(e(U_{\prec_i},\psi_J),\chi)=\tr(e(U_{\prec_{i-1}},\psi_J),\chi)$.

Chaining the steps $i=s,s-1,\dots,1$ with the two end points proves the theorem. For $s=0$ there is no step and $U_{\prec_0}=U_{\prec_s}$.
\end{proof}

\paragraph{Column expansion.}

For $n\ge1$ let $P_n\subseteq G_n$ be the mirabolic subgroup of matrices whose last row is $(0,\dots,0,1)$. Embed $G_{n-1}\subseteq P_n$ by $h\mapsto\operatorname{diag}(h,1)$ and let $Y_n=\{y(x):x\in\F^{n-1}\}$, where $y(x)$ is the unipotent matrix with last column $(x,1)$. Then $P_n=G_{n-1}\ltimes Y_n$, $hy(x)h^{-1}=y(hx)$, and $U_n=U_{n-1}\ltimes Y_n$. For $0\le i\le n-1$ let
\[
\psi^{(i)}_n(u)=\psi_0\Bigl(\sum_{0\le l\le n-2,\ l\ne i-1}u_{l,l+1}\Bigr)\qquad(u\in U_n),
\]
so that $\psi^{(0)}_n=\psi_n$ is the generic character. For a representation $\tau$ of $P_n$ put
\[
A(\tau)=\dim\tau^{(U_{n-1},\psi_{n-1})},\qquad B_i(\tau)=\dim\tau^{(U_n,\psi^{(i)}_n)}\quad(0\le i\le n-1).
\]
For a representation $\rho$ of $G_n$ put $\mathcal W(\rho)=\dim\rho^{(U_n,\psi_n)}$.

\begin{lemma}[orbits of the unipotent group]\label{w39:orbits}
Let $m\ge1$ and let $\sigma$ be a representation of $P_m$. For $0\le j\le m-2$ let $S_j=\{u\in U_{m-1}:u_{jl}=0\ (l>j)\}$, and for $c\in\F^\times$ let $\chi'_{j,c}(s\,y(x))=\psi_{m-1}(s)\psi_0(cx_j)$ on $S_jY_m$. Then $\chi'_{j,c}$ is a character and
\[
A(\sigma)=B_{m-1}(\sigma)+\sum_{j=0}^{m-2}\sum_{c\in\F^\times}X_{j,c}(\sigma),\qquad X_{j,c}(\sigma)=\dim\sigma^{(S_jY_m,\chi'_{j,c})}.
\]
\end{lemma}

\begin{proof}
For a row vector $\mu\in\F^{m-1}$ let $\sigma_\mu=\{v:\sigma(y(x))v=\psi_0(\mu x)v\}$. Since $Y_m$ is abelian, $\sigma=\bigoplus_\mu\sigma_\mu$, and $\sigma(h)\sigma_\mu=\sigma_{\mu h^{-1}}$ for $h\in G_{m-1}$. The group $U_{m-1}$ acts on row vectors by $\mu\mapsto\mu u^{-1}$. Its orbits are $\{0\}$ and, for $0\le j\le m-2$ and $c\in\F^\times$, the sets $O_{j,c}$ of $\mu$ with $\mu_l=0$ for $l<j$ and $\mu_j=c$. Indeed $\mu u^{-1}$ has the same first nonzero coordinate as $\mu$, and $c\,e_j^{T}(1+\sum_{l>j}c^{-1}\mu_lE_{jl})=\mu$ for $\mu\in O_{j,c}$. The stabilizer of $c\,e_j^T$ is $S_j$, which therefore fixes the character $\psi_0(cx_j)$ of $Y_m$, so $\chi'_{j,c}$ is a character of $S_j\ltimes Y_m$. The sum $\sigma_{O}=\bigoplus_{\mu\in O}\sigma_\mu$ over one orbit is $U_{m-1}$-stable and isomorphic to $\operatorname{Ind}_{S_j}^{U_{m-1}}\sigma_{ce_j}$, so by Frobenius reciprocity $\dim\sigma_O^{(U_{m-1},\psi_{m-1})}=\dim\sigma_{ce_j}^{(S_j,\psi_{m-1})}=X_{j,c}(\sigma)$. The orbit $\{0\}$ gives $\sigma^{Y_m}$, and $\dim(\sigma^{Y_m})^{(U_{m-1},\psi_{m-1})}=B_{m-1}(\sigma)$ because $\psi^{(m-1)}_m(uy)=\psi_{m-1}(u)$.
\end{proof}

\begin{lemma}[induction from the mirabolic subgroup]\label{w39:indP}
Let $m\ge1$ and let $\sigma$ be a representation of $P_m$. For $0\le j\le m-1$ let $Y_m^{<j}=\{y(x):x_l=0\ (l\ge j)\}$ and $H_j=U_{m-1}Y_m^{<j}$, and for $c\in\F^\times$ let
\[
\chi_{j,c}(u\,y(x))=\psi_0\Bigl(\sum_{0\le l\le m-3,\ l\ne j-1}u_{l,l+1}+c\,x_{j-1}\Bigr),
\]
with $x_{-1}=0$. Then $\chi_{j,c}$ is a character of the group $H_j$ and
\[
\mathcal W\bigl(\operatorname{Ind}_{P_m}^{G_m}\sigma\bigr)=\sum_{j=0}^{m-1}\sum_{c\in\F^\times}D_{j,c}(\sigma),\qquad D_{j,c}(\sigma)=\dim\sigma^{(H_j,\chi_{j,c})}.
\]
\end{lemma}

\begin{proof}
By Mackey's formula and Frobenius reciprocity,
\[
\mathcal W(\operatorname{Ind}_{P_m}^{G_m}\sigma)=\sum_{x\in U_m\backslash G_m/P_m}\dim\sigma^{(x^{-1}(U_m\cap xP_mx^{-1})x,\ \psi_m(x\,\cdot\,x^{-1}))}.
\]
The map $gP_m\mapsto e_{m-1}^Tg^{-1}$ identifies $G_m/P_m$ with the nonzero row vectors, and $U_m$ acts by $\mu\mapsto\mu u^{-1}$. As in Lemma~\ref{w39:orbits} the orbits are indexed by $(j,c)$ with $0\le j\le m-1$ and $c\in\F^\times$, with representative $c\,e_j^T$ and stabilizer $U_{(j)}=\{u\in U_m:u_{jl}=0\ (l>j)\}$. Let $\beta$ be the bijection of $I$ with $\beta(l)=l$ for $l<j$, $\beta(j)=m-1$ and $\beta(l)=l-1$ for $l>j$, let $w$ be the permutation matrix with $we_l=e_{\beta(l)}$, and let $t=\operatorname{diag}(1,\dots,1,c)$. Then $x$ with $x^{-1}=tw$ satisfies $e_{m-1}^Tx^{-1}=c\,e_j^T$, so $U_m\cap xP_mx^{-1}=U_{(j)}$. Since $(wzw^{-1})_{\beta(a)\beta(b)}=z_{ab}$ and $\beta$ is increasing on $I\setminus\{j\}$, the group $wU_{(j)}w^{-1}$ consists of the matrices of $U_m$ whose block on $\{0,\dots,m-2\}$ is arbitrary in $U_{m-1}$ and whose last column is supported on the rows $l<j$ (these come from the entries $(l,j)$). Conjugation by $t$ rescales the last column, so $x^{-1}U_{(j)}x=H_j$. For $z\in H_j$ put $u=xzx^{-1}=w^{-1}(t^{-1}zt)w$, so $u_{ab}=(t^{-1}zt)_{\beta(a)\beta(b)}$ and $(t^{-1}zt)_{a,m-1}=c\,z_{a,m-1}$ for $a<m-1$. The superdiagonal entries of $u$ are $u_{l,l+1}=z_{l,l+1}$ for $l+1<j$, $u_{j-1,j}=c\,z_{j-1,m-1}$, $u_{j,j+1}=0$, and $u_{l,l+1}=z_{l-1,l}$ for $l>j$. Hence $\psi_m(xzx^{-1})=\chi_{j,c}(z)$.
\end{proof}

\begin{lemma}\label{w39:DX}
In the notation of Lemmas~\ref{w39:orbits} and~\ref{w39:indP}, $D_{j+1,c}(\sigma)=X_{j,c}(\sigma)$ for $0\le j\le m-2$ and $c\in\F^\times$.
\end{lemma}

\begin{proof}
On $I=\{0,\dots,m-1\}$ let $a\prec_1b$ mean $a<b\le m-2$, or $b=m-1$ and $a\le j$, and let $a\prec_2b$ mean $a<b\le m-2$ and $a\ne j$, or $b=m-1>a$. Then $U_{\prec_1}=H_{j+1}$ and $U_{\prec_2}=S_jY_m$. Let $E$ consist of the pairs $(l,l+1)$ with $0\le l\le m-3$, $l\ne j$, with coefficient $1$, and of $(j,m-1)$ with coefficient $c$. On $H_{j+1}$ the character $\lambda_{E,c}$ is $\chi_{j+1,c}$. On $S_jY_m$ it is $\chi'_{j,c}$, because $(s\,y(x))_{j,m-1}=(sx)_j=x_j$ when the row $j$ of $s$ is $e_j^T$. If $j=m-2$ the two relations coincide. Otherwise we apply Lemma~\ref{w39:exchange} with $p=j$, $q=m-1$ and $W=\{j+1,\dots,m-2\}$. The union of the two relations is the natural order $<$. \ref{w39:X1} holds by definition. \ref{w39:X3}: if $x\le j<w\le y$ then $x\prec_1y$, and a pair $(x,x+1)$ with $x\le j<x+1$ has $x=j$ and is not in $E$, so only $(j,m-1)$ occurs. \ref{w39:X4}: if $x\le w$ and $y\ge m-1$ then $y=m-1$ and $x\prec_2y$, and $(x,m-1)\in E$ only for $x=j$. \ref{w39:X5}: the pairs $(l,l+1)$ are covering pairs of $<$. An element $b$ with $j\prec_1b\prec_1m-1$ would satisfy $b>j$ and $b\le j$. An element $b$ with $j\prec_2b$ equals $m-1$. Lemma~\ref{w39:exchange} applies since $P_m\supseteq U_m=U_\prec$.
\end{proof}

\begin{proposition}\label{w39:propM}
For every $n\ge1$ and every representation $\tau$ of $P_n$,
\[
A(\tau)=B_{n-1}(\tau)+\sum_{i=0}^{n-2}(g-1)g^{\,n-2-i}B_i(\tau).
\]
\end{proposition}

\begin{proof}
For $n=1$ both sides equal $\dim\tau$. Let $n\ge2$ and assume the claim for $n-1$. Put $m=n-1$ and let $\theta(y(x))=\psi_0(x_{n-2})$, a character of $Y_n$ whose stabilizer in $G_{m}$ for the action $\mu\mapsto\mu h^{-1}$ on row vectors is $P_m$. Clifford theory as in the proof of Lemma~\ref{w39:orbits}, now with the transitive action of $G_m$ on nonzero row vectors, gives $\tau|_{G_m}\cong\tau^{Y_n}\oplus\operatorname{Ind}_{P_m}^{G_m}\sigma$ with $\sigma=\tau^{(Y_n,\theta)}$, a representation of $P_m$. Since $\psi^{(n-1)}_n$ is trivial on $Y_n$, $\mathcal W(\tau^{Y_n})=B_{n-1}(\tau)$. For $0\le i\le n-2$ the character $\psi^{(i)}_n$ restricts to $\theta$ on $Y_n$ and to $\psi^{(i)}_m$ on $U_m$, so $B_i(\tau)=B_i(\sigma)$, the right side computed in $P_m$.

By Lemmas~\ref{w39:indP}, \ref{w39:DX} and~\ref{w39:orbits},
\[
\mathcal W(\operatorname{Ind}_{P_m}^{G_m}\sigma)=\sum_cD_{0,c}(\sigma)+\sum_{j=0}^{m-2}\sum_cX_{j,c}(\sigma)=(g-1)A(\sigma)+\bigl(A(\sigma)-B_{m-1}(\sigma)\bigr),
\]
because $H_0=U_{m-1}$ and $\chi_{0,c}=\psi_{m-1}$. By the induction hypothesis $gA(\sigma)-B_{m-1}(\sigma)=(g-1)B_{m-1}(\sigma)+\sum_{i\le m-2}(g-1)g^{m-1-i}B_i(\sigma)=\sum_{i=0}^{m-1}(g-1)g^{m-1-i}B_i(\sigma)$. Adding $B_{n-1}(\tau)$ and using $B_i(\sigma)=B_i(\tau)$ and $m-1-i=n-2-i$ gives the claim.
\end{proof}

For a representation $\pi$ of $G_N$ and $0\le i\le N-1$ put $W_i(\pi)=\dim\pi^{(U_N,\psi^{(i)}_N)}$, so that $W_0(\pi)=\Tr\pi(e_U)$.

\begin{theorem}[column expansion]\label{w39:thmB}\label{w2x:der}
Let $1\le k<N$. For every representation $\pi$ of $G_N$,
\[
\Tr\pi(e')=\sum_{i=0}^{k-1}\bigl(g^{k-i}-g^{k-i-1}\bigr)W_i(\pi)+W_k(\pi).
\]
\end{theorem}

\begin{proof}
Put $n=k+1\le N$ and embed $G_n\subseteq G_N$ in the upper left corner. Let $Q\subseteq U_N$ be the group of $u\in U_N$ whose upper left $n\times n$ block is $1$, and $\psi_Q(z)=\psi_0(\sum_{l=n-1}^{N-2}z_{l,l+1})$, a character of $Q$. For $p\in P_n$ the conjugate $pzp^{-1}$ replaces the upper right block $X$ of $z$ by $pX$ and keeps the lower right block, and the row $n-1$ of $pX$ equals that of $X$, so $P_n$ normalizes $Q$ and fixes $\psi_Q$. Hence $\tau=\pi^{(Q,\psi_Q)}$ is a representation of $P_n$. Every $u\in U_N$ factors uniquely as $u=u_0z$ with $u_0\in U_n$ and $z\in Q$, and $u_0$ is the upper left block of $u$. Thus $U'=U_kQ$ with $U_k\subseteq G_k\subseteq P_n$, and $\psi'(u_0z)=\psi_k(u_0)\psi_Q(z)$, because the row $k=n-1$ of $u_0$ is $e_k^T$. Hence $\Tr\pi(e')=\dim\tau^{(U_k,\psi_k)}=A(\tau)$. In the same way $\psi^{(i)}_N(u_0z)=\psi^{(i)}_n(u_0)\psi_Q(z)$ for $0\le i\le k$, since the omitted entry $(i-1,i)$ lies in the upper left block, so $W_i(\pi)=B_i(\tau)$. Proposition~\ref{w39:propM} with $n=k+1$ and $(g-1)g^{k-1-i}=g^{k-i}-g^{k-i-1}$ ends the proof.
\end{proof}

\begin{corollary}[Hypothesis~\ref{w2:H-wt}]\label{w39:Hwt}\label{w2x:idem}\label{w2x:cmp}
Let $(N,k,s)=(r,k,r_S)$ with $k\ge1$. The element $e_J$ is a self-adjoint idempotent with $e_J(1)=g^{k-N(N-1)/2}>0$, and for every family $(E_c)$ of orthogonal central idempotents of $\C[G_N]$ with $\sum_cE_c=1$, in particular for the label idempotents of Section~\ref{w2:sec:irr},
\[
(1-1/g)\,\mathrm{wt}(e_U,c)\le\mathrm{wt}(e_J,c)\qquad\text{for every }c.
\]
\end{corollary}

\begin{proof}
$U_J$ is a subgroup (Lemma~\ref{w39:pattern}, it is $U_{\prec_s}$) defined in $U_N$ by $k$ vanishing entries above the diagonal, so $e_J(1)=\lvert U_J\rvert^{-1}=g^ke_U(1)$. By Theorems~\ref{w39:thmA} and~\ref{w39:thmB} and $W_i\ge0$, $\Tr\pi(e_J)\ge(g^k-g^{k-1})W_0(\pi)=g^k(1-1/g)\Tr\pi(e_U)$ for every irreducible $\pi$. For $a\in\C[G_N]$, $(aE_\pi)(1)=\lvert G_N\rvert^{-1}\dim\pi\,\Tr\pi(a)$. Summing over the $\pi$ with $E_\pi E_c=E_\pi$ gives $(e_JE_c)(1)\ge g^k(1-1/g)(e_UE_c)(1)$, and dividing by $e_J(1)=g^ke_U(1)$ gives the claim.
\end{proof}

\begin{lemma}\label{w2x:trace}
Let $a\in\C[G_N]$ be an idempotent and $c$ a label. Then
\[
(aE_c)(1)=\frac1{\lvert G_N\rvert}\sum_{\pi\text{ of label }c}\dim\pi\cdot\operatorname{rk}\pi(a)\ \ge0,\qquad\sum_c(aE_c)(1)=a(1).
\]
In particular Lemma~\ref{w2:wtprob} holds for every idempotent $a$ with $\operatorname{Re}a(1)>0$, self-adjoint or not.
\end{lemma}

\begin{proof}
The central idempotent of $\pi$ is $E_\pi=\frac{\dim\pi}{\lvert G_N\rvert}\sum_x\overline{\chi_\pi(x)}\,x$, and $\overline{\chi_\pi(x^{-1})}=\chi_\pi(x)$, so $(aE_\pi)(1)=\sum_xa(x)E_\pi(x^{-1})=\frac{\dim\pi}{\lvert G_N\rvert}\Tr\pi(a)$. The trace of an idempotent operator is its rank. Summing over the $\pi$ of label $c$ gives the first identity. Summing over all $c$ gives $\sum_\pi\frac{\dim\pi}{\lvert G_N\rvert}\Tr\pi(a)$, which is the trace of $a$ in the left regular representation divided by $\lvert G_N\rvert$, that is $a(1)$.
\end{proof}

\subsubsection{Rows, labels and the cylinder idempotents}\label{w2x:sec:rows}

Let $(\tau)$ be the orthonormal eigenbasis of $\mathcal H$ of Section~\ref{w2:sec:irr}. For $k\in\{0,1\}$ let $\mathcal G_k=\prod_{w\in\Pi_k}\GL_N(\F_{g_k})$ be the quotient of the level at $g_k$, and for a set $S$ of label tuples at $k$ let $\Pi_S=\sum_{t\in S}\bigotimes_{w\in\Pi_k}E_{t_w}$, a central idempotent of $\C[\mathcal G_k]$. The level idempotent factors as $e=e^{(k)}\otimes\bigotimes_{w\in\Pi_k}e_{k,w}$, with $e^{(k)}$ supported away from $g_k$, and we put $e_S=e\,\Pi_S$.

\begin{hypothesis}[label dictionary]\label{w2x:H-lab}
(i) For $k\in\{0,1\}$ and every set $S$ of label tuples at $k$, the space of automorphic forms of Section~\ref{w2:sec:statements} defined with $e_S$ in place of $e$ has dimension $\#\{\tau:\mathrm{lab}_k(\tau)\in S\}$. (ii) There is an injective map $\iota$ from an index set of the rows of $R$ to the eigenbasis whose image is the set of $\tau$ whose label tuples at $g_0$ and at $g_1$ are both good, and $\mathrm{lab}_k(\iota(\sigma))$ is the label tuple of the row $\sigma$.
\end{hypothesis}

Indeed the mutually orthogonal central idempotents $E_t$ indexed by pairs of label tuples sum to $1$ and commute with $e$. Hence $\mathcal H$ is the orthogonal direct sum of the spaces $R(E_t)\mathcal H$. An orthonormal basis adapted to this sum gives (i) by taking dimensions. For (ii), suppose each row basis vector lies in a good summand and every good summand lies in the row space. Then the row space is precisely the sum over good $t$. Choosing a bijection of the two bases separately in each summand gives the required row indices.

\begin{lemma}\label{w2x:DS}
The element $e_S$ is a self-adjoint idempotent, and $D^*(e_S)=D^*\,w_k(S)$, where $D^*(e')$ denotes the product of the formal degree, the volume of $[G]$ and $\operatorname{Re}e'(1)$.
\end{lemma}

\begin{proof}
$\Pi_S$ is a self-adjoint central idempotent of $\C[\mathcal G_k]$, so it commutes with $\bigotimes_we_{k,w}$, and the product of two commuting self-adjoint idempotents is a self-adjoint idempotent. The coefficient at $1$ of a tensor product is the product of the coefficients at $1$, so
\[
e_S(1)=e^{(k)}(1)\sum_{t\in S}\prod_{w\in\Pi_k}(e_{k,w}E_{t_w})(1)=e(1)\sum_{t\in S}\prod_{w\in\Pi_k}\frac{(e_{k,w}E_{t_w})(1)}{e_{k,w}(1)}.
\]
By Lemma~\ref{w2x:trace} every factor is real, and $e_{k,w}(1)>0$ by Corollary~\ref{w39:Hwt}, so the right side is $e(1)\,w_k(S)$ and $D^*(e_S)=D^*w_k(S)$.
\end{proof}

\subsubsection{The multiplicity formula}\label{w2x:sec:count}

We prove Hypotheses~\ref{w2:H-cnt} and~\ref{w2:H-DIM} from three analytic statements about the pre-trace formula, which are not proved here. Let $\varphi_\infty(y)=d_\infty\prod_{v}\overline{\langle\pi_v(y_v)w_v,w_v\rangle}$, where $w_v$ is a unit lowest $K$-type vector of the discrete series $\pi_v$ at the real place $v$ and $d_\infty$ is the product of the formal degrees. For a self-adjoint idempotent $e'$ of the Hecke algebra of $G(\Af)$ put $\varphi=\varphi_\infty\otimes e'$, let $R(\varphi)$ be its right convolution operator on $L^2([G])$, let $D(e')$ be the dimension of the space of Section~\ref{w2:sec:statements} defined with $e'$, and let $D^*(e')=d_\infty\vol([G])\operatorname{Re}e'(1)$. Let $\mathfrak F$ be a fundamental domain for $G(F)$ in $G(\A)$ and $n=\dim V=r$.

\begin{hypothesis}[pre-trace inputs, not proved here]\label{w2x:H-pre}
There are absolute constants $c_1>0$, $c_b>0$, and for every $A\ge1$ constants $n_0$, $C_A$ and $A_1$, such that the following holds for every datum with $n\ge n_0$ whose level and Harish-Chandra parameters satisfy the bounds of the construction with the constant $A$, whose idempotent $e'$ is of positive type and supported on a compact open subgroup that is unipotent modulo an auxiliary prime split in $E$ and larger than $4n$, and whose Harish-Chandra parameters have gap at least $h$ at every real place, with
\begin{equation}\label{w2x:gap}
h\ge C_A\,n^{19}f^5\log(b\,\mathrm{rd}_E),\qquad\log f\le\sqrt n,\qquad\log\log b\le c_b\,nf\log n.
\end{equation}
\begin{enumerate}[label=(\roman*),leftmargin=*]
\item $R(\varphi)$ is an orthogonal projection of rank $D(e')$.
\item For almost every $x\in\mathfrak F$ the series $\sum_{\gamma\in G(F)}\varphi(x^{-1}\gamma x)$ converges absolutely, and its sum is $\sum_j\lvert F_j(x)\rvert^2$ for every orthonormal basis $(F_j)$ of the range of $R(\varphi)$.
\item There is a measurable $M\subset\mathfrak F$ such that, with $\eta=\exp(-c_1nf\log n)$,
\[
\begin{gathered}
\vol(\mathfrak F\setminus M)\le\eta\vol(\mathfrak F),\\
\Bigl\lvert\int_M\sum_{\gamma\ne1}\varphi(x^{-1}\gamma x)\,dx\Bigr\rvert\le A_1\eta D^*(e'),\\
\Bigl\lvert\int_{\mathfrak F\setminus M}\sum_{\gamma}\varphi(x^{-1}\gamma x)\,dx\Bigr\rvert\le A_1e^{-n^2f}D(e').
\end{gathered}
\]
\end{enumerate}
\end{hypothesis}

\begin{theorem}\label{w2x:count}
Under Hypothesis~\ref{w2x:H-pre} put $c=c_1/2$. There is $n_1(A)$ such that every datum as in Hypothesis~\ref{w2x:H-pre} with $n\ge n_1(A)$ satisfies $\lvert D(e')-D^*(e')\rvert\le\exp(-c\,nf\log n)\,D^*(e')$.
\end{theorem}

\begin{proof}
Write $D=D(e')$, $D^*=D^*(e')$ and $K(x)=\sum_\gamma\varphi(x^{-1}\gamma x)$. By (i) and (ii), integrating $\sum_j\lvert F_j\rvert^2$ over $\mathfrak F$ gives $D=\int_{\mathfrak F}K$. The term $\gamma=1$ is $\varphi(1)=d_\infty e'(1)$, because $\lVert w_v\rVert=1$. Since $e'$ is a self-adjoint idempotent, $e'(1)=\lVert e'\delta_1\rVert^2$ is real and nonnegative, so $\varphi(1)\vol(\mathfrak F)=D^*$. Splitting $\mathfrak F=M\cup(\mathfrak F\setminus M)$ gives
\[
D-D^*=-\varphi(1)\vol(\mathfrak F\setminus M)+\int_M\sum_{\gamma\ne1}\varphi(x^{-1}\gamma x)\,dx+\int_{\mathfrak F\setminus M}K,
\]
and $\varphi(1)\vol(\mathfrak F\setminus M)\le\eta D^*$, so by (iii) $\lvert D-D^*\rvert\le(1+A_1)\eta D^*+A_1\theta D$ with $\theta=e^{-n^2f}$. Take $n$ so large that $A_1\theta\le\frac12$ and $n\ge c_1\log n$. Then $D\le2(1+(1+A_1)\eta)D^*\le2(2+A_1)D^*$, and $\theta\le\eta$, so $\lvert D-D^*\rvert\le B\eta D^*$ with $B=1+A_1+2A_1(2+A_1)$. Finally $B\eta\le\eta^{1/2}=\exp(-c\,nf\log n)$ once $c\,nf\log n\ge\log B$.
\end{proof}

\begin{corollary}\label{w2x:cnt}
Assume Hypotheses~\ref{w2x:H-lab} and~\ref{w2x:H-pre}, and that the construction satisfies \eqref{w2x:gap} with $n=r$ for $e$ and for every $e_S$, the unipotent condition for $e_S$ being imposed at the prime $g_{1-k}$. Then for $x$ large Hypothesis~\ref{w2:H-cnt} holds with $\varepsilon_C=\exp(-c\,rf\log r)\le\frac12$, and Hypothesis~\ref{w2:H-DIM} holds with the same $c$.
\end{corollary}

\begin{proof}
The set $S$ of all tuples has $e_S=e$ and $D(e)=D_{\mathrm{all}}$, so Theorem~\ref{w2x:count} gives Hypothesis~\ref{w2:H-DIM} and the first clause of Hypothesis~\ref{w2:H-cnt}. For general $S$, Lemma~\ref{w2x:DS} shows that $e_S$ is a self-adjoint idempotent with $D^*(e_S)=D^*w_k(S)$, and Hypothesis~\ref{w2x:H-lab}(i) gives $D(e_S)=\#\{\tau:\mathrm{lab}_k(\tau)\in S\}$. Theorem~\ref{w2x:count} applied to $e_S$ is the second clause. Since $r\ge100$ and $f\ge x^2$, $\exp(-c\,rf\log r)\le\frac12$ for $x$ large.
\end{proof}

\subsubsection{Sizes}\label{w2x:sec:sizes}

\begin{lemma}\label{w2x:sizes}
Let $0<\varepsilon<\frac1{52}$, $J\ge1$ and $\varepsilon_z=\varepsilon/J$. Suppose $x^{2+8\varepsilon}\le n$, $x^{9/2+10\varepsilon}\le k\le2x^{9/2+10\varepsilon}$, $m\le2n\le k$, $g_0,g_1\ge(kb)^8$ and $\min(q,d)\ge\exp(x^{\varepsilon_z/6})$. Then \eqref{w2:sizes} holds with $\varepsilon'=\varepsilon_z$, $C_s=6$, $N=r=k+m-1$ and $f=2n$, for $x\ge4$ large in terms of $\varepsilon$.
\end{lemma}

\begin{proof}
$g_k\ge(kb)^8\ge b^8=e^{8x^\varepsilon}\ge e^{8x^{\varepsilon_z}}$. Next $f=2n\ge x^2$, and $f\le k\le2x^{9/2+10\varepsilon}\le x^5$ since $\varepsilon<\frac1{52}$ and $x$ is large, so $f\le x^6$ and $m\le f\le x^6$. The rank satisfies $k\le r\le2k\le4x^{9/2+10\varepsilon}\le x^6$, and $r\ge k\ge x^{9/2}\ge100$. The last condition is the hypothesis on $q$ and $d$.
\end{proof}



\subsection{Near orthogonality of Bessel Poincar\'e series}\label{w2:sec:S}

\subsubsection{Assembly}

\begin{proposition}\label{w2:assembly}
Let $C_0\ge1$, $\Gamma_0>1$, $D^*>0$, $\eta_0\ge0$, $\xi\ge0$, $\eta_1\ge0$, and assume \eqref{w2:window}. Let $K_{uu'}=K^0_{uu'}+K^{\deg}_{uu'}$ with
\begin{enumerate}[label=(\roman*),leftmargin=*]
\item $K^{\deg}_{uu'}=0$ for $u\ne u'$,
\item $\lVert K^{\deg}_{uu}\rVert_\HS^2\le\xi d^S_u$,
\item $\lVert K^0_{uu'}\rVert_\HS^2\le(1+\eta_0)d^S_ud^S_{u'}/D^*$ for $u\ne u'$,
\item $\lVert K^0_{uu}-1\rVert_\HS^2\le\Gamma_0(d^S_u)^2/D^*$,
\item $\lvert D_{\mathrm{all}}-D^*\rvert\le\eta_1D^*$.
\end{enumerate}
Then the conclusion of Theorem~\ref{w2:S} holds with these $\eta_0,\eta_1$ and $e_0=2C_0/\Gamma_0+2\xi$.
\end{proposition}

\begin{proof}
Off the diagonal use (i) and (iii). On the diagonal, $\lVert A+B\rVert^2\le2\lVert A\rVert^2+2\lVert B\rVert^2$ with $A=K^0_{uu}-1$, $B=K^{\deg}_{uu}$ and (iv), (ii) give $\lVert K_{uu}-1\rVert_\HS^2\le2\Gamma_0(d^S_u)^2/D^*+2\xi d^S_u$, and \eqref{w2:window} gives $\Gamma_0(d^S_u)^2/D^*\le(C_0/\Gamma_0)d^S_u$.
\end{proof}

\subsubsection{The pair bound}

Unfolding $\langle\mathcal P_uT,\mathcal P_{u'}T'\rangle$ as in Proposition~\ref{w2:unfold}, with $\theta=\overline{P_{T}}$, writes $K_{uu'}$ as a sum over $S_u(F)\backslash G(F)/S_{u'}(F)$ of archimedean orbital integrals of $W_u\otimes\overline{W_{u'}}$. We organize the terms as follows, and the next theorem is the Hilbert space lemma that closes the off-diagonal clause.

\begin{definition}[term families]\label{w2:terms}
A term family consists of a countable set of terms $t$, each with an operator $\mathrm{op}(t)$ in a Hilbert space of operators, a vector $\mathrm{vec}(t)$ in a Hilbert space $H_1$, a row $\mathrm{row}(t)$, a Gram value $G(t)$ and two flags, $\mathrm{isId}(t)$ and $\mathrm{gen}(t)$; a finite set $I$ with a map $\iota$ from rows to $I$; a hermitian matrix $(\Gamma_{ij})_{i,j\in I}$ and a number $\Gamma_1>0$; for pairs of terms in one row a corank $s_0(t,t')\in\mathbb Z_{\ge0}$ and a flag $\mathrm{scal}(t,t')$. Put $\mathrm{id}=\{\mathrm{isId}\}$, $\mathrm{deg}=\{G=0,\neg\mathrm{isId}\}$, $\mathrm{ng}=\{G\ne0,\neg\mathrm{gen}\}$, $\mathrm{gen}=\{G\ne0,\mathrm{gen}\}$, $\mathrm{part}(S)=\sum_{t\in S}\mathrm{op}(t)$, $v_i(S)=\sum_{t\in S,\ \iota(\mathrm{row}\,t)=i}\mathrm{vec}(t)$. For a set $Q$ of pairs in one row put $\Sigma(Q)=\Gamma_1\sum_Q\operatorname{Re}\langle\mathrm{vec}(t),\mathrm{vec}(t')\rangle$ and $\lvert\Sigma\rvert(Q)=\Gamma_1\sum_Q\lvert\langle\mathrm{vec}(t),\mathrm{vec}(t')\rangle\rvert$. Let $\mathrm{big}(S)$ be the pairs in $S\times S$ in one row with $s_0\ge2$, $\mathrm{one}(S)$ those with $s_0=1$, split into $\mathrm{sc}(S)$ (with $\mathrm{scal}$) and $\mathrm{nsc}(S)$.
\end{definition}

\begin{theorem}[pair bound]\label{w2:pair}
Assume $\sum_t\lVert\mathrm{op}(t)\rVert<\infty$, $\sum_t\lVert\mathrm{vec}(t)\rVert<\infty$, and
\begin{enumerate}[label=(\alph*),leftmargin=*]
\item $\lVert\mathrm{part}(S)\rVert^2=\operatorname{Re}\sum_{i,j}\Gamma_{ij}\langle v_i(S),v_j(S)\rangle$ for $S=\mathrm{gen}$ and $S=\mathrm{ng}$,
\item $P_g\Gamma_1\cdot1-(\Gamma_{ij})$ is positive semidefinite for some $P_g\ge0$,
\item $\langle\mathrm{vec}(t),\mathrm{vec}(t')\rangle=0$ if $\iota(\mathrm{row}\,t)=\iota(\mathrm{row}\,t')$ and $\mathrm{row}\,t\ne\mathrm{row}\,t'$,
\item $s_0(t,t')\ge1$ if $t,t'$ lie in one row and $G(t),G(t')\ne0$,
\item $\mathrm{part}(\mathrm{id})=0$ and $\mathrm{part}(\mathrm{deg})=0$,
\item with $dd\ge0$: $\lvert\Sigma\rvert(\mathrm{big}(\mathrm{gen}\cup\mathrm{ng}))\le\varepsilon\,dd$, $\lvert\Sigma\rvert(\mathrm{nsc}(\mathrm{gen}))\le\varepsilon\,dd$, $\Sigma(\mathrm{sc}(\mathrm{gen}))\le A\,dd$ and $\Sigma(\mathrm{one}(\mathrm{ng}))\le\varepsilon'\,dd$.
\end{enumerate}
Then $\lVert\sum_t\mathrm{op}(t)\rVert^2\le\bigl(\sqrt{P_g(A+2\varepsilon)}+\sqrt{P_g(\varepsilon+\varepsilon')}\bigr)^2dd$.
\end{theorem}

\begin{proof}
By (b), $\operatorname{Re}\sum_{ij}\Gamma_{ij}\langle a_i,a_j\rangle\le P_g\Gamma_1\sum_i\lVert a_i\rVert^2$ for all vectors $a_i$. With (a), $\lVert\mathrm{part}(S)\rVert^2\le P_g\Phi(S)$ where $\Phi(S)=\Gamma_1\sum_i\lVert v_i(S)\rVert^2$. Expanding $\lVert v_i(S)\rVert^2$ and using (c), only pairs in one row remain, so $\Phi(S)=\Sigma(\text{pairs of }S\text{ in one row})$. By (d), $\Phi(\mathrm{gen})=\Sigma(\mathrm{big}(\mathrm{gen}))+\Sigma(\mathrm{nsc}(\mathrm{gen}))+\Sigma(\mathrm{sc}(\mathrm{gen}))\le(A+2\varepsilon)dd$ and $\Phi(\mathrm{ng})=\Sigma(\mathrm{big}(\mathrm{ng}))+\Sigma(\mathrm{one}(\mathrm{ng}))\le(\varepsilon+\varepsilon')dd$, using $\Sigma\le\lvert\Sigma\rvert$ and monotonicity of $\lvert\Sigma\rvert$. By (e), $\sum_t\mathrm{op}(t)=\mathrm{part}(\mathrm{gen})+\mathrm{part}(\mathrm{ng})$, and the triangle inequality ends the proof.
\end{proof}

\begin{lemma}\label{w2:partition}
For $g\ge3$ and every $k$, $\prod_{i=1}^k(1-g^{-i})^{-1}\le1+3/g$.
\end{lemma}

\begin{proof}
With $y=1/g$, $\prod_{i\ge1}(1-y^i)^{-1}=\sum_{n\ge0}p(n)y^n$, and $p(n)\le2^{n-1}$ for $n\ge1$ because sorting the parts of a composition gives every partition. So the product is at most $1+y/(1-2y)\le1+3y$ for $y\le1/3$.
\end{proof}

Lemma~\ref{w2:partition} is the estimate behind the value $P_g=(1+3/g_1)^f(1+3/g_2)^f$ in Hypothesis~\ref{w2:H-KE} below: the matrix $(\Gamma_{ij})$ contains the finite group factors at the typed places, and its diagonal dominance costs a $p$-adic volume factor of the form $\prod_i(1-g^{-i})^{-1}$ at each of them.

\begin{lemma}[closing numerics]\label{w2:close}
Let $b\ge16$, $0\le P_g\le e^{6b^{-7}}$, $A\le(1+b^{-7})e^{b^{-14}}$ and $0\le\varepsilon,\varepsilon'\le b^{-14}$. Then $\bigl(\sqrt{P_g(A+2\varepsilon)}+\sqrt{P_g(\varepsilon+\varepsilon')}\bigr)^2\le1+b^{-6}$.
\end{lemma}

\begin{proof}
Put $t=b^{-7}\le16^{-7}$. Then $P_g(A+2\varepsilon)\le e^{6t}(1+t)e^{t^2}+2e^{6t}t^2\le1+8t$ and $P_g(\varepsilon+\varepsilon')\le2e^{6t}t^2\le3t^2$. The square is at most $1+8t+3t^2+2\sqrt{3t^2(1+8t)}\le1+12t$, and $12t=12b^{-1}b^{-6}\le b^{-6}$.
\end{proof}

\begin{corollary}[off-diagonal and diagonal clauses from the row estimates]\label{w2:offdiag}
Let $b\ge16$. Suppose that for each pair of words the terms of $K_{uu'}$ form a term family with (a) to (d) of Theorem~\ref{w2:pair}, with $P_g\le e^{6b^{-7}}$, and with $\mathrm{part}(\mathrm{deg})=K^{\deg}_{uu'}$, $\mathrm{part}(\mathrm{id})=0$ for $u\ne u'$ and $\mathrm{part}(\mathrm{id})=1$ for $u=u'$. Suppose the row estimates (f) hold with $\varepsilon,\varepsilon'\le b^{-14}$ and
\[
dd=d^S_ud^S_{u'}/D^*,\qquad A\le(1+b^{-7})e^{b^{-14}}\ (u\ne u'),\qquad A\le\gamma_B\ (u=u'),
\]
with $\gamma_B\ge1$ and $\Gamma_0\ge6\gamma_B$. Then (iii) of Proposition~\ref{w2:assembly} holds with $\eta_0=b^{-6}$, and (iv) holds.
\end{corollary}

\begin{proof}
For $u\ne u'$ apply Theorem~\ref{w2:pair} to the terms other than the degenerate ones, whose sum is $K^0_{uu'}$, and Lemma~\ref{w2:close}. For $u=u'$ apply Theorem~\ref{w2:pair} to the terms other than the identity and degenerate ones, whose sum is $K^0_{uu}-1$. With $(\sqrt a+\sqrt b)^2\le2a+2b$ the bound is $2P_g(\gamma_B+4b^{-14})(d^S_u)^2/D^*\le6\gamma_B(d^S_u)^2/D^*\le\Gamma_0(d^S_u)^2/D^*$.
\end{proof}

\subsubsection{The inputs from the construction}\label{w2:sec:Sinputs}

Put $r_S=m/2-1$, the rank parameter of a row. The ambient rank remains $r=k+m-1$. Proposition~\ref{w2:assembly} and Corollary~\ref{w2:offdiag} reduce Theorem~\ref{w2:S} to the following statements about the concrete Gram operators.

\begin{hypothesis}[S-KE, kernel expansion]\label{w2:H-KE}
For every pair of words the unfolding of $\langle\mathcal P_uT,\mathcal P_{u'}T'\rangle$ converges absolutely and produces a term family with (a) to (d) of Theorem~\ref{w2:pair}, $P_g=(1+3/g_1)^f(1+3/g_2)^f$, $\mathrm{part}(\mathrm{id})=0$ and $\mathrm{part}(\mathrm{deg})=0$ for $u\ne u'$, $\mathrm{part}(\mathrm{id})=1$ for $u=u'$.
\end{hypothesis}

\begin{hypothesis}[S-MR, main row]\label{w2:H-MR}
$\Sigma(\mathrm{sc}(\mathrm{gen}))\le(1+\delta)e_1\,d^S_ud^S_{u'}/D^*$ for $u\ne u'$, with $0\le\delta\le b^{-7}$ and $1\le e_1\le\exp(12f2^{-(k-2r_S-2)})$; and $\Sigma(\mathrm{sc}(\mathrm{gen}))\le\gamma_B(d^S_u)^2/D^*$ for $u=u'$, with $\gamma_B=\exp(C_Bm^5f\log b)$.
\end{hypothesis}

\begin{hypothesis}[S-LR, loops]\label{w2:H-LR}
Each of the sums $\lvert\Sigma\rvert(\mathrm{big}(\mathrm{gen}\cup\mathrm{ng}))$ and $\lvert\Sigma\rvert(\mathrm{nsc}(\mathrm{gen}))$ is at most $\exp(-\tfrac13kf\log k)\,dd$.
\end{hypothesis}

\begin{hypothesis}[S-NG, nongeneric terms]\label{w2:H-NG}
$\Sigma(\mathrm{one}(\mathrm{ng}))\le2^{-(k-2r_S-5)}\exp(C_4m^5f\log b)\,dd$.
\end{hypothesis}

\begin{hypothesis}[S-DG, degenerate terms]\label{w2:H-DG}
$\lVert K^{\deg}_{uu}\rVert_\HS^2\le\exp(-kf\log k)\,d^S_u$.
\end{hypothesis}

\begin{hypothesis}[S-DIM, multiplicity formula]\label{w2:H-DIM}
$\lvert D_{\mathrm{all}}-D^*\rvert\le\exp(-c\,rf\log r)\,D^*$ for an absolute $c>0$.
\end{hypothesis}

\begin{hypothesis}[S-SZ, sizes]\label{w2:H-SZ}
$b=e^{x^\varepsilon}\ge16$, $(kb)^8\le g_1,g_2$, $f\le k^8b$, $70\log b\le k$, $k\ge8\log\Gamma_0$, $C_\Gamma\ge\max(C_B,C_4)+1$, the quantities $\exp(-\tfrac13kf\log k)$ and $2^{-(k-2r_S-5)}\exp(C_4m^5f\log b)$ are at most $b^{-14}$, $12f2^{-(k-2r_S-2)}\le b^{-14}$, and $c\,rf\log r\ge x^{\varepsilon/8}$.
We also require $32\le k$, $4r_S+12\le k$ and $8C_4m^5f\log b\le k$.
\end{hypothesis}

\begin{lemma}[sizes]\label{w39:sizes}
Let $b=e^{x^\varepsilon}$, and let $m,k,n$ be as in Proposition~\ref{ranks} with $f=2n$, $r_S=m/2-1$. For $x$ large in terms of $\varepsilon$ and $C_4$: $32\le k$, $4r_S+12\le k$, $f\le k$, $70\log b\le k$, $8C_4m^5f\log b\le k$, and consequently
\[
12f2^{-(k-2r_S-2)}\le b^{-14},\qquad2^{-(k-2r_S-5)}\exp(C_4m^5f\log b)\le b^{-14}.
\]
\end{lemma}

\begin{proof}
By Proposition~\ref{ranks}, $m\le2\sqrt x$, $f=2n\le4x^{2+8\varepsilon}$ and $k\ge x^{9/2+10\varepsilon}$. So $4r_S+12=2m+8\le k$, $f\le k$, $70x^\varepsilon\le k$, and $8C_4m^5f\log b\le1024\,C_4x^{9/2+9\varepsilon}\le k$ for large $x$. From $4r_S+12\le k$ we get $k-2r_S-2\ge k/2+4$ and $k-2r_S-5\ge k/2+1$. Then $12f2^{-(k-2r_S-2)}\le f2^{-k/2}\le e^{-k/5}$ since $f\le k\le e^{(\frac{\log2}2-\frac15)k}$ for $k\ge32$, and $2^{-(k-2r_S-5)}e^{C_4m^5f\log b}\le2^{-k/2}e^{k/8}\le e^{-k/5}$. Finally $e^{-k/5}\le b^{-14}$ because $70\log b\le k$.
\end{proof}

\begin{theorem}\label{w2:Sproof}
Hypotheses~\ref{w2:H-KE} to~\ref{w2:H-SZ} imply Theorem~\ref{w2:S} for $x$ large in terms of $\varepsilon$ and $C_0$.
\end{theorem}

\begin{proof}
By \ref{w2:H-SZ}, $P_g\le\exp(3f(g_1^{-1}+g_2^{-1}))\le\exp(6f(kb)^{-8})\le e^{6b^{-7}}$, $(1+\delta)e_1\le(1+b^{-7})e^{b^{-14}}$, and $\varepsilon=\exp(-\tfrac13kf\log k)$, $\varepsilon'$ of \ref{w2:H-NG} are at most $b^{-14}$. Also $\Gamma_0/\gamma_B=\exp((C_\Gamma-C_B)m^5fL)\ge e^{L}\ge6$ since $L=\log b$. Corollary~\ref{w2:offdiag} gives (iii) with $\eta_0=b^{-6}$ and (iv) of Proposition~\ref{w2:assembly}; \ref{w2:H-KE} gives (i); \ref{w2:H-DG} gives (ii) with $\xi=\exp(-kf\log k)$; \ref{w2:H-DIM} gives (v) with $\eta_1=\exp(-crf\log r)$. Proposition~\ref{w2:assembly} gives the three clauses with $e_0=2C_0/\Gamma_0+2\xi$. With $y=x^{\varepsilon/8}$: $\eta_0=e^{-6y^8}\le e^{-y}$; $\eta_1\le e^{-y}$ by \ref{w2:H-SZ}; $\log\Gamma_0\ge C_\Gamma y^8$ and $k\ge8\log\Gamma_0$ give $2C_0/\Gamma_0+2\xi\le e^{-2y}$ for large $y$, so $\sqrt{e_0}\le e^{-y}$.
\end{proof}

\begin{remark}[the window]\label{w2:rem:window}
Theorem~\ref{w2:S} assumes \eqref{w2:window}. Theorem~\ref{r3:A3} uses its conclusion, so the construction of the words must establish \eqref{w2:window}. The lower half of \eqref{w2:window} is the hypothesis $d^S_u\ge D^*/(C_0\Gamma_0^2)$ of Lemma~\ref{r3:sizecond}.
\end{remark}

\begin{remark}[no circularity]\label{w2:rem:circ}
Lemma~\ref{w2:punz} derives the nonvanishing used in the proof of Theorem~\ref{w2:B} from the diagonal clause of Theorem~\ref{w2:S}. Theorem~\ref{w2:S} is a statement about Poincar\'e series alone and its inputs \ref{w2:H-KE} to \ref{w2:H-SZ} do not involve the period maps, so there is no cycle. Theorem~\ref{r3:A3} uses all three theorems.
\end{remark}

\subsubsection{The error in the multiplicity formula}\label{count:error}

Let $n$ denote the rank in the counting theorem, $f$ the degree of the
totally real field, $D$ the multiplicity and $M$ the identity contribution.
With the Haar measures used for the formal degrees $d_v$, one has
$M=(\prod_{v\mid\infty}d_v)\operatorname{vol}(\mathfrak F)
\operatorname{Re}e(1)$ for the level idempotent $e$. In particular $M\ge0$.
The trace estimate also requires finite dimensionality of the multiplicity
space. In the varying-rank estimate the gap in the weights is at least
\[
  8\cdot10^{14}(C_g+4)n^{19}f^5\log(b\,\mathrm{rd}_E),
\]
where the level parameter $g$ satisfies $1\le g\le b^{C_g}$.
The hypotheses include a common signature $(p,q)$ with $q^2\le n$,
a prime $g_0>4n$ split completely in $E$, and the bounds
$\log f\le\sqrt n$, $\log\log\mathrm{rd}_E\le n$,
$\log\log b\le nf\log n/16$, and $|\lambda_{v,i}|\le b^2$.



\subsubsection{Compactly supported Gram-fibre identities}

Let $E/F$ be a quadratic CM extension with involution $\sigma$. Let $V=E^M$ be equipped with a nonsingular Hermitian form $h_J$ and unitary group $U_J$, and let $\mathcal H_n$ denote Hermitian $n\times n$ matrices. For $X=(X_1,\ldots,X_n)\in V^n$, define $\operatorname{Gram}_J(X)_{ab}=h_J(X_a,X_b)$ and $\mathcal H_n^\times=\{S:\det S\ {\rm is\ a\ unit}\}$. Fix an $F$-basis of $E$. At every place $v$ of $F$, choose additive Haar measure with integral covolume one and transport it to matrix and Hermitian coordinates as follows. The diagonal Hermitian coordinates are $\tfrac12\operatorname{Tr}_{E/F}(S_{ii})$, and the off-diagonal coordinates are the coordinates of $S_{ij}$ in the chosen basis. Write $dX_v^\circ,dS_v^\circ$ for these measures and $L_v$ for the standard local lattice of matrices.

On $\mathcal H_n^\times$, disintegrate $dX_v^\circ$ over $dS_v^\circ$ along $\operatorname{Gram}_J$. Let $\omega_{v,S}^\circ$ be the fibre measures and define
\[
 D_v^\circ(\varphi,S)=\int_{\operatorname{Gram}_J^{-1}(S)}\varphi\,d\omega_{v,S}^\circ,
 \qquad \ell_v^\circ(T)=\omega_{v,T}^\circ(L_v).
\]
Thus
\[
 \int_{\operatorname{Gram}_J^{-1}(K)}\varphi\,dX_v^\circ
 =\int_KD_v^\circ(\varphi,S)\,dS_v^\circ
\]
for every compact $K\subset\mathcal H_n^\times$. In particular, $D_v^\circ(\varphi,\cdot)$ is continuous on that locus for continuous compactly supported $\varphi$. Empty fibres have measure zero.

Let $\psi$ be a nontrivial additive character of $\mathbb A_F$. If $dX_v,dS_v$ are self-dual for the trace pairings used in the Schrödinger model, write $dX_v=c_vdX_v^\circ$ and $dS_v=e_vdS_v^\circ$. Fourier inversion gives
\[
 c_v=e_v=1\quad\hbox{for almost all }v,\qquad
 c_\infty\prod_vc_v=e_\infty\prod_ve_v=1.
\]
Uniqueness of fibre disintegration implies
\[
 D_v^{\rm sd}={c_v\over e_v}D_v^\circ
\]
for the densities taken with $dX_v,dS_v$.

\begin{proposition}\label{prop:normalized-unitary-sw}
Suppose $2n<M$ and $T\in\mathcal H_n(E)$ is nonsingular. Let $\Phi_\infty$ be continuous, rapidly decreasing, smooth and compactly supported, and let $\Phi_v$ be Schwartz--Bruhat, with $\Phi_v=1_{L_v}$ outside a finite set. Put $\Phi=\Phi_\infty\prod_v\Phi_v$. For every Haar measure $\mu$ on $U_J(\mathbb A)$ and every fundamental domain $\mathcal F$ for $U_J(F)$, one has $0<\mu(\mathcal F)<\infty$ and
\[
 \mu(\mathcal F)^{-1}\int_{\mathcal F}
 \sum_{X\in V^n:\ \operatorname{Gram}_J(X)=T}\Phi(h^{-1}X)\,d\mu(h)
 =D_\infty^{\rm sd}(\Phi_\infty,T)\prod_vD_v^{\rm sd}(\Phi_v,T).
\]
The finite product multiplies. Moreover, the fibre measures form an Euler product: $\ell_v^\circ(T)<\infty$ for every finite $v$, $\sum_v|\ell_v^\circ(T)-1|<\infty$, and, for every sufficiently large finite set $S$,
\[
 \int\Phi_\infty\prod_{v\in S}\Phi_v\,d\mathcal M_T
 =D_\infty^\circ(\Phi_\infty,T)\prod_{v\in S}D_v^\circ(\Phi_v,T)
 \prod_{v\notin S}\ell_v^\circ(T),
\]
where $\mathcal M_T$ is the normalized pushforward of $\mu|_{\mathcal F}$ along the rational fibres, and the functions are evaluated on the same orbit variables on both sides.
\end{proposition}

The proof uses only the global convergent-range Siegel--Weil formula as an external analytic theorem. All local normalizations and tails are derived below.

\subsubsection{The normalized global formula}

For the adelic Weil representation $\omega$, write
$\Theta(\Phi,g,h)=\sum_{X\in V^n}\omega(g,h)\Phi(X)$, and let $E(g,s,\Phi)$ be the associated Siegel Eisenstein series. Weil's convergent-range formula \cite[Theorem 5]{Weil1965}, in the unitary formulation also recalled in \cite[Theorem 2.3.1]{LiTheta}, asserts for $M>2n$ that
\[
 {1\over2}\int_{[U_J]}\Theta(\Phi,g,h)\,dh=E(g,(M-n)/2,\Phi),
\]
where $dh$ is Tamagawa measure. Any admissible auxiliary splitting character may be used for the dual pair; the big-cell comparison below removes that choice from the identity at the identity. No signature restriction is needed. If the Witt index $r$ is at most $n$, then $M-r>n$; otherwise $M\ge2r$ and again $M-r>n$. The full unitary Tamagawa covolume is two: $\tau(SU_J)=1$ \cite{Mars1969}, Ono's Main Theorem gives covolume two for the quadratic norm-one torus \cite{Ono1963}, and the relative theorem transports this through the determinant sequence \cite{Ono1965,WeilAdeles}. Hence division by the covolume removes the displayed factor $1/2$. Taking a nonsingular $T$-Fourier coefficient at the identity gives
\[
 \overline I_T(\Phi)=\prod_vW_{v,T}(1,\Phi_v).
\]
Borel's theorem \cite[Theorem 5.6(i)]{borel63} applies because the abelian quotient of $U_J$ is the norm-one torus, which has no nontrivial $F$-rational character. This gives finite covolume. Comparison of fundamental domains and Haar scaling give it for every $\mu,\mathcal F$ in the proposition, and covolume-normalized means are independent of those choices.

In the large Bruhat cell, the Schrödinger action and Fubini give
\begin{equation}\label{eq:whittaker-coarea}
 W_{v,T}(1,\varphi)=D_v^\circ(\varphi,T).
\end{equation}
Indeed, after integrating over the unipotent Hermitian variable, one obtains the Fourier transform at $T$ of $\operatorname{Gram}_*(\varphi\,dX_v^\circ)$; the coarea definition identifies this with $D_v^\circ(\varphi,T)$. The same equality holds with any consistent scaling of input and target measures. For $n=1$, the Hermitian scalar is identified with an element of $F$ by half the trace, and \eqref{eq:whittaker-coarea} is the rank-one Gram-fibre identity.

Varying the test functions determines an identity of locally finite Radon measures. Compact finite-place cylinders are approximated by locally constant functions; compactly supported Borel functions at infinity by smooth functions; and general nonnegative measurable functions by increasing compact truncations and monotone convergence. The scalar identity therefore includes the measurable functions in the eventual Euler-product statement.

\subsubsection{The cofinite error estimate}

It remains to bound $\ell_v^\circ(T)-1$. Begin with $n=1$. At almost every finite place, the trace form in the chosen coordinates is an integral unimodular quadratic form $Q$ in $2M$ variables, while $2$ and the fixed nonzero scalar $t$ are both units. Every residue-field solution is then smooth: otherwise nondegeneracy forces $x=0$, contradicting $Q(0)=0\ne t$. Hensel lifting identifies the local mass with the first-level count:
\[
 \ell_v^\circ(t)=q_v^{1-2M}\#\{x\in k_v^{2M}:Q(x)=t\}.
\]
The character-sum evaluation of this count gives
\[
 N_1=q_v^{2M-1}+q_v^{-1}\sum_{a\in k_v^\times}G(aQ)e(-at).
\]
The Gauss sum of a nondegenerate quadratic form in $2M$ variables has modulus $q_v^M$, so
\[
 |\ell_v^\circ(t)-1|\le q_v^{1-M}\le2q_v^{1-M}.
\]
For $M\ge3$, the error is at most $2q_v^{-2}$.

At good inert places, every unimodular Hermitian lattice has an orthonormal basis. A unit-norm vector exists: otherwise polarization, with $1/2$, would make every Hermitian pairing nonunit, contrary to unimodularity. Such a vector splits off after scaling, using surjectivity of the norm on units, and induction applies to its complement. The split case follows from a basis and its dual basis.

Now split $T\simeq(t)\perp T_1$. Two-stage coarea over the first represented vector gives
\[
 \ell_v^\circ((t)\perp T_1)
 =\delta_v(t)^n\ell_v^\circ(T_1)\ell_v^\circ(t),
\]
where $\delta_v(t)$ is the coordinate Jacobian of multiplication by $t$. At almost every place $\delta_v(t)=1$, and all congruence and Gram determinants are units. The identity
$xy-1=(x-1)(y-1)+(x-1)+(y-1)$ shows that an error $2q_v^{-M}$ for the first factor and $C_1q_v^{-2}$ for the second gives an error $(3|C_1|+2)q_v^{-2}$ for their product. Induction yields a finite set $S_0$ and a constant $C$ such that
\[
 |\ell_v^\circ(T)-1|\le Cq_v^{-2}\qquad(v\notin S_0).
\]
All lattice masses are finite by compactness. Hence $\sum_v|\ell_v^\circ(T)-1|<\infty$. Since $D_v^\circ(\Phi_v,T)=\ell_v^\circ(T)$ whenever $\Phi_v=1_{L_v}$, the finite local densities form a multipliable family for every test function in the proposition. Empty fibres cause no convergence problem.

Combining the normalized global formula, \eqref{eq:whittaker-coarea}, and this tail estimate proves the coordinate version of the proposition in every rank. Finally,
\[
 \biggl({c_\infty\over e_\infty}\biggr)\prod_v{c_v\over e_v}=1,
\]
so replacing $D^\circ$ by $D^{\rm sd}$ leaves the completed product unchanged. The scaled densities have the required continuity and compact-set identities by definition of fibre disintegration. Rank zero is elementary: all relevant spaces are singletons, all tail masses are one, both sides equal $\Phi(0)$, and a self-dual measure on a singleton has mass one.

\subsubsection{The Witt-block form}

The same normalization is compatible with induction on the target rank. Let $J'=(t)\perp J$ and $T'=(t)\perp T$. Witt extension \cite[Theorem 42:17]{OMeara1973} and measurable selection on the locally closed orbit give a section $s(y)$ from the distinguished vector $e$ to $y$ on $q(y)=t$. Define
\[
 \Psi(y)=\int_X\varphi(s(y),X)\,d\omega_T^\circ(X).
\]
Stabilizer invariance makes $\Psi$ independent of the section, and two-stage coarea gives
\[
 \int\varphi\,d\omega_{T'}^\circ
 =\delta_v(t)^n\int\Psi(y)\,d\omega_{(t),1}^\circ(y).
\]
Here $\delta_v(t)=1$ for almost all $v$ and $\delta_\infty(t)\prod_v\delta_v(t)=1$. If $\mathcal M_{J,T}$ is the averaged orbit measure and $\mathcal M^{(1)}_{J',(t)}$ is its rank-one analogue, decomposition of global frames by their distinguished vector gives
\[
 \int\varphi\,d\mathcal M_{J',T'}
 =\Big(\prod_{v\notin S}\ell_v^\circ(T)\Big)
 \int\Psi\,d\mathcal M^{(1)}_{J',(t)}.
\]
For each fixed first vector, the inner integral is the lower-rank orbit integral; stabilizer and fundamental-domain invariance identify it with the lower Euler formula. Nonnegative truncation followed by real/imaginary decomposition permits the interchange. The product formula then cancels $\delta_\infty^n\prod_{v\in S}\delta_v^n$ before passage to the limit, so the block Euler product has constant one.

Congruence of $J$ or $T$ preserves the statement: the induced group conjugacy identifies fundamental domains, while the source and target Jacobians enter in the same ratio and have adelic product one. If $q_J$ does not represent $t$ rationally, the scalar and block global measures vanish. The scalar Euler identity forces an archimedean or a finite local scalar fibre measure to vanish; the block fibre at the same place also vanishes by projecting to the first vector. The absolutely convergent zero-Euler-product statement follows. Peeling off an anisotropic scalar and applying these three operations proves the block-induction form without excluding unrepresented targets.

\subsubsection{Measure of an adapted flag class}

Let $\mathfrak A=(I,I_0,e,x_f,S,\Lambda,x_0,T,T',T_0)$ be an adapted
model. Thus $e$ is compatible with the involutions, $n=n_0+2J$ where
$J=\sum_i j_i$, each $\Lambda_i$ is a full lattice, and the blocks
satisfy
\[
 T_i^*JT_k=(T_i')^*JT_k'=0,\qquad T_i^*JT_k'=0\quad(i\ne k),
\]
together with $T_i^*JT_i'=1$, $T_i^*JT_0=(T_i')^*JT_0=0$, and
$T_0^*JT_0=J_0$. If $L,L_0$ are the two self-dual lattices and
$\Lambda_i^\vee$ is the trace dual, the adapted-lattice condition is
\[
 z\in L\quad\Longleftrightarrow\quad
 z=\sum_iT_iu_i+\sum_iT_i'u_i'+T_0e(y),\qquad
 u_i\in\Lambda_i,\quad u_i'\in\Lambda_i^\vee,\quad y\in L_0,
\]
and $S_k$ is the span of the initial $T$-blocks. Assume $j_i\ge1$.
Fix the common signature $(p,q)$ of $I$. The whole-group and subgroup
gauges are normalized by their Hilbert--Schmidt small-ball asymptotics:
in dimension $d$, the measure of the radius-$r$ ball, divided by $r^d$,
tends to the volume of the Euclidean unit ball. These limits use
$n^2,d_N,(n-2J)^2+2\sum_ij_i^2$, respectively; the compact subgroup
measure is a probability. Fix gauge constants $c=(c_v)$ and
$c_0=(c_{0,v})$ (necessarily positive). Isotropy gives $J\le\min(p,q)$.

Write $f=[F:\mathbf Q]$, $v=v_{\mathrm{inv}}$, and
$d_N=4\sum_{i<k}j_ij_k+2J(n-2J)+\sum_i j_i^2$. Let
$\mathcal H_a$ consist of those $\tau$ with
$0<\tau_1<\cdots<\tau_a<1/2$. Define
\[
 \mathcal L_f(\tau)=
 \prod_{i<k}(\tau_i/\tau_k)^{2j_ij_kf}
                (\tau_i\tau_k)^{2j_ij_kf}
 \prod_i\tau_i^{2j_i(n-2J)f}(\tau_i^2)^{j_i^2f},
 \qquad
 D(\tau)=\mathcal L_f(\tau)\prod_i\frac{2fj_i}{\tau_i}.
\]
Let $\ell_v(m)$ be the same expression with $f=1$, evaluated at the
local determinant ratios
$(d_{v,i+1}(m)/d_{v,i}(m))^{-1/j_i}$.

For a group $\Gamma$ acting on the left, put
\[
 \operatorname{CV}_\Gamma(\mu;A)=
 \inf_D\mu(D),
\]
where $D$ runs through measurable sets meeting every left
$\Gamma$-orbit in $A$. Let $\Gamma,\Gamma_0$ be the flag and residual
lattice stabilizers, and let $A_B$ be the points whose successive flag
sizes lie in $B$. Let $\mathsf V_i$ be the volume of the determinant
slab $1\le d_i<\mathrm e$ modulo $\operatorname{Aut}(\Lambda_i)$, and let
$\mathsf V_0$ be $(\prod_vc_{0,v})$ times the covering covolume of
$\Gamma_0$. Denote by $\iota(x)$ the value of a nonnegative real $x$ in
$[0,\infty]$. Finally, with $U_r$ normalized by
$U_0=1$ and $U_r/U_{r-1}=2^{r-1}(2\pi^r)/(r-1)!$, put
\[
 c_{\mathrm{Iw}}=2^{-d_N/2}\frac{U_pU_q}
 {\bigl(\prod_iU_{j_i}2^{j_i^2/2}\bigr)U_{p-J}U_{q-J}}.
\]

\begin{theorem}
For every adapted model, every $a$ and $(j_i)$ with $j_i\ge1$, every
choice of the above gauges, and every measurable
$B\subseteq\mathcal H_a$,
\begin{equation}\label{eq:adapted-class-bound}
\begin{split}
&\iota\!\left(\prod_vc_v\right)\operatorname{CV}_\Gamma(A_B)\\
&\quad\le \iota\!\left(c_{\mathrm{Iw}}^f(2v)^{fd_N/2}\right)
 \left(\prod_i\iota(2^{fj_i^2})\mathsf V_i\right)\mathsf V_0
 \int_B\iota(D(\tau))\,d\tau .
\end{split}
\end{equation}
\end{theorem}

\begin{proof}[Proof]
The index $a=0$ and the case $A_B=\varnothing$ are included, the latter
by taking the empty covering domain.

\textbf{1. The real-place measure identity.} Let $N,M,K$ be the
unipotent radical, the flag Levi subgroup, and the standard maximal
compact subgroup of the real unitary group. Orthogonal projections
adapted to the flag give $G=NMK$ and show that $um\in K$ with
$(u,m)\in N\times M$ forces $u=1$. Since the ambient group is
reductive, $N$ is nilpotent, and $K$ is compact, the closed subgroups
satisfy the hypotheses of the Haar-decomposition theorem
\cite[Theorem 8.32 and Corollary 8.31]{Knapp}. The modular function is
computed by the determinant of
the adjoint action \cite[Proposition 8.27]{Knapp}.

To evaluate it here, double the flag by duality. If $D_s$ is the
determinant on the $s$-th step and $r_s$ its complex dimension, the
adjoint determinant is
$\prod_sD_s^{r_{s+1}}D_{s+1}^{-r_s}$. Self-duality identifies opposite
steps and yields $\dim F_{2a+1-s}+\dim F_s=n$. Consequently the
pushforward of unipotent Haar measure by $u\mapsto mum^{-1}$ is
$\iota(\ell(m))$ times itself. Hence there is a finite nonzero
$\kappa$ with
\[
 \int\phi\,dg=\kappa\int_N\int_M\iota(\ell(m))
 \int_K\phi(umk)\,dk\,dm\,dn.
\]

The value of $c\kappa$ is obtained from small balls, not from a generic
measure comparison. Put $d_M=(n-2J)^2+2\sum_ij_i^2$ and
$d_L=(p-J)^2+(q-J)^2+\sum_ij_i^2$. Then
$n^2=2d_N+d_M$ and $p^2+q^2=d_N+d_L$. The volume of $K$ is $U_pU_q$,
and the volume of $K\cap M$ is
$U_L=(\prod_iU_{j_i}2^{j_i^2/2})U_{p-J}U_{q-J}$; the powers of $2$
come from the two copies of each isotropic block in the
Hilbert--Schmidt metric. Orthogonal diagonalization of the Lie-algebra
Jacobian gives a unit body of volume
$2^{d_N/2}\operatorname{ballVol}(2d_N+d_M)
 \operatorname{ballVol}(d_L)$. If $\mu_K$ is the probability measure
and $\mu_0$ is the Riemannian gauge on $K$, then
$\mu_0=(U_pU_q)\mu_K$ by Haar uniqueness. Comparing the two small-ball
asymptotics, with a slab on $K\cap M$ squeezed between radii
$r/(1+\eta)$ and $(1+\eta)r$, gives
$c\kappa U_L2^{d_N/2}=U_pU_q$. Thus $c\kappa=c_{\mathrm{Iw}}$. Taking
the finite product over the $f$ places and applying Tonelli gives
\begin{equation}\label{eq:iwasawa-product}
\begin{split}
&\iota\!\left(\prod_vc_v\right)\int\phi\,dg\\
&\quad=\iota(c_{\mathrm{Iw}}^f)\int_{N_G}\int_{M_G}
 \prod_v\iota(\ell_v(m_v))\int_{K_G}\phi(umk)\,dk\,dm\,dn .
\end{split}
\end{equation}

\textbf{2. Integral radical covolume.} Complete the flag by duality and
let $r_b$ and $h_b$ be its graded ranks and heights. Setting
$\tau_i$ equal to the size of the $i$-th graded piece at the identity,
we have $h_i=\tau_i^{2fj_i}$, $h_{2a-i}=\tau_i^{-2fj_i}$, $h_a=1$,
$r_i=r_{2a-i}=j_i$, and $r_a=n-2J$. For $0\le c<2a+1$, define
\[
 A_c=\sum_br_br_{b+c+1},\qquad
 Q_c=\prod_bh_b^{r_{b+c+1}}h_{b+c+1}^{-r_b}.
\]
The $c$-th nilpotent layer consists of homomorphism blocks. If a
biorthogonal block frame changes $\Lambda_b$ to $G_b\Lambda_b$, then
$\operatorname{cov}(\Lambda_b)\prod_\nu|\det G_{\nu,b}|^2
 =v^{fr_b}h_b$. Evaluating the source and target homomorphism
covolumes therefore gives $\operatorname{cov}(L_c)=v^{fA_c}Q_c$.

Now separate real and imaginary coordinates. Conjugation symmetry gives
$\operatorname{cov}(A)\operatorname{cov}(B)
 \le2^{fA_c}\operatorname{cov}(L_c)$. The two covolumes are equal. To
see this, for each rational prime $p$ choose a skew algebraic integer
$\delta$, purely imaginary at every selected embedding, whose real
squared determinant $N$ is prime to $p$. Such a $\delta$ is obtained by
the Chinese remainder theorem away from the finitely many maximal
ideals above $p$. Multiplication by $\delta$ sends $A$ to $B$ and $B$
to $A$. If $T$ denotes this map, integral bases yield
$|\det T|\operatorname{cov}(A)=n_1\operatorname{cov}(B)$ and
$|\det T|\operatorname{cov}(B)=n_2\operatorname{cov}(A)$ with
$n_1n_2=N^{A_c}$. If $\rho$ is the covolume ratio, then for every $p$
there are natural $a,b$ such that $\rho^2b=a$ and $p\nmid ab$. Varying
$p$ proves $\rho=1$. A real fundamental parallelepiped $D_c$ therefore
satisfies $|D_c|\le\sqrt{(2v)^{fA_c}Q_c}$.

The exponential chart has unipotent triangular coordinate changes
between successive layers. Correcting one layer at a time therefore
gives a covering domain with
$\operatorname{CV}_{\Gamma_N}\le\prod_c|D_c|$. Reindexing by ordered
pairs gives $\sum_cA_c=d_N$ and $\prod_cQ_c=\mathcal L_f(\tau)^2$.
It follows that
\begin{equation}\label{eq:radical-covolume}
 \operatorname{CV}_{\Gamma_N}
 \le\iota\!\left((2v)^{fd_N/2}\mathcal L_f(\tau)\right).
\end{equation}

\textbf{3. Levi integration.} The adapted basis gives a measurable
Levi chart $\Theta$ for which
\[
 \operatorname{size}_i(\Theta(z,g_0))
 =(d_i(z_i)/\kappa_i)^{-1/(2fj_i)},\qquad
 \kappa_i=\tau_i^{2fj_i},
\]
and $d\mu_M=(\prod_i2^{fj_i^2})(\prod_vc_{0,v})$ times the pushforward
of the product of the matrix Haar measures with $d\mu_{G_0}$. The
singular locus has measure zero, because $j_i\ge1$ makes the matrix
Haar density vanish there, so restriction to $d_i>0$ loses no mass.
Choose a measurable scale-stable fundamental domain $\mathcal D_i$ for
$\operatorname{Aut}(\Lambda_i)$, using a shortest orbit representative
and a real-coordinate lexicographic tie-break. If $N_i=2fj_i$, then
\begin{equation}\label{eq:slab}
 \int_{\mathcal D_i}H((d_i/\kappa_i)^{-1/N_i})\,dz_i
 =\mathsf V_i\int_0^\infty H(t)\frac{N_i}{t}\,dt .
\end{equation}
Indeed, translation in $\log d_i$ reduces the left side to one slab
with the weight $dt/t$, and $t=\kappa_i\tau^{-N_i}$ has
$|dt/t|=N_i\,d\tau/\tau$. Multiplication over $i$ and covering the
residual factor by a $\Gamma_0$-domain shows that the weighted Levi
covolume is at most
\[
 \left(\prod_i\iota(2^{fj_i^2})\mathsf V_i\right)\mathsf V_0
 \int_{\tau_i>0}F(\tau)\iota\!\left(\prod_i\frac{2fj_i}{\tau_i}\right)
 d\tau .
\]
The residual covolume is nonzero because $\Gamma_0$ is discrete, so
passing to the infimum over residual domains is legitimate. Put
$F=\mathbf1_B\iota(\mathcal L_f)$. On $B$, every factor is positive,
and the last integral is $\int_B\iota(D(\tau))\,d\tau$.

\textbf{4. Assembly.} If $A_B\ne\varnothing$, the size-movement formula
\[
 \operatorname{size}_i(umk)^f
 =\tau_i^f\prod_v\operatorname{locSize}_v(m)_i
\]
and membership of one point in $A_B$ prove $\tau_i>0$. Also
$\mathcal L_f(\operatorname{size}(m))
 =\mathcal L_f(\tau)\prod_v\ell_v(m_v)$. Take
$Z=\{m:\operatorname{size}(m)\in B\}$. A flag lattice stabilizer
fixes $L$ and every $S_k$, hence leaves every graded size unchanged;
in particular, $A_B$ is invariant under its action. The radical estimate
and the size relation allow the base $\mathcal L_f(\tau)$ to be taken outside
the $N$-integral and cancelled against the indicator-weighted Levi
integrand.

More explicitly, fix measurable $N$- and $M$-covering domains and a
measurable section of $K/(K\cap M)$. Form the corresponding measurable
$NMK$ slice. Its orbit meets $A_B$, and \eqref{eq:iwasawa-product}
bounds the flag covering covolume by
$c_{\mathrm{Iw}}^f\mu_N(D_N)\int_{D_M\cap Z}\prod_v\ell_v\,dm$.
Taking the infimum over $D_N$, then using
\eqref{eq:radical-covolume} and the slab computation, and finally
taking the infimum over $D_M$, gives \eqref{eq:adapted-class-bound}.
The zero and infinite $M$-integral cases are handled directly; all
other scalar--infimum interchanges involve finite nonzero constants.
Every real quantity embedded by $\iota$ is nonnegative on the indicated
domain.
\end{proof}

\subsubsection{Definite unimodular Hermitian planes}\label{sec:rank-two-planes}

\paragraph{Notation.} Let $E$ be a totally complex CM field with maximal real
subfield $F$, and let $x\mapsto\bar x$ be complex conjugation. Put
\[
\mathcal O=\mathcal O_E,\qquad f=[F:\mathbb Q],\qquad D=|\Delta_E|,
\qquad \rho=D^{1/(2f)},\qquad Y=\rho^f,\qquad L=1+f\log\rho .
\]
Thus $[E:\mathbb Q]=2f$, $\sqrt D=Y$, $\rho\geq1$, $Y\geq1$, and $L\geq1$.
For a number field $K$, $h_K$ is its ideal class number and
$N\mathfrak a=[\mathcal O_K:\mathfrak a]$ is the absolute norm of a nonzero
integral ideal. A definite unimodular plane is a full $\mathcal O$-lattice
$\Lambda$ in a two-dimensional $E$-vector space, equipped with a Hermitian
form $h$ which is positive definite at every complex embedding and satisfies
\[
\Lambda=\{x:h(y,x)\in\mathcal O\quad(y\in\Lambda)\}.
\]
Let $\operatorname{defClassNum}(E,2)$ be the number of its isometry classes.
We use Minkowski's convex-body theorem, the Minkowski ideal bound, and
Dirichlet's unit theorem in the precise forms of Marcus, \emph{Number Fields},
2nd ed., Chapter 5: the lemma on pp.~96--97, Theorem 37, Corollary~2, and
Theorem~38, respectively.

\begin{lemma}[Ideal counting]\label{lem:hyperbola-ideals}
For every number field $K$ of degree $d$ and every $X\geq1$,
\[
\#\{\mathfrak a\subseteq\mathcal O_K:\mathfrak a\neq0,\,
N\mathfrak a\leq X\}\leq X(1+\log X)^{d-1}.
\]
Consequently $h_K\leq\sqrt{|\Delta_K|}(1+\log\sqrt{|\Delta_K|})^d$; in
particular $h_E\leq YL^{2f}$.
\end{lemma}
\begin{proof}
Unique factorization of ideals and $\sum_{\mathfrak p\mid p}e_{\mathfrak p}
f_{\mathfrak p}=d$ are Marcus, Chapter~3, Theorems~16 and~21. Colour the primes
over each rational prime $p$ by $d$ colours, injectively for fixed $p$. If
$\mathfrak a=\prod_{\mathfrak p}\mathfrak p^{e_{\mathfrak p}}$, form
$A_i=\prod_{\operatorname{colour}(\mathfrak p)=i}(N\mathfrak p)^{e_{\mathfrak p}}$.
The tuple $(A_1,\ldots,A_d)$ determines $\mathfrak a$. Hence at most $d_d(m)$
ideals have norm $m$, where $d_d(m)$ counts ordered $d$-tuples of product $m$.

Set $D_1(N)=N$ and
$D_{r+1}(N)=\sum_{m=1}^{N}D_r(\lfloor N/m\rfloor)$. Induction on $r$, using
$\sum_{m\leq N}m^{-1}\leq1+\log N$, gives
$D_r(N)\leq N(1+\log X)^{r-1}$ for $N\leq X$. Since
$D_d(N)=\sum_{m\leq N}d_d(m)$, this proves the first bound. Every ideal class
has an integral ideal of norm at most
$(4/\pi)^{r_2}(d!/d^d)\sqrt{|\Delta_K|}\leq\sqrt{|\Delta_K|}$ by the
Minkowski bound, so selecting one such ideal for each class proves the second
bound. Finally $\sqrt D=Y$ and $[E:\mathbb Q]=2f$.
\end{proof}

\begin{lemma}[Gram determinant and index]\label{lem:gram-index}
Let $X\subseteq E^r$ be an integral self-dual Hermitian lattice containing
$\mathcal O^r$. If $a$ is the determinant of its standard-basis Gram matrix
and $n=[X:\mathcal O^r]$, then $|N_{E/\mathbb Q}(a)|=n^2$.
\end{lemma}
\begin{proof}
Successive coordinate projections and splitting by invertible fractional ideals
give a pseudobasis $X=\bigoplus_iC_i z_i$, where $\mathcal O\subseteq C_i$
and the rows $z_i$ are unipotent triangular in the standard coordinates. Thus
$n=\prod_i[C_i:\mathcal O]$. Let $C=\prod_iC_i$ and
$\mathfrak J=\bar C(a)C$. If $\bar s_i\in C_i$ and $t_i\in C_i$, the matrix
$(H(\bar s_i z_i,t_jz_j))_{ij}$ is integral; its determinant is
$(\prod_i s_i)a(\prod_it_i)$. Hence $\mathfrak J\subseteq\mathcal O$.
Conversely, for $u_i\in\bar C_i^{-1}$, form the standard dual vectors with
coordinates $Q^{-1}(u_ie_i)$, where $Q$ is the conjugate triangular matrix
times the Gram matrix. Self-duality puts these vectors in $X$. Comparing their
pseudobasis coefficient determinants gives $\prod_i u_i\in(a)C$. Therefore
$\bar C^{-1}\subseteq(a)C$, so $\mathfrak J=\mathcal O$. The identities
$N(C_i)[C_i:\mathcal O]=N(\bar C_i)[C_i:\mathcal O]=1$ now give
$N(\mathfrak J)n^2=|N_{E/\mathbb Q}(a)|$.
\end{proof}

\begin{lemma}[A shortest vector at every embedding]\label{lem:short-all}
Every definite unimodular plane contains a nonzero $u\in\Lambda$ such that
$h_\varphi(u,u)\leq\rho$ for every $\varphi:E\hookrightarrow\mathbb C$.
\end{lemma}
\begin{proof}
Choose an $E$-basis $v_1,v_2$ lying in $\Lambda$ and set
$\Gamma=\mathcal O v_1\oplus\mathcal O v_2$ and $n=[\Lambda:\Gamma]$.
For each complex place $w$, let $G_w$ be the Gram matrix in this basis and
$\delta_w=\det G_w$. Positive definiteness on $E$-rational vectors, density of
$E$ in $\mathbb C$, and nonvanishing of the algebraic Gram determinant show
that $G_w$ is positive definite over $\mathbb C$.

Let $a$ be the algebraic Gram determinant and $P=\prod_w\delta_w$.
Lemma~\ref{lem:gram-index} and the principal-ideal norm formula give
$P^2=|N_{E/\mathbb Q}(a)|=n^2$, hence $P\leq n$. The mixed embedding of
$\Gamma$ has covolume $(2^{-f}\sqrt D)^2=D/4^f$, so the image of $\Lambda$
in $\prod_w\mathbb C^2$ has covolume $D/(4^fn)$. For $R^2=\rho$, define
\[
\mathcal E_R=\{(z_w)_w:z_w^*G_wz_w\leq R^2\quad(w)\}.
\]
The volume of a complex Hermitian ellipsoid follows from its Cholesky
factorization:
$\operatorname{vol}\{z:z^*G_wz\leq R^2\}=\pi^2R^4/(2\delta_w)$.
Since $R^{4f}=\rho^{2f}=D$,
\[
\operatorname{vol}(\mathcal E_R)=\frac{(\pi^2/2)^fD}{P},
\qquad 2^{4f}\operatorname{covol}(\Lambda)=\frac{4^fD}{n}.
\]
Minkowski's theorem therefore gives
\[
\#(\mathcal E_R\cap\Lambda)\geq
\frac{\operatorname{vol}(\mathcal E_R)}{2^{4f}\operatorname{covol}(\Lambda)}
=\frac nP\Bigl(\frac{\pi^2}{8}\Bigr)^f>1.
\]
This supplies the required nonzero vector at the place representatives.
Conjugate embeddings have the same real Hermitian value, proving the bound
at every embedding.
\end{proof}

\paragraph{Diagonal splitting.} Choose $u$ as in
Lemma~\ref{lem:short-all}, and complete it to an orthogonal basis $(u,w)$.
Writing $\gamma=(u,w)$ gives
$\gamma^*J\gamma=\operatorname{diag}(\beta_1,\beta_2)$, where
$\beta_1=h(u,u)$, $\beta_2=h(w,w)$ are totally positive. Put
$L=\gamma^{-1}\Lambda$,
\[
A=\{a:(a,0)\in L\},\qquad B=\{b:(0,b)\in L\},
\]
and let $P,Q$ be the two coordinate projections of $L$. Then $1\in A$ and
\[
t\in A\iff\bar p\beta_1t\in\mathcal O\quad(p\in P),\qquad
t\in B\iff\bar q\beta_2t\in\mathcal O\quad(q\in Q).
\tag{1}
\]
Both projections are nonzero because $L$ is full. Using inverse fractional
ideals in (1) shows that $A$ and $B$ are nonzero full rank-one lattices, just
as in the dual-colon calculation below. Define
\[
\mathfrak a_1=\operatorname{span}_{\mathcal O}
\{\bar x\beta_1y:x,y\in A\},\qquad
\mathfrak a_2=\operatorname{span}_{\mathcal O}
\{\bar x\beta_2y:x,y\in B\}.
\]
For $c\in E$, the identity
\[
\bar b\beta_1((\bar x\beta_1y)a)
=(\bar b\beta_1y)\overline{(\bar a\beta_1x)}
\]
and inversion of the fractional ideal $P$ give
$c\in\mathfrak a_1\iff cP\subseteq A$. Similarly
$c\in\mathfrak a_2\iff cQ\subseteq B$. If $cP\subseteq A$ and $(p,q)\in L$,
subtracting $(cp,0)$ from $c(p,q)$ proves $(0,cq)\in L$. Hence
$\mathfrak n:=\mathfrak a_1=\mathfrak a_2$.

The ideal $\mathfrak n$ is generated by its conjugation-fixed elements. Indeed,
let $\mathfrak b$ be the ideal generated by those elements. If
$(\mathfrak b:\mathfrak n)\neq\mathcal O$, choose a maximal $\mathfrak p$
containing it. By (1) and $P^{-1}P=\mathcal O$, there are $x,y\in A$ and
$a,b\in P$ with $\bar x\beta_1a$ and $\bar b\beta_1y$ outside $\mathfrak p$.
Examining the other two cross products shows that one of $x$, $y$, or $x+y$
can be used so that both
$s=\bar x\beta_1a$ and $t=\bar b\beta_1x$ are outside $\mathfrak p$. For every
generator $\bar y\beta_1z$ of $\mathfrak n$,
\[
st(\bar y\beta_1z)=(\bar y\beta_1a)(\bar b\beta_1z)(\bar x\beta_1x)
\]
belongs to $\mathfrak b$. Thus $st\in\mathfrak p$, a contradiction.
Consequently $\mathfrak n=(\mathfrak n\cap\mathcal O_F)\mathcal O$.
Finally $1\in A$ gives $\beta_1\in\mathfrak n$, and the shortest-vector
estimate at all $2f$ embeddings gives
\[
N\mathfrak n\leq|N_{E/\mathbb Q}(\beta_1)|
=\prod_\varphi|h_\varphi(u,u)|\leq\rho^{2f}=Y^2.
\tag{2}
\]

\begin{lemma}[Real ideals and lines]\label{lem:real-lines}
For $Y\geq1$, there are at most $YL^f$ nonzero fixed-generated ideals with
$N\mathfrak n\leq Y^2$. For each such $\mathfrak n$, there are at most
$h_E2^f$ equivalence classes of pairs $X=(\beta,N)$ with $\beta$ totally
positive, $N$ a full $\mathcal O$-lattice in $E$, and
$\operatorname{span}_{\mathcal O}\{\bar x\beta y:x,y\in N\}
=\mathfrak n$. Here $(\beta,N)\sim(\beta',N')$ if
$\bar c\beta'c=\beta$ and $cN=N'$ for some $c\in E^\times$.
\end{lemma}
\begin{proof}
Contraction $\mathfrak n\mapsto\mathfrak a=\mathfrak n\cap\mathcal O_F$ is
injective by fixed-generation, and
$N_{E/\mathbb Q}(\mathfrak a\mathcal O)=(N_{F/\mathbb Q}\mathfrak a)^2$.
The first claim follows from Lemma~\ref{lem:hyperbola-ideals} applied to $F$.

For lines, choose a reference $(\beta^0,N^0)$ in each occurring fractional
ideal class. For $X$ in that class choose $aN=N^0$. Equality of norm ideals
gives $u_X\in\mathcal O^\times$ with $\bar a\beta^0a\,u_X=\beta$: the multiplier
$t=\bar a\beta^0a/\beta$ and its inverse stabilize the nonzero finitely
generated ideal $\mathfrak n$, hence are integral. Record the parity of the
exponents of $u_X$ in a fundamental unit basis. Dirichlet's theorem gives
$f-1$ free generators, hence at most $2^{f-1}$ parity vectors.

If $X,Y$ have the same ideal class and parity, then $u_X=u_Y\zeta v^2$ with
$\zeta$ torsion and $v$ a unit. Put $c=a_Xv/a_Y$. This gives $cN_X=N_Y$ and
\[
\bar c\beta_Yc=\bar a_X\beta^0a_Xu_Y\bar v v.
\]
At any embedding, $|\zeta|=1$ and $|v^2|=|\bar v v|$, so this positive fixed
value equals $\beta_X=\bar a_X\beta^0a_Xu_X$. Thus $X\sim Y$. There are at
most $h_E2^{f-1}\leq h_E2^f$ line classes.
\end{proof}

\begin{lemma}[Unimodular gluings]\label{lem:glue-count}
Fix line representatives $(\beta_1,N_1)$ and $(\beta_2,N_2)$, with the norm
ideal of the first equal to $\mathfrak n$. The lattices $M\subseteq E^2$ which
are self-dual for $\operatorname{diag}(\beta_1,\beta_2)$ and have prescribed
axis intersections $N_1,N_2$ number at most $N\mathfrak n$.
\end{lemma}
\begin{proof}
Write $N^{\vee_i}=\{a:\bar s\beta_i a\in\mathcal O\ (s\in N)\}$. Choose
$0\neq s\in N_1$ and $d=\bar s\beta_1$. For the integral ideals
$I=\{c:d^{-1}c\in N_1\}$ and $J=\{c:d^{-1}c\in N_1^{\vee_1}\}$,
prime factorization gives $y$ with $J=I+(y)$. Taking $u=d^{-1}y$ gives
$N_1^{\vee_1}=N_1+\mathcal O u$. Also
$\mathfrak nN_1^{\vee_1}\subseteq N_1$, since
$(\bar x\beta_1y)u'=(\bar x\beta_1u')y$.

If $P$ is the first projection of an admissible $M$ and
$a\in N_1^{\vee_1}$, then $a\in P$. Indeed, for $q\in P^{-1}$,
$\bar q\beta_1^{-1}\in N_1$; hence $qa$ is integral and
$a(pq)=(qa)p\in P$. Summing over $PP^{-1}=\mathcal O$ proves $a\in P$.
Thus there is $b$ with $(a,b)\in M$, and similarly in the second coordinate.

Fix $M_0$ if the family is nonempty, and choose $(u,v_M)\in M$ for every $M$.
For each $M$, lift $v_M$ to $(a,v_M)\in M_0$ and choose $c_M\in\mathcal O$
with $a-c_Mu\in N_1$. Then $v_M-c_Mv_0\in N_2$. If
$c_M-c_{M'}\in\mathfrak n$, it follows that $v_M-v_{M'}\in N_2$. For any
$(a,b)\in M$, choose $d$ with $a-du\in N_1$; then $b-dv_M\in N_2$ and
\[
(a,b)=(a-du,0)+d(u,v_{M'})
 +(0,b-dv_M+d(v_M-v_{M'}))\in M'.
\]
Symmetry gives $M=M'$. Therefore $M\mapsto c_M\bmod\mathfrak n$ is an
injection into $\mathcal O/\mathfrak n$.
\end{proof}

\begin{theorem}\label{thm:rank-two-class-bound}
With the notation above,
\[
\operatorname{defClassNum}(E,2)\leq4^f\rho^{5f}(1+f\log\rho)^{5f}.
\]
\end{theorem}
\begin{proof}
By Lemma~\ref{lem:real-lines} and $h_E\leq YL^{2f}$, there are at most $YL^f$
ideals and at most $2^fYL^{2f}$ classes for each line. By
Lemma~\ref{lem:glue-count} and (2), each pair of line representatives has at
most $\lfloor Y^2\rfloor$ gluings. Thus the number of tuples
$(\mathfrak n,[X_1],[X_2],M)$ is at most
\[
(YL^f)(h_E2^f)^2\lfloor Y^2\rfloor
\leq4^fY^5L^{5f}.
\]
Every class has such a tuple. Namely, after splitting, choose for each
resulting line $X_i=(\beta_i,N_i)$ its reference representative
$X_i^0=(\beta_i^0,N_i^0)$, and scalars with
$\bar c_i\beta_i c_i=\beta_i^0$ and $c_iN_i^0=N_i$.
With $\delta=\gamma\operatorname{diag}(c_1,c_2)$ and $M=\delta^{-1}\Lambda$,
the diagonal form is $\operatorname{diag}(\beta_1^0,\beta_2^0)$ and $M$ has
exactly the required axis intersections. Fixed tuple data determine this
diagonal pair, and the original plane is isometric to it; hence the map from
classes to tuples is injective. Since $Y=\rho^f$ and $L=1+f\log\rho$, this is
the claimed bound.
\end{proof}

\paragraph{Absolute-volume substitution.} Let $\ell_E(M)$ denote the local
$L$-factor in the rank recurrence and set
$\operatorname{rankFac}(E,M)=\rho^{f(M-1/2)}\ell_E(M)$. The local comparison
$\ell_E(M)\leq\zeta(M)^f$ and the explicit estimate
$\prod_{M=3}^{n}\zeta(M)<1.4$ give
\[
\prod_{2<M\leq n}\operatorname{rankFac}(E,M)
\leq\rho^{f(n^2-4)/2}1.4^f,
\qquad\sum_{2<M\leq n}(M-\tfrac12)=\frac{n^2-4}{2}.
\tag{3}
\]
For the zeta product, truncate the $\zeta(3)$ and $\zeta(5)$ series and use
$\zeta(4)=\pi^4/90$, $\zeta(6)=\pi^6/945$; for $M\geq6$, comparison term by
term gives $\zeta(M)-1\leq2^{6-M}(\zeta(6)-1)$, hence
$\sum_{M\geq6}(\zeta(M)-1)<0.0347$. Together with $\zeta(3)<1.2021$,
$\zeta(4)<1.0824$, $\zeta(5)<1.0385$, and $e^{0.0347}<1.0354$, this gives the
product bound $1.3991<1.4$.

Assume now that the rank recurrence has reduced the volume $V_{p,q}$ to the
definite signature $(2,0)$ and that $n=p+q\geq8$, $q^2\leq n$, $X\geq2$,
$\rho\leq X$, and $2f\rho\leq X^2$. (The first two imply $p\geq2$.) Since $\operatorname{volU}(2)=8\pi^3$,
Theorem~\ref{thm:rank-two-class-bound} and (3) yield
\[
V_{p,q}\leq(8\pi^3)^f\operatorname{defClassNum}(E,2)
\prod_{2<M\leq n}\operatorname{rankFac}(E,M)
\leq(8\pi^3\cdot4\cdot1.4)^f\rho^{5f}L^{5f}\rho^{f(n^2-4)/2}.
\]
The remaining comparisons are explicit:
\[
8\pi^3\cdot4\cdot1.4\leq8(3.15)^3(4)(1.4)<2^{11}\leq X^{11},
\]
and, using $\log\rho\leq\rho-1$ and $f\geq1$,
\[
L=1+f\log\rho\leq1+f(\rho-1)\leq f\rho\leq X^2/2.
\]
Consequently
\[
V_{p,q}\leq X^{11f}X^{5f}X^{10f}X^{f(n^2-4)/2}
=X^{f(26+(n^2-4)/2)}\leq X^{n^2f},
\]
where the last inequality uses $n^2\geq64$, hence $48\leq n^2$.

\subsubsection{Decay and the far orbits}

Write $r=\min(p,q)\ge1$. Let $h$ be a lower bound for the gaps between opposite-sign weights and suppose $h\ge120nr^{3/2}$. In polar coordinates $k_1e^Xk_2$ put $t=\|X\|_{\rm HS}$. For unit vectors in the lowest $K$-type the absolute value of the matrix coefficient is at most
\[
 \exp\{-\tfrac h8\min(t^2,t/\sqrt{2r})\}.
\]
Indeed the compression to the lowest $K$-type is the compact integral of
$\prod_{j<n}\Delta_j(k^{-1}e^Xk)^{-(\lambda_j-\lambda_{j+1}+1)}$
against the two highest-vector matrix coefficients, with factor $\dim S$ for the lowest type $S$. The normalization follows by Schur orthogonality at the identity. The integral and the discrete-series coefficient satisfy the same lowering equations: commuting a compact derivative past a noncompact one gives the Lie bracket, and an ordered lowering derivative kills the highest vector. The gap $h>2n$ excludes the complementary types in the successive jets. Induction on the noncompact degree gives equality of the Taylor expansions, hence of the analytic coefficients along $\exp(tX)$.

For $Y=k^{-1}Xk$ the leading minor $\Delta_j(e^Y)$ is a convex sum of exponentials of sums of $j$ eigenvalues, by Cauchy--Binet. These sums have mean zero, absolute value at most $\sqrt{2r}\,t$, and variance $B_j=\sum_{a>j,\,b\le j}|Y_{ab}|^2$. The centered exponential estimate at $\min(1,1/(\sqrt{2r}\,t))$ and convexity give
$\log\Delta_j(e^Y)\ge B_j\min(1,1/(\sqrt{2r}\,t))/4$.
Summation over the cuts gives
$\sum_j(\lambda_j-\lambda_{j+1}+1)B_j\ge ht^2/2$.
Cauchy--Schwarz and Schur orthogonality in the compact integral give the asserted bound. In the operator radius $s=\|X\|_{\rm op}$ the paired nonzero eigenvalues give $t\ge\sqrt2s$, hence the exponent is at most $-(h/8)\min(2s^2,s/\sqrt r)$.

Multiply the local estimates by the product $d_\infty$ of the formal degrees and the idempotent bound $a=e(1)$. For least lattice scale $\tau$ put $A(x)=(3\sqrt n/\min(1/2,\tau)^2)^{2n^2f}$. The number of elements with radii at most $s_v$ is bounded by $A(x)\exp(2n^2\sum_v s_v)$, by packing lattice points in the product of balls. Put $\kappa=2n^2$, $T=h\varepsilon^2/16$, and take $\varepsilon\le1/\sqrt{2r}$, $\kappa\varepsilon\le\log2$, $32\kappa\le h\varepsilon$, $T\ge1$. In a dyadic shell with upper radius $2^j\varepsilon$, $j\ge1$, decay times the counting factor is at most $\exp(-T2^{j-1})$; the bottom shell contributes at most $2$. The sum of the positive-index terms is at most $2e^{-T}\le1$. A far orbit has at least one positive index, so the product over $f$ places gives an absolutely convergent sum bounded by $d_\infty a A(x)\,2f3^fe^{-T}$.

For the deep region put $N=n^2f$, $P=8qN$, and write $\tau$ for the least lattice scale, $\rho=n^2\sqrt f\,g\,\mathrm{rd}_E(2\tau)^{1/q}/4$. Let $\beta$ bound the one-step derivative constants through order $P+N$. The local Sobolev estimate on jets has squared evaluation bound
\[
 (6\beta)^N e^f(17\sqrt n/\tau^2)^{2N}
 D\,\prod_v c_v/\mu_f(K_g).
\]
Here $c_v$ are the local gauge constants. The lattice scale on the radius-$\rho$ neighbourhood is at least $\tau/2$. At each of the $P$ steps the constant term along the isotropic radical vanishes by cuspidality; comparison with its average costs at most $2\rho\sqrt f\,\beta$. Squaring gives the additional factor $(2\rho\sqrt f\,\beta)^{2P}$. Applying the frame estimate first to finite subspaces and integrating with the absolute volume bound gives finite dimensionality. An orthonormal basis then gives the diagonal kernel as the squared evaluation norm.

\subsubsection{Near orbits and exceptional flags}\label{fr:sec:near-orbits}

Throughout this subsection let $E/F$ be the totally imaginary quadratic extension of the totally real field in the automorphic datum, and put $f=[F:\mathbb Q]$. Thus $E$ has $f$ complex places. Let $n$ be the hermitian rank and write $J=\sum_i j_i$ for a type $(j,\tau)$ of length $a$. Define
\[
 \mathcal P_\delta(j,\tau)=
 \prod_{i<k}\!\left[G_\delta\!\left(\frac{\tau_i}{\tau_k}\right)^{2j_ij_kf}
 G_\delta(\tau_i\tau_k)^{2j_ij_kf}\right]
 \prod_i\!\left[G_\delta(\tau_i)^{2j_i(n-2J)f}
 G_\delta(\tau_i^2)^{j_i^2f}\right],
 \qquad G_\delta(u)=1+\frac{12\delta}{u}.
\]
Write $\mathcal F$ for the fundamental domain with measure $\mu$, and $\varphi$ for the product of the archimedean matrix coefficients, their formal degrees and the finite level idempotent. Write $s_{\min}(x)$ for the least lattice scale. For the lattice at $x$, let $S_{\le}(t)$ and $S_{<}(t)$ be the spans of sublattices of largest size at most $t$ and strictly less than $t$, respectively. Let $\mathcal N(x,\delta)$ consist of the rational elements whose conjugates by $x$ lie in the finite level and within Frobenius distance $\delta$ of the identity at each infinite place. Write $\mathcal N^*(x)$ for the nonidentity near orbits at the original Cartan cutoff.

\begin{proposition}[Near unipotence]\label{fr:prop:near-unipotent}
Assume $n\ge1$, $4n<g_0$, and that $(g_0)\mathcal O_E$ is radical. For $0<c\le10^{-5}$ put
\[
 \varepsilon=\frac{c}{n^8f^2},\qquad \delta=3\sqrt{n\varepsilon}.
\]
Then $\delta<1/4$, and $\delta qnf\le10^{-2}$ whenever $q^2\le n$. Moreover every automorphic $\gamma$ for which $\varphi(x^{-1}\gamma x)\ne0$ and whose archimedean Cartan distances after conjugation by $x$ are at most $\varepsilon$ has
\[
 (\gamma_f-1)^n=0,\qquad (\gamma_v-1)^n=0\quad(v\mid\infty),
\]
and belongs to $\mathcal N(x,\delta)$.
\end{proposition}

\begin{proof}
The scalar assertions are uniform in $n,f\ge1$. With $a=n\varepsilon$ one has
\[
 a(n^3f^2)=c/n^4\le c,\qquad a\le c\le10^{-5}.
\]
Hence $3\sqrt a<1/4$. If $q^2\le n$, then
\[
 (3\sqrt a\,qnf)^2=9a(q^2n^2f^2)\le9a(n^3f^2)\le9c<10^{-4},
\]
which gives the last radius bound. Also $n\varepsilon\le1/4$.

For an invertible complex matrix $A$, write $A=ke^X$ with $k$ unitary and $X$ hermitian. This is the unitary polar decomposition: $A^*A=e^{2X}$ and $k=Ae^{-X}$. The Cartan distance in the statement is $\|X\|_F$. Suppose $\|X\|_F\le\varepsilon$ and put $E=e^X$. Then
\[
 \|E-1\|_F\le e^\varepsilon-1,\qquad
 \|(kE)^r\|_F\le e^{r\varepsilon}\sqrt n,
\]
by expanding the exponential and inducting on $r$. Cauchy--Schwarz on diagonal entries gives $|\operatorname{tr}M|\le\sqrt n\,\|M\|_F$. Since $r\varepsilon\le n\varepsilon\le1/4$ and $e^{1/4}<3/2$, it follows for $0\le r\le n$ that
\begin{equation}\label{fr:eq:archtrace}
 |\operatorname{tr}((kE)^r)-n|\le ne^{r\varepsilon}+n\le3n.
\end{equation}
Now suppose in addition $\operatorname{tr}(kE)=n$. Using $\varepsilon\le1$ gives $\|E-1\|_F\le2\varepsilon$. The identities
\[
 \|k-1\|_F^2=2n-2\Re\operatorname{tr}k,
 \qquad \operatorname{tr}k=n-\operatorname{tr}(k(E-1))
\]
imply $\|k-1\|_F\le2\sqrt{n\varepsilon}$. Since $\sqrt\varepsilon\le\sqrt n$ and $\sqrt{n\varepsilon}\le1/2$,
\[
 \|kE-1\|_F\le2\sqrt{n\varepsilon}+2\varepsilon\le3\sqrt{n\varepsilon}=\delta.
\tag{1}
\]

Nonvanishing of the level idempotent at the conjugated finite component puts that component in the lattice stabilizer and makes it unipotent modulo $g_0$, by the two support clauses of admissibility. Let $G$ be the rational matrix supplied by rationality of the automorphic orbit, and let $\Lambda_x$ be its self-dual $\mathcal O_E$-lattice. Thus $G\Lambda_x\subseteq\Lambda_x$ and
\[
 (G-1)^n\Lambda_x\subseteq g_0\Lambda_x.
\]
For fixed $r$, set $S_r=\sum_{s=0}^{r-1}G^s$. Since $(G-1)S_r=G^r-1$ and $S_r$ commutes with $G-1$, the same divisibility holds for $(G^r-1)^n$.

We use the following trace congruence. If $A\Lambda\subseteq\Lambda$ and $A^m\Lambda\subseteq g_0\Lambda$, then $\operatorname{tr}A\in g_0\mathcal O_E$. Indeed, for an eigenvalue $z$ of $A$ in an algebraic closure, a nonzero left eigenvector and a common denominator in $\Lambda$ show that $z$ is integral over $\mathcal O_E$, while $z^m/g_0$ is integral. Thus $z$ lies in the radical of $g_0$ in the integral closure. Summing eigenvalues gives $(\operatorname{tr}A)^N\in g_0\mathcal O_E$ for some $N$; radicality of $(g_0)$ gives the assertion. Consequently
\[
 \operatorname{tr}(G^r)-n=g_0t_r\qquad(t_r\in\mathcal O_E).
\]

The archimedean realizations of $G^r$ are conjugate to the matrices in \eqref{fr:eq:archtrace}, so $|\sigma(g_0t_r)|\le3n$ at every complex embedding $\sigma$. If $t_r\ne0$, the product formula for its nonzero algebraic norm gives some $|\sigma(t_r)|\ge1$; hence $g_0\le3n$, contrary to $4n<g_0$. Thus $\operatorname{tr}(G^r)=n$ for $0\le r\le n$.

For $N=G-1$, the binomial expansion now yields $\operatorname{tr}(N^k)=n(1-1)^k=0$ for $1\le k\le n$. Newton's identities over the characteristic-zero field $E$ force the characteristic polynomial of $N$ to be $X^n$. Cayley--Hamilton gives $N^n=0$. Ring homomorphisms and conjugation preserve this identity, giving unipotence at every finite and infinite component. In particular the trace at $r=1$ is $n$, so (1) supplies exactly the archimedean radius in $\mathcal N(x,\delta)$. The other two membership clauses are rationality and the lattice stabilizer already established.
\end{proof}

At the operative hypotheses, $g_0$ is prime and the number of primes of $\mathcal O_E$ over $g_0$ equals $[E:\mathbb Q]$. The identity $\sum_{\mathfrak p\mid g_0}e_{\mathfrak p}f_{\mathfrak p}=[E:\mathbb Q]$ therefore forces every ramification index and residue degree to be $1$. In particular $(g_0)$ is squarefree, hence radical. This is the total-splitting input used above.

\begin{lemma}[Geometric-mean packing]\label{fr:lem:gm-packing}
Let $U_v$, $v=1,\dots,f$, be real vector subspaces of a finite-dimensional normed space, with $\dim_{\mathbb R}U_v\le N$. Suppose $y_t(v)\in U_v$ has norm at most $X>0$, and distinct $t,t'$ satisfy
\[
 \left(\prod_v\|y_t(v)-y_{t'}(v)\|\right)^{1/f}\ge\lambda>0.
\]
Then the number of indices $t$ is at most $(1+2X/\lambda)^{Nf}$.
\end{lemma}

\begin{proof}
Put $G=(\prod_vX)^{1/f}=X$ and, in the general version with radii $X_v$, rescale coordinate $v$ by $G/X_v$. These scalars have product $1$, so the geometric mean of every difference is unchanged. After rescaling all coordinates have norm at most $G$, while the geometric mean of a difference is at most its sup norm. Thus the points are $\lambda$-separated in sup norm in a ball of radius $G$. Their balls of radius $\lambda/2$ are disjoint and lie in the ball of radius $G+\lambda/2$. Haar-volume comparison in dimension at most $Nf$ proves the bound; padding each $U_v$ to dimension $N$ removes the word ``at most''.
\end{proof}

We next isolate the slope estimate behind the block count. For an $\mathcal O_E$-lattice $L$ of rank $r$, write
\[
 H(L)=2^{fr}|d_E|^{-r/2}\operatorname{covol}(L),\qquad
 \operatorname{size}(L)=H(L)^{1/(2fr)},
\]
and use the analogous normalized covolume $H(A/A_0)$ and $\operatorname{size}(A/A_0)$ for a saturated $A_0\subsetneq A$. Let $\operatorname{leastSize}(B)$ be the infimum of sizes of nonzero sublattices of $B$, and $\operatorname{largestSize}(A)$ the supremum of quotient sizes over proper saturated $A_0$.

If $\phi:A\to B$ has nonzero image and local operator norms $c_v$, use the singular-value decomposition of the induced map from $A/(A\cap\ker\phi)$ to its image at each complex place. Its $r=\operatorname{rk}\phi(A)$ singular values are at most $c_v$, so the determinant of that induced map has squared absolute value at most $c_v^{2r}$. Applied to a free full-rank sublattice of the quotient, this gives, with $P=\prod_vc_v$,
\[
 H(\phi A)\le P^{2r}H(A/(A\cap\ker\phi)).
\]
Taking the $(2fr)^{-1}$-st root gives
\[
 \operatorname{size}(\phi A)\le P^{1/f}\operatorname{size}(A/(A\cap\ker\phi)).
\]
The image itself is a test lattice for $\operatorname{leastSize}(B)$, and $\ker\phi\cap A$ is a test quotient for $\operatorname{largestSize}(A)$. Therefore
\begin{equation}\label{fr:eq:size-inequality}
 \frac{\operatorname{leastSize}(B)}{\operatorname{largestSize}(A)}
 \le\left(\prod_vc_v\right)^{1/f}.
\end{equation}
Quotienting by a saturated sublattice and replacing operator norms by distance-to-subspace bounds gives the same estimate; this is the form used below.

\begin{proposition}[Filtration-aware near count]\label{fr:prop:near-count}
For every point $x$ of type $(j,\tau)$ and every $0<\delta<1/4$,
\[
 \mathcal N(x,\delta)\text{ is finite},\qquad
 \#\mathcal N(x,\delta)\le\mathcal P_\delta(j,\tau).
\]
No additional numerical hypothesis on the type is needed.
\end{proposition}

\begin{proof}
The type supplies a saturated flag $0=S_0\subsetneq\cdots\subsetneq S_a$ with rank increments $j_i$ and
\[
 S_{i+1}=S_{\le}(\tau_i),\qquad S_i=S_{<}(\tau_i),\qquad S_a=S_{<}(1/2).
\]
Complete it by hermitian duality: $G_p=S_p$ for $p\le a$ and $G_p=G_{2a+1-p}^{\perp}$ for $p>a$. Put $t_i=\tau_i$ for $i<a$, $t_a=2$, and $t_p=\tau_{2a-p}^{-1}$ for $p>a$. Self-duality gives $G_{2a+1}=\Lambda_x$, $\operatorname{rk}G_p+\operatorname{rk}G_{2a+1-p}=n$, and the strict-step descriptions of every $G_p$. If $r_p=\operatorname{rk}G_{p+1}-\operatorname{rk}G_p$, then
\[
 r_i=r_{2a-i}=j_i\ (i<a),\qquad r_a=n-2J.
\]

The condition $\delta<1/4$ is used at the central step. A near matrix lowers $S_{\le}(t)$ to $S_{\le}(\delta t)$. If it did not, applying \eqref{fr:eq:size-inequality} to the induced nonzero map on the quotient would give a ratio at most $\delta$, whereas the strict-step lower size and the largest-size upper bound make that ratio greater than $\delta$. At $p=a$, $2\delta<1/2$ gives $S_{\le}(2\delta)\subseteq S_{<}(1/2)=G_a$. Consequently a rational near matrix lies in
\[
 \mathcal L=\{\gamma:\gamma\text{ isometric},\ ( \gamma-1)G_{p+1}\subseteq G_p\ (p\le2a),\
 \|\gamma_v-1\|\le\delta\ \forall v\}.
\]
The rational representation is injective: its archimedean and finite matrices determine the adelic element. It therefore remains to count $\mathcal L$.

Index blocks by $\mathcal B=\{(q,p):0\le q<a,\ q<p,\ q+p\le2a\}$ and order them by decreasing height $p-q$. The relation that $D=\gamma-\gamma'$ maps $G_{p+1}$ to $G_q$ is an equivalence. Unitarity gives the exact identity
\[
 h(Du,\gamma v)+h(\gamma'u,Dv)=0.
\]
A strong induction on $p$, with a second induction on the dual index, propagates the relation from smaller heights to all pairs $(q,p)$ of height $p-q<c$ and proves $D(G_p)\subseteq G_{p-c}$. It also proves
\[
 h(Du,v)+h(u,Dv)=0\qquad
 (u\in G_{p+1},\ v\in G_{p'+1},\ p+p'\le2a+c).
\]
At height $2a+1$, this makes $D$ zero on $G_{2a+1}=\Lambda_x$, hence zero because $\Lambda_x$ spans.

For a fixed block $(q,p)$, subtract a basepoint $\gamma_0$. The corresponding archimedean block is the projection of $(\gamma-\gamma_0)_v$ from the graded piece of rank $r_p$ to the dual graded piece of rank $r_{2a-q}$. It has real dimension $2r_pr_{2a-q}$, except when $q+p=2a$; in that diagonal case the skew identity restricts it to a skew-hermitian space of real dimension $r_p^2$. These dimensions are exactly
\[
 d_{q,p}=\begin{cases}r_p^2,&q+p=2a,\\2r_pr_{2a-q},&q+p<2a.\end{cases}
\]
All block norms are at most $2\delta$.

If two elements differ, the relation at level $(q,p)$ fails for a quotient-separating subset unless they agree. By the previous paragraph, nontriviality of the induced map supplies a nonzero quotient image. The quotient form of \eqref{fr:eq:size-inequality}, the strict-step lower bound $t_q$ and the largest-step upper bound $t_p$ give geometric-mean separation at least $t_q/t_p$. Lemma~\ref{fr:lem:gm-packing} therefore bounds this fibre by
\[
 \left(1+\frac{2(2\delta)}{t_q/t_p}\right)^{d_{q,p}f}.
\]
Taking fibre products in decreasing height proves finiteness and
\begin{equation}\label{fr:eq:raw-block-product}
 \#\mathcal L\le\prod_{(q,p)\in\mathcal B}
 \left(1+\frac{4\delta}{t_q/t_p}\right)^{d_{q,p}f}.
\end{equation}

It remains to evaluate the product. For $i<k<a$ the two pairs $(i,k)$ and $(i,2a-k)$ have dimensions $2j_ij_k$ and ratios $\tau_i/\tau_k$ and $\tau_i\tau_k$, respectively. The central pairs $(i,a)$ and $(i,2a-i)$ have dimensions $2j_i(n-2J)$ and $j_i^2$ and ratios $\tau_i/2$ and $\tau_i^2$. All these identifications use $t_a=2$. Finally, for positive $\ell,\ell'$ with $\ell\le3\ell'$,
\[
 \left(1+\frac{4\delta}{\ell'}\right)^m
 \le\left(1+\frac{12\delta}{\ell}\right)^m,
\]
because $4\delta/\ell'\le12\delta/\ell$. It applies with equality of ratios away from $(i,a)$, and with $\ell=\tau_i$, $\ell'=\tau_i/2$ at $(i,a)$. Substitution in \eqref{fr:eq:raw-block-product} gives exactly $\mathcal P_\delta(j,\tau)$, including every exponent and both factors for each off-diagonal pair.
\end{proof}

For the integrated estimate write $L=nf\log n$, $X=nfrd\,g\,b$, $\ell_*=X^{-11q}$, and
\[
 Y=2fq(11q\log X).
\]
Let $\mathrm M=(\prod_vd_v)\mu(\mathcal F)\operatorname{Re}e(1)$ be the main term, and let $\kappa_{p,q,j}$ be the flag-class constant. The pointwise orbit input gives $|\varphi(y)|\le M_0=(\prod_vd_v)\operatorname{Re}e(1)$ and $\mathrm M=M_0\mu(\mathcal F)\ge0$.

\begin{proposition}[Near terms on the moderate region]\label{fr:prop:near-moderate}
Use the flag-measure bound, its constant bound, near unipotence, and the pointwise bound with their stated hypotheses. For each $C_g\ge0$ there is an $n_0$, which can be taken to be
\[
 n_0=\max\!\left\{3^{15},\left\lceil e^{16(1+22(|C_g|+4))}\right\rceil\right\},
\]
such that, under the full counting hypotheses with $q\ge1$,
\[
 \int_{\mathcal F\cap\{\ell_*\le s_{\min}\}}
 \sum_{\gamma\in\mathcal N^*(x)}|\varphi(x^{-1}\gamma x)|\,d\mu(x)
 \le e^{-L/4}\mathrm M.
\]
The left side is interpreted as an extended-nonnegative-real integral and the right side through $\mathbb R_{\ge0}\cup\{\infty\}$-valued realification.
\end{proposition}

\begin{proof}
Set $\delta=3\sqrt{n\,10^{-6}/(n^8f^2)}$. Near unipotence and the split-prime argument supply $0<\delta<1/4$ and $\delta qnf\le10^{-2}$. If $s_{\min}(x)\ge1/2$, the same lowering argument as in the preceding proof forces $\mathcal N(x,\delta)=\{1\}$, so the nonidentity near sum is zero.

Fix a composition type $j$ with $J\le q$. On its moderate part, Proposition~\ref{fr:prop:near-count}, finiteness of the near set and the pointwise bound give the pointwise majorant $M_0\mathcal P_\delta(j,\tau)$. The flag-measure bound turns integration over the type into
\[
 M_0\mu(\mathcal F)\kappa_{p,q,j}
 \int_{\tau\in\mathcal S_a,\ \ell_*\le\tau_i}
 \mathcal P_\delta(j,\tau)\operatorname{dens}_{n,f,j}(\tau)\,d\tau.
\]
Here $\operatorname{dens}=\mathcal P_0^{\mathrm{id}}\prod_i(2fj_i/\tau_i)$, where $\mathcal P_0^{\mathrm{id}}$ is the product defining $\mathcal P_\delta$ with $G_\delta(u)$ replaced by $u$. Every type parameter satisfies $0<\tau_i<1/2$ and is increasing. For any argument $\lambda$ occurring in this product,
\[
 (1+12\delta/\lambda)\lambda=\lambda+12\delta\le e^{12\delta}.
\]
The sum of all exponents is at most $3Jnf\le3qnf$, so
\[
 \mathcal P_\delta(j,\tau)\mathcal P_0^{\mathrm{id}}(j,\tau)
 \le e^{36\delta qnf}\le e^{36/100}\le2.
\]
The last inequality uses $36/100<\log2$. Consequently the one-type integral is at most
\[
 \mathrm M\,2\kappa_{p,q,j}\prod_i\int_{\ell_*}^{1/2}\frac{2fj_i}{\tau_i}\,d\tau_i
 =\mathrm M\,2\kappa_{p,q,j}\prod_i\left(2fj_i\log\frac1{2\ell_*}\right).
\]
Since $\ell_*=X^{-11q}$, $\log(1/(2\ell_*))\le11q\log X$, and $j_i\le q$, this is at most
\begin{equation}\label{fr:eq:one-type}
 \mathrm M\,2\kappa_{p,q,j}Y^a.
\end{equation}
There is no lost volume factor: $\mathrm M=M_0\mu(\mathcal F)$. If $\kappa_{p,q,j}\le0$, the realification is zero and the inequality is immediate; otherwise the same real calculation applies.

The constant bound, with $p+q=n$, $f_{\mathrm{deg}}(E)=f$ and $2v_{\mathrm{inv}}(E)=rd$, gives
\[
 \kappa_{p,q,j}\le e^{-JL/2}.
\]
Also $a\le J$, since every block is positive. The next paragraph proves $Y\le e^{L/8}$. Thus
\[
 2\kappa_{p,q,j}Y^a
 \le2\exp(-JL/2+aL/8)\le2e^{-3L/8}.
\]

To prove $Y\le e^{L/8}$, put $K=22(|C_g|+4)$ and $E_0=e^{L/16}$. The chosen threshold gives $\log n\ge16(1+K)$. Hence $nf+K\le L/16$ and $(nf)K\le E_0$. Using $q^2\le n$, $\log rd\le e^n\le E_0$, $\log f\le\sqrt n\le E_0$, $\log b\le E_0$, and $\log g\le C_g\log b\le|C_g|E_0$ gives
\[
 \log X\le(|C_g|+4)E_0,
 \qquad Y=22fq^2\log X\le(nfK)E_0\le E_0^2=e^{L/8}.
\]
Moreover $2^q\le e^q$ and $1+q\le L/8$, so
\begin{equation}\label{fr:eq:final-absorption}
 2\,2^q e^{-3L/8}\le e^{1+q-3L/8}\le e^{-L/4}.
\end{equation}

Finally, points outside the finite type union contribute zero. Canonically filtering a lattice with $s_{\min}<1/2$ (Grayson's canonical filtration \cite[Theorem 1.18 and Definition 1.19, p. 609]{r3grayson}) gives a type with $1\le J\le q$; isotropy gives the upper bound. There are $2^{J-1}$ compositions of $J$. Subadditivity and \eqref{fr:eq:one-type} therefore leave the factor
\[
 \sum_{J=1}^q2^{J-1}=2^q-1<2^q.
\]
Applying \eqref{fr:eq:final-absorption} completes the proof with the original exponent $-L/4$.
\end{proof}

\begin{proposition}[Rare flags]\label{fr:prop:rare-flags}
The nonzero-flag bound implies the rare-flag bound with exponent $-L/4$, with threshold $\lceil e^{15}\rceil$.
\end{proposition}

\begin{proof}
If $n\ge\lceil e^{15}\rceil$, then $\log n\ge15$. Apply the nonzero-flag bound with the unchanged signature, $q^2\le n$, $\log rd\le e^n$, and $\log f\le\sqrt n$ hypotheses. The step characterization gives
\[
 \{s_{\min}<1/2\}=\{S_{<}(1/2)\ne0\}.
\]
Thus its measure divided by $\mu(\mathcal F)$ is at most $e^{-L/3}$. Since $n,f\ge1$ and $\log n\ge15$, one has $L>0$ and
\[
 e^{-L/3}\le e^{-L/4}.
\]
No strengthening of the degree or discriminant hypotheses is required.
\end{proof}

\paragraph{The total error.}
Put $L=nf\log n$, $\delta=\exp(-n^2f)$, and $X=|D-M|$.
The required orbit estimates give
\[
 X\le (2e^{-L/4}+\delta)M+\delta D .
\]
For $n\ge e^{15}$ put $\eta=e^{-L/16}$. Then $240\le L\le n^2f$,
so $e^{-L/4},\delta\le\eta/8$ and $\delta\le1/2$. Since $D\le M+X$,
\[
 X\le (2e^{-L/4}+2\delta)M+\delta X
   \le \tfrac12\eta M+\tfrac12X .
\]
Thus $|D-M|\le e^{-nf\log n/16}M$.
